% Les piles — Chimie 1ère D — Cours signé OnBuch+
% Programme officiel MINESEC, Première C, D, E, TI. Compiler avec : tectonic N07-les-piles.tex
\newcommand{\DOCMATIERE}{Chimie}
\newcommand{\DOCNIVEAU}{1ère D}
\newcommand{\DOCMODULE}{Module 2 : Oxydoréduction}
\newcommand{\DOCLECON}{07}
\newcommand{\DOCTITRE}{Les piles}
% OnBuch+ — préambule « cours signés » ENRICHI (compilé avec tectonic / XeLaTeX)
% Système visuel « Pop » d'OnBuch+ : fond crème (jamais blanc), bordures encre
% épaisses, ombres « dures » décalées, titres Archivo Black en capitales.
% Version enrichie pour les cours de chimie : chimie (mhchem, chemfig),
% graphes (pgfplots), boîtes pédagogiques supplémentaires (définition,
% propriété, méthode, attention, le savais-tu, exercice résolu, exercice,
% TP, bilan), page de garde refaite et signature OnBuch+ ancrée en bas.
%
% Macros attendues dans main.tex AVANT \input{preamble} :
%   \DOCMATIERE  \DOCNIVEAU  \DOCMODULE  \DOCLECON (numéro)  \DOCTITRE
\documentclass[11pt]{article}

\usepackage{fontspec}
\usepackage[french]{babel}
\usepackage[a4paper,margin=2cm,top=3.4cm,bottom=2.8cm,headheight=1.4cm,footskip=1.3cm]{geometry}
\usepackage{amsmath,amssymb,amsfonts}
\usepackage[table]{xcolor} % [table] : fournit aussi \rowcolors (alternance de lignes)
\usepackage{pagecolor}
\usepackage{tikz}
\usetikzlibrary{arrows.meta,positioning,calc,shapes.geometric,shapes.misc,shapes.arrows,decorations.pathmorphing,decorations.markings,patterns,shadows,mindmap,trees,fit,backgrounds,matrix,intersections}
\usepackage{pgfplots}
\pgfplotsset{compat=1.18}
\usepackage[version=4]{mhchem}
\usepackage{chemfig}
\usepackage{enumitem}
\usepackage{fancyhdr}
% [most] charge toutes les bibliothèques tcolorbox (listings, minted, raster,
% documentation...) et gonfle inutilement la mémoire TeX sur un document long
% (~300 boîtes) : on ne charge que ce qui est réellement utilisé.
\usepackage[skins,breakable]{tcolorbox}
\usepackage{graphicx}
\usepackage{colortbl}
\usepackage{booktabs}
\usepackage{tabularx}
\usepackage{array}
\usepackage{multicol}
\usepackage{siunitx}
% parse-numbers=false : accepte aussi \SI{2,5 \times 10^{-3}}{...}
\sisetup{locale=FR,output-decimal-marker={,},group-minimum-digits=4,parse-numbers=false}
\DeclareSIUnit\ohm{\ensuremath{\symup{\Omega}}} % U+2126 absent de la police
\usepackage{qrcode}
% \degree : commande fréquemment utilisée par les modèles pour un angle ou une
% température en dehors de \SI{}{} (non fournie par les paquets chargés).
\providecommand{\degree}{^{\circ}}
\usepackage{lastpage}
\usepackage{hyperref}
\hypersetup{hidelinks}

% ============================================================
% Palette « Pop » OnBuch+ — mêmes hex que la classe P (plus_ui.dart)
% ============================================================
\definecolor{poporange}{HTML}{F59321}
\definecolor{popdark}{HTML}{E07A0C}
\definecolor{popcream}{HTML}{FFF6E8}
\definecolor{popink}{HTML}{1C1714}
\definecolor{poppurple}{HTML}{7A5AE0}
\definecolor{popgreen}{HTML}{2E9E6A}
\definecolor{poppink}{HTML}{E0457B}
\definecolor{popblue}{HTML}{2F7DE1}
\definecolor{popgold}{HTML}{C98A12}
\definecolor{popmuted}{HTML}{8A7F76}
\definecolor{popline}{HTML}{EADFD2}
\definecolor{popgray}{HTML}{8A7F76}
\colorlet{popgrey}{popgray}
% Alias tolérants (le modèle nomme parfois les couleurs différemment)
\colorlet{popwhite}{white}
\colorlet{popblack}{popink}
\colorlet{popred}{poppink}
\colorlet{popyellow}{popgold}
% Teintes claires pour fonds de tableaux / zones
\colorlet{popblueL}{popblue!10!white}
\colorlet{poporangeL}{poporange!14!white}
\colorlet{popgreenL}{popgreen!12!white}
\colorlet{poppurpleL}{poppurple!12!white}
\colorlet{poppinkL}{poppink!12!white}
\colorlet{popgoldL}{popgold!14!white}

\pagecolor{popcream}

% ============================================================
% Typographie
% ============================================================
\defaultfontfeatures{Ligatures=TeX}
\setmainfont{PlusJakartaSans-Medium.ttf}[Path=../../../fonts/,BoldFont=PlusJakartaSans-Bold.ttf]
% Maths en sans-serif (Fira Math) avec chiffres et lettres droites de Plus
% Jakarta, pour que les formules se fondent dans le texte.
\usepackage[math-style=ISO,bold-style=ISO]{unicode-math}
\setmathfont{FiraMath-Regular.otf}[Path=../../../fonts/]
\setmathfont{PlusJakartaSans-Medium.ttf}[Path=../../../fonts/,range={up/num,up/{Latin,latin},\mathup}]
\usepackage{icomma} % virgule décimale sans espace en mode math
% Caractères absents de la police de texte (sinon carré vide) : rendus par
% leur équivalent mathématique, en texte comme en formule.
\usepackage{newunicodechar}
\newunicodechar{α}{\ensuremath{\alpha}}
\newunicodechar{Α}{\ensuremath{\alpha}}
\newunicodechar{β}{\ensuremath{\beta}}
\newunicodechar{δ}{\ensuremath{\delta}}
\newunicodechar{✓}{\popcheck{popgreen}}
\newfontfamily\popdisplay{ArchivoBlack-Regular.ttf}[Path=../../../fonts/]
\newfontfamily\poplogo{SpaceGrotesk-Bold.ttf}[Path=../../../fonts/]
\newfontfamily\popsemi{PlusJakartaSans-SemiBold.ttf}[Path=../../../fonts/]
\newfontfamily\popxbold{PlusJakartaSans-ExtraBold.ttf}[Path=../../../fonts/]
\newfontfamily\popxboldls{PlusJakartaSans-ExtraBold.ttf}[Path=../../../fonts/,LetterSpace=3.0]

\setlist[enumerate]{leftmargin=1.7em,itemsep=5pt,topsep=4pt}
\setlist[itemize]{leftmargin=1.7em,itemsep=4pt,topsep=4pt,label={\color{poporange}\textbullet}}
\setlength{\parindent}{0pt}
\setlength{\parskip}{5pt}
\renewcommand{\arraystretch}{1.25}

% chemfig : liaisons un peu plus épaisses, encre Pop
\setchemfig{atom sep=2em,bond style={line width=0.9pt,popink},cram width=3pt}

% pgfplots : style Pop par défaut
\pgfplotsset{
  every axis/.append style={axis line style={popink,line width=1pt},tick style={popink},
    grid style={popline,line width=0.5pt},label style={font=\small},tick label style={font=\footnotesize}},
  popcurve/.style={poporange,line width=1.8pt,no markers,smooth},
}

% ============================================================
% Pill / squiggle
% ============================================================
\newcommand{\poppill}[3]{%
  \tikz[baseline=-0.55ex]\node[draw=popink,line width=1.2pt,rounded corners=2.6pt,%
    fill=#2,inner xsep=6.5pt,inner ysep=3pt]{%
    \popxboldls\fontsize{7}{7}\selectfont\color{#3}\MakeUppercase{#1}};%
}
\newcommand{\popsquiggle}[1]{%
  \tikz[baseline=-0.4ex]{
    \draw[#1,line width=2.6pt,line cap=round]
      plot [smooth] coordinates {(0,0.09) (0.34,0.01) (0.68,0.21) (1.02,0.01) (1.36,0.10)};
  }%
}
% Mot-clé surligné (marqueur orange sous le texte)
\newcommand{\cle}[1]{{\popxbold\color{popdark}#1}}

% ============================================================
% En-tête / pied
% ============================================================
\pagestyle{fancy}
\fancyhf{}
\renewcommand{\headrulewidth}{0pt}
\renewcommand{\footrulewidth}{0pt}
\lhead{\raisebox{-0.15ex}{\poplogo\fontsize{13}{13}\selectfont\color{popink}OnBuch\color{poporange}+}}
\rhead{\poppill{\DOCMATIERE\ \textbullet\ \DOCNIVEAU\ \textbullet\ Leçon \DOCLECON}{white}{popink}}
\lfoot{\popxbold\fontsize{7.6}{7.6}\selectfont\color{popmuted}\MakeUppercase{Cours signé OnBuch+}\ \textbullet\ {\popsemi\DOCTITRE}}
\rfoot{\tikz[baseline=-0.6ex]\node[rounded corners=8.5pt,fill=popink,minimum height=17pt,minimum width=17pt,inner xsep=5pt,inner sep=0]{\popxbold\fontsize{8}{8}\selectfont\color{popcream}\thepage\,/\,\pageref*{LastPage}};}

% ============================================================
% Titres
% ============================================================
\newcommand{\chaptitre}[2]{%
  \begin{tcolorbox}[enhanced,colback=poporange,colframe=popink,boxrule=2.6pt,arc=11pt,
      drop shadow=popink,left=15pt,right=15pt,top=12pt,bottom=11pt,before skip=3pt,after skip=18pt]
    \poppill{#2}{popink}{popcream}\par\vspace{9pt}
    {\raggedright\hyphenpenalty=10000\exhyphenpenalty=10000\popdisplay\fontsize{22}{24}\selectfont\color{popcream}\MakeUppercase{#1}\par}
  \end{tcolorbox}
}
% Section numérotée : pastille orange avec le numéro + titre Archivo
\newcounter{popsec}
\newcommand{\coursec}[1]{%
  \stepcounter{popsec}\setcounter{popsub}{0}%
  \par\vspace{16pt}\noindent
  \begin{minipage}{\linewidth}
  \tikz[baseline=-0.7ex]\node[draw=popink,line width=1.6pt,fill=poporange,rounded corners=4pt,minimum size=22pt,inner sep=0]{\popdisplay\fontsize{12}{12}\selectfont\color{popcream}\thepopsec};%
  \hspace{8pt}{\popdisplay\fontsize{15}{17}\selectfont\color{popink}\MakeUppercase{#1}}\par
  \vspace{4pt}\noindent{\color{poporange}\rule{36pt}{3.6pt}}
  \end{minipage}\par\nopagebreak\vspace{8pt}
}
\newcounter{popsub}
\newcommand{\courssub}[1]{%
  \stepcounter{popsub}%
  \par\vspace{9pt}\noindent{\popxbold\fontsize{11.5}{13}\selectfont\color{popink}{\color{poporange}\thepopsec.\thepopsub}\hspace{6pt}#1}\par\nopagebreak\vspace{4pt}
}

% ============================================================
% Boîtes pédagogiques — même recette : fond blanc, bordure épaisse colorée,
% ombre dure encre, badge plein. Titre optionnel [#1].
% ============================================================
\tcbset{popbase/.style={breakable,enhanced,colback=white,boxrule=2.3pt,arc=9pt,
  shadow={3pt}{-3pt}{0pt}{popink},left=14pt,right=14pt,top=11pt,bottom=11pt,before skip=11pt,after skip=15pt,
  fontupper=\popsemi\fontsize{10.6}{14.5}\selectfont\color{popink},
  fontlower=\popsemi\fontsize{10.6}{14.5}\selectfont\color{popink}}}
% Coche dessinée (le glyphe ✓ n'existe pas dans les polices du document)
\newcommand{\popcheck}[1]{\tikz[baseline=-0.5ex]\draw[#1,line width=1.6pt,line cap=round,line join=round](-0.13,0.0)--(-0.04,-0.09)--(0.13,0.1);}
\newcommand{\popboxhead}[3]{\poppill{#1}{#2}{white}\ifx\relax#3\relax\else\hspace{6pt}{\popxbold\fontsize{10.6}{12}\selectfont\color{popink}#3}\fi\par\vspace{7pt}}

\newtcolorbox{aretenir}[1][]{popbase,colframe=popgreen,before upper={\popboxhead{À retenir}{popgreen}{#1}}}
\newtcolorbox{exemplebox}[1][]{popbase,colframe=poporange,before upper={\popboxhead{Exemple}{poporange}{#1}}}
\newtcolorbox{definition}[1][]{popbase,colframe=popblue,before upper={\popboxhead{Définition}{popblue}{#1}}}
\newtcolorbox{propriete}[1][]{popbase,colframe=poppurple,before upper={\popboxhead{Propriété}{poppurple}{#1}}}
\newtcolorbox{methode}[1][]{popbase,colframe=popgold,before upper={\popboxhead{Méthode}{popgold}{#1}}}
\newtcolorbox{attention}[1][]{popbase,colframe=poppink,before upper={\popboxhead{Attention}{poppink}{#1}}}
\newtcolorbox{experience}[1][]{popbase,colframe=popblue,colback=popblueL,before upper={\popboxhead{Expérience}{popblue}{#1}}}
% « Le savais-tu ? » : carte violette pleine (contexte, culture, Cameroun)
\newtcolorbox{savaistu}[1][]{popbase,colback=poppurple,colframe=popink,
  fontupper=\popsemi\fontsize{10.4}{14.2}\selectfont\color{white},
  before upper={\poppill{Le savais-tu\,?}{white}{poppurple}\ifx\relax#1\relax\else\hspace{6pt}{\popxbold\color{white}#1}\fi\par\vspace{7pt}}}
% Exercice résolu : énoncé en haut, \tcblower puis la solution sur fond crème
\newcounter{popexo}\newcounter{popexr}
\newtcolorbox{exoresolu}[1][]{popbase,colframe=poporange,colbacklower=popcream,
  segmentation style={popink,line width=1.2pt,dashed},
  before upper={\stepcounter{popexr}\popboxhead{Exercice résolu \thepopexr}{poporange}{#1}},
  before lower={\poppill{Solution}{popink}{popcream}\par\vspace{6pt}}}
% Exercice (non résolu en place) — la correction est dans la partie Corrigés
\newtcolorbox{exercice}[1][]{popbase,colframe=popink,
  before upper={\stepcounter{popexo}\popboxhead{Exercice \thepopexo}{popink}{#1}}}
% Corrigé
\newtcolorbox{corrige}[1][]{popbase,colframe=popgreen,colback=popgreenL,
  before upper={\popboxhead{Corrigé}{popgreen}{#1}}}
% Objectifs / prérequis / bilan
\newtcolorbox{objectifs}{popbase,colframe=popink,colback=poporangeL,
  before upper={\popboxhead{Objectifs de la leçon}{popink}{}}}
\newtcolorbox{prerequis}{popbase,colframe=popink,colback=popblueL,
  before upper={\popboxhead{Prérequis}{popblue}{}}}
\newtcolorbox{bilan}{popbase,colframe=popink,colback=white,boxrule=2.8pt,
  before upper={\popboxhead{Fiche bilan}{poporange}{L'essentiel à connaître pour l'examen}}}
% Figure encadrée avec légende
\newtcolorbox{popfigure}[1][]{enhanced,breakable=false,colback=white,colframe=popline,boxrule=1.4pt,arc=8pt,
  left=10pt,right=10pt,top=10pt,bottom=8pt,before skip=10pt,after skip=12pt,halign=center,
  lower separated=false,halign lower=center,
  fontlower=\popsemi\fontsize{9}{11.5}\selectfont\color{popmuted},
  #1}
\newcommand{\legende}[1]{\par\vspace{6pt}{\popsemi\fontsize{9}{11.5}\selectfont\color{popmuted}#1}}

% ============================================================
% Signature OnBuch+
% ============================================================
\newcommand{\poplogoinline}[1]{{\poplogo\fontsize{#1}{#1}\selectfont\color{popink}OnBuch\color{poporange}+}}
\newcommand{\signatureonbuch}{%
  \begin{tcolorbox}[enhanced,colback=popink,colframe=popink,boxrule=2.2pt,arc=10pt,
      drop shadow=poporange,left=14pt,right=14pt,top=10pt,bottom=10pt,before skip=0pt,after skip=0pt]
    \tikz[baseline=-0.7ex]\node[circle,fill=popgreen,draw=popcream,line width=1.3pt,minimum size=22pt,inner sep=0]{\popcheck{white}};%
    \hspace{9pt}\begin{minipage}[c]{0.62\linewidth}
      {\popxbold\fontsize{12}{13}\selectfont\color{popcream}Signé\ \textbullet\ {\poplogo OnBuch\color{poporange}+}}\par\vspace{2pt}
      {\raggedright\popsemi\fontsize{8}{10.5}\selectfont\color{popline}\MakeUppercase{Équipe pédagogique OnBuch+}\\\MakeUppercase{Conforme au programme officiel MINESEC}\par}
    \end{minipage}\hfill
    {\poplogo\fontsize{22}{22}\selectfont\color{popcream}OnBuch\color{poporange}+}
  \end{tcolorbox}%
}

% ============================================================
% Page de garde
% #1 titre · #2 matière · #3 niveau · #4 module · #5 n° leçon · #6 durée ·
% #7 description / objectifs (texte ou itemize)
% ============================================================
\newcommand{\pagegarde}[7]{%
  \thispagestyle{empty}
  \newgeometry{a4paper,margin=1.7cm,top=1.6cm,bottom=1.4cm}%
  \begin{tikzpicture}[remember picture,overlay]
    \fill[poporange!18] ([xshift=-3.2cm,yshift=-3.6cm]current page.north east) circle (4.6cm);
    \fill[poppurple!14] ([xshift=2.4cm,yshift=7.6cm]current page.south west) circle (3.8cm);
  \end{tikzpicture}%
  {\poplogo\fontsize{30}{30}\selectfont\color{popink}OnBuch\color{poporange}+}\hfill
  \raisebox{0.6ex}{\poppill{Cours signé \textbullet\ Premium}{popink}{popcream}}\par
  \vspace{1pt}\popsquiggle{poppurple}\par
  \vspace{8pt}
  \begin{tcolorbox}[enhanced,colback=poporange,colframe=popink,boxrule=2.8pt,arc=13pt,
      drop shadow=popink,left=18pt,right=18pt,top=12pt,bottom=13pt,after skip=10pt]
    \poppill{#2 \textbullet\ #3}{popink}{popcream}\hspace{5pt}\poppill{#4}{white}{popink}\par\vspace{8pt}
    {\popdisplay\fontsize{14}{14}\selectfont\color{popink}LEÇON #5}\par\vspace{4pt}
    {\raggedright\hyphenpenalty=10000\exhyphenpenalty=10000\popdisplay\fontsize{25}{27}\selectfont\color{popcream}\MakeUppercase{#1}\par}
    \vspace{8pt}
    {\popxbold\fontsize{9.5}{9.5}\selectfont\color{popink}\MakeUppercase{Durée conseillée : #6}}
  \end{tcolorbox}
  \begin{tcolorbox}[enhanced,colback=white,colframe=popink,boxrule=2.2pt,arc=10pt,
      drop shadow=popink,left=16pt,right=16pt,top=10pt,bottom=10pt,after skip=8pt]
    \poppill{Dans cette leçon}{poppurple}{white}\par\vspace{5pt}
    \popsemi\fontsize{9.3}{12.6}\selectfont\color{popink}#7
  \end{tcolorbox}
  \noindent\begin{minipage}[t]{0.487\linewidth}
    \begin{tcolorbox}[enhanced,colback=popgreen,colframe=popink,boxrule=2.2pt,arc=10pt,
        drop shadow=popink,left=11pt,right=11pt,top=10pt,bottom=10pt,height=3.1cm,valign=center]
      \tcbox[colback=white,colframe=popink,boxrule=1.3pt,arc=5pt,boxsep=0pt,nobeforeafter,
        left=3pt,right=3pt,top=3pt,bottom=3pt,valign=center]{\qrcode[height=1.9cm]{https://play.google.com/store/apps/details?id=com.luvvix.onbuch&hl=fr}}%
      \hspace{8pt}\begin{minipage}[c]{0.5\linewidth}
        {\popdisplay\fontsize{11}{12.5}\selectfont\color{white}\MakeUppercase{Télécharge OnBuch}}\par\vspace{4pt}
        {\raggedright\popsemi\fontsize{8.4}{11}\selectfont\color{white}Gratuit \textbullet\ Google Play\\Tuteur IA, annales, concours, résultats\par}
      \end{minipage}
    \end{tcolorbox}
  \end{minipage}\hfill
  \begin{minipage}[t]{0.487\linewidth}
    \begin{tcolorbox}[enhanced,colback=poppurple,colframe=popink,boxrule=2.2pt,arc=10pt,
        drop shadow=popink,left=11pt,right=11pt,top=10pt,bottom=10pt,height=3.1cm,valign=center]
      \tikz[baseline=-0.7ex]\node[circle,fill=white,draw=popink,line width=1.3pt,minimum size=1.6cm,inner sep=0]{\popdisplay\fontsize{24}{24}\selectfont\color{poppurple}+};%
      \hspace{8pt}\begin{minipage}[c]{0.6\linewidth}
        {\popdisplay\fontsize{11}{12.5}\selectfont\color{white}\MakeUppercase{Passe à OnBuch+}}\par\vspace{4pt}
        {\raggedright\popsemi\fontsize{8.4}{11}\selectfont\color{white}Tous les cours signés, Tuteur IA illimité, fiches PDF — onglet OnBuch+ de l'app\par}
      \end{minipage}
    \end{tcolorbox}
  \end{minipage}
  \vspace{8pt}
  \popcontenu
  \vfill
  \signatureonbuch
  \restoregeometry
  \newpage
}

% Rangée « Ce cours contient » (page de garde)
\newcommand{\poptile}[3]{% #1 couleur #2 gros texte #3 légende
  \begin{tcolorbox}[enhanced,colback=#1,colframe=popink,boxrule=2pt,arc=9pt,shadow={2.5pt}{-2.5pt}{0pt}{popink},
      left=6pt,right=6pt,top=9pt,bottom=9pt,halign=center,valign=center,height=1.75cm,width=\linewidth,nobeforeafter]
    {\popdisplay\fontsize{12.5}{13}\selectfont\color{white}\MakeUppercase{#2}}\par\vspace{3pt}
    {\centering\popsemi\fontsize{7.8}{9.5}\selectfont\color{white}#3\par}
  \end{tcolorbox}}
\newcommand{\popcontenu}{%
  \par\noindent{\popxboldls\fontsize{7.5}{7.5}\selectfont\color{popmuted}CE COURS SIGNÉ CONTIENT}\par\vspace{7pt}
  \noindent\begin{minipage}[t]{0.235\linewidth}\poptile{popblue}{Cours}{complet, illustré, pas à pas}\end{minipage}\hfill
  \begin{minipage}[t]{0.235\linewidth}\poptile{poporange}{Méthodes}{et exercices résolus}\end{minipage}\hfill
  \begin{minipage}[t]{0.235\linewidth}\poptile{poppink}{Exercices}{type examen + corrigés}\end{minipage}\hfill
  \begin{minipage}[t]{0.235\linewidth}\poptile{popgreen}{Fiche bilan}{l'essentiel à retenir}\end{minipage}\par}

% Sommaire
\newtcolorbox{sommaire}{enhanced,colback=white,colframe=popink,boxrule=2.2pt,arc=10pt,
  drop shadow=popink,left=18pt,right=18pt,top=13pt,bottom=13pt,before skip=6pt,after skip=20pt,
  before upper={\poppill{Sommaire}{poppurple}{white}\par\vspace{9pt}\popsemi\fontsize{10.6}{14.5}\selectfont\color{popink}}}

% Signature de fin de cours — poussée tout en bas de la dernière page
\newcommand{\findecours}{%
  \par\vfill\nopagebreak
  \begin{center}\popsquiggle{poporange}\end{center}\vspace{4pt}
  \signatureonbuch
}
% Glyphes absents de Fira Math : redéfinis à partir de glyphes disponibles
\AtBeginDocument{%
  \def\setminus{\mathbin{\backslash}}\let\smallsetminus\setminus
  \def\vdots{\mathord{\vbox{\baselineskip3.2pt\lineskiplimit0pt\kern4pt\hbox{.}\hbox{.}\hbox{.}}}}%
  \def\ddots{\mathinner{\mkern1mu\raise6.5pt\hbox{.}\mkern2mu\raise3.6pt\hbox{.}\mkern2mu\raise0.7pt\hbox{.}\mkern1mu}}%
  \def\triangle{\mathord{\scalebox{0.9}{$\symup{\Delta}$}}}%
  \def\nabla{\mathord{\raisebox{\depth}{\scalebox{1}[-1]{$\symup{\Delta}$}}}}%
  \def\checkmark{\popcheck{popdark}}%
  \def\star{\mathbin{\ast}}\def\bigstar{\mathord{\ast}}%
  \def\diamond{\mathbin{\scalebox{0.6}{$\Diamond$}}}%
  \def\bigcup{\mathop{\mathchoice{\raisebox{-0.3em}{\scalebox{1.7}{$\cup$}}}{\scalebox{1.25}{$\cup$}}{\cup}{\cup}}\displaylimits}%
  \def\bigcap{\mathop{\mathchoice{\raisebox{-0.3em}{\scalebox{1.7}{$\cap$}}}{\scalebox{1.25}{$\cap$}}{\cap}{\cap}}\displaylimits}%
  \def\lesseqqgtr{\mathrel{\lessgtr}}\def\bigoplus{\mathop{\oplus}\displaylimits}%
}
\providecommand{\ssi}{si et seulement si} % abréviation fréquente

%% FIGFIX-BEGIN v5 — correctifs globaux des figures (docs/onbuch-generation/tools/figfix)
\usepackage{adjustbox}\usepackage{etoolbox}\tcbuselibrary{raster}
% toute figure TikZ/pgfplots plus large que la ligne est réduite à la largeur disponible
\BeforeBeginEnvironment{tikzpicture}{\begin{adjustbox}{max width=\linewidth}}
\AfterEndEnvironment{tikzpicture}{\end{adjustbox}}
% légendes pgfplots lisibles par-dessus les courbes
\pgfplotsset{every axis/.append style={legend style={fill=white,fill opacity=0.92,text opacity=1,draw=black!15,inner sep=3pt}}}
% étiquettes d'arêtes (syntaxe edge["..."]) avec fond blanc
\tikzset{every edge quotes/.append style={fill=white,inner sep=1.2pt}}
%% FIGFIX-END
\begin{document}

\pagegarde{Les piles}{Chimie}{1ère D}{Module 2 : Oxydoréduction}{07}{4 h}{Cette leçon étudie les piles électrochimiques à travers l'exemple fondateur de la pile Daniell : constitution, polarité, notation et fonctionnement. Elle introduit le potentiel d'oxydoréduction et la classification quantitative des couples, outils permettant de prévoir les réactions, les pôles et la f.é.m. de toute pile.
\vspace{4pt}
\begin{multicols}{2}\small
\begin{itemize}
\item À la fin de cette leçon, je sais décrire la constitution d'une demi-pile et de la pile Daniell (électrodes, solutions, pont salin).
\item À la fin de cette leçon, je sais identifier les pôles d'une pile et justifier le sens de circulation des électrons et du courant.
\item À la fin de cette leçon, je sais écrire les équations aux électrodes et l'équation-bilan de fonctionnement d'une pile.
\item À la fin de cette leçon, je sais représenter une pile par sa notation conventionnelle et la schématiser.
\item À la fin de cette leçon, je sais définir le potentiel d'oxydoréduction et le potentiel standard d'un couple, en référence à l'électrode à hydrogène.
\item À la fin de cette leçon, je sais utiliser la classification quantitative des couples pour prévoir la réaction spontanée entre deux couples.
\end{itemize}\end{multicols}}

\begin{sommaire}
\begin{multicols}{2}
\begin{enumerate}
\item Étude expérimentale de la pile Daniell
\item Fonctionnement, équations et notation conventionnelle
\item La demi-pile de référence : l'électrode à hydrogène
\item Potentiel d'oxydoréduction et classification quantitative
\item Prévision de la polarité et calcul de la f.é.m.
\item Généralisation, piles usuelles et synthèse
\item Travaux pratiques
\item Méthodes et exercices résolus
\item Exercices
\item Corrigés des exercices
\item Fiche bilan
\end{enumerate}
\end{multicols}
\end{sommaire}

\begin{objectifs}
\begin{itemize}
\item À la fin de cette leçon, je sais décrire la constitution d'une demi-pile et de la pile Daniell (électrodes, solutions, pont salin).
\item À la fin de cette leçon, je sais identifier les pôles d'une pile et justifier le sens de circulation des électrons et du courant.
\item À la fin de cette leçon, je sais écrire les équations aux électrodes et l'équation-bilan de fonctionnement d'une pile.
\item À la fin de cette leçon, je sais représenter une pile par sa notation conventionnelle et la schématiser.
\item À la fin de cette leçon, je sais définir le potentiel d'oxydoréduction et le potentiel standard d'un couple, en référence à l'électrode à hydrogène.
\item À la fin de cette leçon, je sais utiliser la classification quantitative des couples pour prévoir la réaction spontanée entre deux couples.
\item À la fin de cette leçon, je sais déterminer la polarité et calculer la f.é.m. d'une pile à partir des potentiels standard.
\item À la fin de cette leçon, je sais réaliser expérimentalement une pile, mesurer sa f.é.m. et déterminer ses pôles avec un voltmètre.
\end{itemize}
\end{objectifs}

\begin{prerequis}
\begin{itemize}
\item Couple oxydant/réducteur, demi-équations électroniques et équations d'oxydoréduction (leçons 05 et 06 du module 2).
\item Règle du gamma (classification qualitative) pour prévoir le sens d'une réaction.
\item Notions d'électricité de Seconde : courant, tension, utilisation du voltmètre et de l'ampèremètre.
\item Ions en solution aqueuse et électroneutralité (dissolution des sels, ex. CuSO4, ZnSO4).
\item Sens conventionnel du courant et sens réel des électrons dans un circuit.
\end{itemize}
\end{prerequis}

\newpage

\coursec{Étude expérimentale de la pile Daniell}

\courssub{Situation d'accroche : l'électricité cachée dans la chimie}

Au quartier, quand le courant ENEO se coupe, on sort la lampe torche ou la radio à piles : mais comment une simple pile transforme-t-elle des produits chimiques en électricité ? Pourquoi une pile usée ne « donne »-t-elle plus de courant ? Et comment les techniciens des centres de santé choisissent-ils la bonne pile pour un appareil ? En réalité, chaque pile cache une réaction d'\cle{oxydoréduction} soigneusement organisée. Cette leçon nous apprend à construire une pile, à prévoir ses pôles et à calculer sa tension.

\textbf{Problématique.} Une réaction d'oxydoréduction est un transfert d'électrons d'un réducteur vers un oxydant. Peut-on \emph{canaliser} ce transfert d'électrons dans un circuit extérieur pour produire un courant électrique utilisable ?

\courssub{Rappel : la réaction spontanée entre le zinc et les ions cuivre(II)}

Reprenons l'expérience classique vue au chapitre précédent : on plonge un clou de zinc dans une solution bleue de sulfate de cuivre \ce{CuSO4} (c'est-à-dire contenant des ions \ce{Cu^2+} et \ce{SO4^2-}). Au bout de quelques minutes, on observe :

\begin{itemize}
\item un dépôt \textbf{rougeâtre} de cuivre métallique sur le clou de zinc ;
\item la \textbf{décoloration progressive} de la solution (les ions \ce{Cu^2+}, responsables de la couleur bleue, disparaissent) ;
\item un léger dégagement de chaleur : la réaction est exothermique.
\end{itemize}

L'équation-bilan de cette réaction \textbf{spontanée} s'écrit :

\[ \ce{Zn(s) + Cu^2+(aq) -> Zn^2+(aq) + Cu(s)} \]

Analysons-la en termes d'oxydoréduction, avec les deux couples mis en jeu \ce{Zn^2+}/\ce{Zn} et \ce{Cu^2+}/\ce{Cu} :

\begin{align*}
\text{Oxydation du zinc : } & \ce{Zn -> Zn^2+ + 2e-} \quad \text{(le zinc, réducteur, cède 2 électrons)}\\
\text{Réduction des ions cuivre : } & \ce{Cu^2+ + 2e- -> Cu} \quad \text{(l'ion cuivre, oxydant, capte 2 électrons)}
\end{align*}

\begin{attention}[L'essentiel à comprendre]
Dans ce mélange direct, le transfert d'électrons se fait \textbf{au contact} entre le zinc et les ions \ce{Cu^2+}, directement à la surface du métal. Les électrons ne parcourent \emph{aucun circuit extérieur} : l'énergie chimique est dégradée en \textbf{chaleur}, et non récupérée sous forme d'électricité. C'est exactement comme une rivière qui dévale une pente sans passer par une turbine : l'énergie est perdue.
\end{attention}

\courssub{L'idée fondamentale : séparer les deux couples}

Voici l'intuition géniale qui a donné naissance à toutes les piles : si les électrons passent spontanément du zinc vers les ions \ce{Cu^2+}, alors \textbf{obligeons-les à faire un détour} !

Séparons physiquement le réducteur (le zinc) de l'oxydant (les ions \ce{Cu^2+}) dans deux récipients distincts, et relions-les uniquement par un fil électrique. Le transfert d'électrons reste spontané, mais il ne peut plus se faire qu'\textbf{à travers le circuit extérieur}. Ce flux d'électrons canalisé, c'est précisément un \textbf{courant électrique} : on peut y brancher une lampe, un moteur, une radio.

\begin{definition}[Demi-pile et pile]
Une \cle{demi-pile} est l'association d'un conducteur métallique (l'\textbf{électrode}) et d'une solution contenant l'oxydant ou le réducteur du \textbf{même couple} d'oxydoréduction. Exemple : une lame de zinc plongée dans une solution de sulfate de zinc réalise la demi-pile du couple \ce{Zn^2+}/\ce{Zn}.

Une \cle{pile} électrochimique est l'association de \textbf{deux demi-piles} mettant en jeu deux couples différents, reliées par un circuit électrique extérieur et par un \textbf{pont salin} (ou une paroi poreuse). Une pile transforme l'énergie chimique d'une réaction d'oxydoréduction spontanée en énergie électrique.
\end{definition}

\courssub{Constitution de la pile Daniell}

La pile Daniell associe les deux couples \ce{Zn^2+}/\ce{Zn} et \ce{Cu^2+}/\ce{Cu}. Elle comprend quatre éléments indispensables :

\begin{enumerate}
\item \textbf{Demi-pile zinc} : une lame de zinc plongée dans une solution de sulfate de zinc \ce{ZnSO4} de concentration $c = \SI{1,0}{mol.L^{-1}}$ (ions \ce{Zn^2+}).
\item \textbf{Demi-pile cuivre} : une lame de cuivre plongée dans une solution de sulfate de cuivre \ce{CuSO4} de concentration $c = \SI{1,0}{mol.L^{-1}}$ (ions \ce{Cu^2+}, solution bleue).
\item \textbf{Le pont salin} : un tube en U rempli d'un gel conducteur contenant une solution ionique concentrée, par exemple de chlorure de potassium \ce{KCl} ou de nitrate de potassium \ce{KNO3}. Ses extrémités plongent dans les deux solutions.
\item \textbf{Le circuit extérieur} : des fils reliant les deux lames métalliques, avec un appareil de mesure (voltmètre, galvanomètre) ou un récepteur (DEL, moteur).
\end{enumerate}

\begin{propriete}[Rôle du pont salin]
Le pont salin joue \textbf{deux rôles} indispensables :
\begin{enumerate}
\item il \textbf{ferme le circuit} électrique en permettant le passage du courant entre les deux solutions (sans lui, le circuit serait ouvert et aucun courant ne circulerait) ;
\item il assure l'\textbf{électroneutralité} de chaque solution par \textbf{migration des ions} : les anions (ici \ce{Cl-}) migrent vers la demi-pile du zinc (qui s'enrichit en cations \ce{Zn^2+}), tandis que les cations (ici \ce{K+}) migrent vers la demi-pile du cuivre (qui s'appauvrit en cations \ce{Cu^2+}).
\end{enumerate}
Les ions du pont salin ne participent \emph{pas} à la réaction d'oxydoréduction : ils sont dits « spectateurs » ou « indifférents ». On les choisit précisément pour cela (ni oxydants, ni réducteurs dans le domaine de potentiel concerné).
\end{propriete}

\begin{attention}[Pourquoi pas un simple fil entre les deux solutions ?]
Un fil métallique conduit les \textbf{électrons}, mais pas les \textbf{ions}. Or, si l'on reliait les deux solutions par un fil, les électrons ne pourraient pas « traverser » les solutions. Seul un conducteur ionique (pont salin ou paroi poreuse) peut fermer le circuit à l'intérieur des solutions. Sans pont salin, la pile ne débite \textbf{aucun} courant : c'est une cause classique d'échec de l'expérience au laboratoire !
\end{attention}

\begin{popfigure}
\begin{center}
\begin{tikzpicture}[scale=0.95, >=Stealth]
% Bécher gauche (zinc)
\draw[thick, popdark] (-4.6,0) -- (-4.6,2.6) (-1.6,0) -- (-1.6,2.6);
\draw[thick, popdark] (-4.6,0) -- (-1.6,0);
\fill[popblue!25] (-4.5,0.05) rectangle (-1.7,1.9);
\draw[popblue, thick] (-4.5,1.9) -- (-1.7,1.9);
\node[popblue, font=\small] at (-3.1,0.5) {$\ce{ZnSO4}$};
\node[popblue, font=\footnotesize] at (-3.1,1.0) {$\ce{Zn^2+} + \ce{SO4^2-}$};
% Lame de zinc
\fill[popmuted!60] (-3.35,1.2) rectangle (-2.85,3.6);
\node[font=\small\bfseries, popdark] at (-3.1,3.85) {Zn};
\node[font=\footnotesize, popdark, align=center] at (-5.3,3.1) {pôle $-$\\ (anode)};
% Bécher droit (cuivre)
\draw[thick, popdark] (1.6,0) -- (1.6,2.6) (4.6,0) -- (4.6,2.6);
\draw[thick, popdark] (1.6,0) -- (4.6,0);
\fill[popblue!50] (1.7,0.05) rectangle (4.5,1.9);
\draw[popblue!70!black, thick] (1.7,1.9) -- (4.5,1.9);
\node[popblue!70!black, font=\small] at (3.1,0.5) {$\ce{CuSO4}$};
\node[popblue!70!black, font=\footnotesize] at (3.1,1.0) {$\ce{Cu^2+} + \ce{SO4^2-}$};
% Lame de cuivre
\fill[poporange!80] (2.85,1.2) rectangle (3.35,3.6);
\node[font=\small\bfseries, poporange!80!black] at (3.1,3.85) {Cu};
\node[font=\footnotesize, popdark, align=center] at (5.3,3.1) {pôle $+$\\ (cathode)};
% Pont salin
\draw[line width=5pt, popgreen!60] (-2.2,1.6) -- (-2.2,3.0) .. controls (-2.2,3.7) and (2.2,3.7) .. (2.2,3.0) -- (2.2,1.6);
\node[font=\footnotesize, popgreen!50!black, align=center] at (0,4.15) {pont salin (gel de \ce{KCl})};
% Migrations ioniques dans le pont
\draw[->, popgreen!50!black, thick] (-0.5,3.35) -- (-1.5,3.35) node[midway, above, font=\footnotesize] {$\ce{Cl-}$};
\draw[->, popgreen!50!black, thick] (0.5,3.35) -- (1.5,3.35) node[midway, above, font=\footnotesize] {$\ce{K+}$};
% Circuit extérieur + voltmètre
\draw[thick, popdark] (-3.1,3.6) -- (-3.1,4.9) -- (-0.6,4.9);
\draw[thick, popdark] (3.1,3.6) -- (3.1,4.9) -- (0.6,4.9);
\draw[thick, popdark, fill=white] (0,4.9) circle (0.6);
\node[font=\large\bfseries, popdark] at (0,4.9) {V};
\node[font=\footnotesize, popdark] at (0,5.75) {voltmètre};
% Sens des électrons
\draw[->, poppink, very thick] (-2.6,5.25) -- (-1.0,5.25) node[midway, above, font=\footnotesize\bfseries] {$e^-$};
% Légendes réactions
\node[font=\footnotesize, align=center, popdark] at (-3.1,-0.7) {oxydation : $\ce{Zn -> Zn^2+ + 2e-}$\\ la lame s'amincit};
\node[font=\footnotesize, align=center, popdark] at (3.1,-0.7) {réduction : $\ce{Cu^2+ + 2e- -> Cu}$\\ dépôt rougeâtre};
\end{tikzpicture}
\end{center}
\legende{Schéma annoté de la pile Daniell en fonctionnement. Les électrons circulent dans le circuit extérieur de la lame de zinc vers la lame de cuivre ; les ions migrent dans le pont salin pour garantir l'électroneutralité des solutions.}
\end{popfigure}

\courssub{Observations expérimentales : la pile débite un courant}

\begin{experience}[Réalisation de la pile Daniell au laboratoire]
\textbf{Protocole.} Dans deux béchers, on verse respectivement \SI{50}{mL} de solution de sulfate de zinc (\SI{1,0}{mol.L^{-1}}) et \SI{50}{mL} de solution de sulfate de cuivre (\SI{1,0}{mol.L^{-1}}). On plonge une lame de zinc bien décapée dans le premier bécher, une lame de cuivre dans le second. On relie les deux béchers par un pont salin au \ce{KCl} gélifié, puis les deux lames par des fils à un galvanomètre (ou à une DEL avec sa résistance de protection).

\textbf{Observations.}
\begin{itemize}
\item L'aiguille du galvanomètre \textbf{dévie} (ou la DEL \textbf{brille}) : la pile débite un courant électrique dans le circuit extérieur.
\item Un voltmètre branché aux bornes indique une tension $U \approx \SI{1,10}{V}$, la borne COM (ou $-$) étant reliée au zinc et la borne V (ou $+$) au cuivre : la tension $U_{\text{Cu}} - U_{\text{Zn}} = +\SI{1,10}{V}$.
\item Après un fonctionnement prolongé : la lame de zinc \textbf{s'amincit} (sa masse diminue), tandis que la lame de cuivre \textbf{s'épaissit} d'un dépôt \textbf{rougeâtre} de cuivre ; la solution bleue pâlit.
\end{itemize}

\textbf{Interprétation.} La réaction spontanée $\ce{Zn + Cu^2+ -> Zn^2+ + Cu}$ a bien lieu, mais le transfert d'électrons s'effectue désormais \emph{à travers le circuit extérieur} : c'est ce flux d'électrons qui fait dévier l'aiguille ou briller la DEL.
\end{experience}

\begin{popfigure}
\begin{center}
\begin{tikzpicture}[scale=0.85, >=Stealth]
% --- Mélange direct (gauche) ---
\draw[thick, popdark] (-6.2,0) -- (-6.2,2.4) (-3.4,0) -- (-3.4,2.4);
\draw[thick, popdark] (-6.2,0) -- (-3.4,0);
\fill[popblue!40] (-6.1,0.05) rectangle (-3.5,1.8);
\fill[popmuted!60] (-5.1,0.05) rectangle (-4.5,2.9);
\node[font=\footnotesize\bfseries, popdark] at (-4.8,3.15) {Zn};
\node[font=\footnotesize, popblue!70!black] at (-5.0,1.0) {$\ce{Cu^2+}$};
\draw[->, poppink, thick] (-4.45,1.6) -- (-3.75,1.6) node[right, font=\footnotesize, poppink] {$e^-$ direct};
\draw[decorate, decoration={coil, aspect=0.5, amplitude=2pt}, poporange] (-5.9,2.0) -- (-5.9,2.7);
\node[font=\footnotesize, poporange!80!black] at (-5.9,2.95) {chaleur};
\node[font=\small\bfseries, align=center, popdark] at (-4.8,-0.9) {Mélange direct\\ \emph{énergie perdue en chaleur}};
% --- Pile (droite) ---
\draw[thick, popdark] (0.4,0) -- (0.4,2.0) (2.2,0) -- (2.2,2.0) (0.4,0) -- (2.2,0);
\fill[popblue!20] (0.5,0.05) rectangle (2.1,1.5);
\fill[popmuted!60] (1.15,0.8) rectangle (1.45,2.7);
\node[font=\footnotesize\bfseries, popdark] at (1.3,2.95) {Zn};
\draw[thick, popdark] (4.0,0) -- (4.0,2.0) (5.8,0) -- (5.8,2.0) (4.0,0) -- (5.8,0);
\fill[popblue!45] (4.1,0.05) rectangle (5.7,1.5);
\fill[poporange!80] (4.75,0.8) rectangle (5.05,2.7);
\node[font=\footnotesize\bfseries, poporange!80!black] at (4.9,2.95) {Cu};
\draw[line width=4pt, popgreen!60] (1.9,1.3) -- (1.9,2.2) .. controls (1.9,2.7) and (4.3,2.7) .. (4.3,2.2) -- (4.3,1.3);
\draw[thick, popdark] (1.3,2.7) -- (1.3,3.7) -- (2.8,3.7);
\draw[thick, popdark] (4.9,2.7) -- (4.9,3.7) -- (3.6,3.7);
\draw[thick, popgold!80!black, fill=poppinkL] (3.2,3.7) circle (0.35);
\node[font=\footnotesize, popdark] at (3.2,4.35) {DEL};
\draw[->, poppink, very thick] (1.7,4.0) -- (2.6,4.0) node[midway, above, font=\footnotesize\bfseries] {$e^-$};
\node[font=\small\bfseries, align=center, popdark] at (3.1,-0.9) {Pile Daniell\\ \emph{énergie récupérée en électricité}};
\end{tikzpicture}
\end{center}
\legende{Comparaison : dans le mélange direct, les électrons passent au contact et l'énergie se dissipe en chaleur ; dans la pile, la séparation des couples force les électrons à traverser le circuit extérieur, où leur énergie est récupérée.}
\end{popfigure}

\courssub{Détermination des pôles et sens de circulation}

\textbf{Au voltmètre.} On branche la borne « V » (ou $+$) du voltmètre sur l'électrode de cuivre et la borne « COM » (ou $-$) sur l'électrode de zinc : l'appareil affiche une valeur \emph{positive} ($U \approx +\SI{1,10}{V}$). On en déduit :

\begin{itemize}
\item l'électrode de \textbf{cuivre} est le \cle{pôle positif} de la pile ;
\item l'électrode de \textbf{zinc} est le \cle{pôle négatif}.
\end{itemize}

\textbf{Circulation des électrons.} Les électrons, chargés négativement, sont attirés par le pôle positif : ils \textbf{quittent la lame de zinc}, traversent le circuit extérieur (fils, voltmètre, DEL) et \textbf{arrivent à la lame de cuivre}, où ils sont captés par les ions \ce{Cu^2+}.

\textbf{Courant conventionnel.} Par convention historique, le courant électrique $I$ circule dans le circuit extérieur \textbf{en sens inverse} des électrons : il sort du pôle $+$ (cuivre) et rentre par le pôle $-$ (zinc).

\begin{attention}[Erreur fréquente : inverser électrons et courant]
Ne confonds jamais les deux sens ! Dans le \textbf{circuit extérieur} d'une pile :
\begin{itemize}
\item les \textbf{électrons} vont du pôle $-$ vers le pôle $+$ (ici : du zinc vers le cuivre) ;
\item le \textbf{courant conventionnel} $I$ va du pôle $+$ vers le pôle $-$ (ici : du cuivre vers le zinc).
\end{itemize}
Moyen mnémotechnique : « les électrons, négatifs, fuient le pôle négatif ». À l'intérieur des solutions et du pont salin, ce sont les \textbf{ions} qui assurent le passage du courant, jamais les électrons : un électron ne « nage » pas dans une solution !
\end{attention}

\courssub{Notation conventionnelle de la pile Daniell}

Pour représenter une pile de façon compacte et normalisée (norme IUPAC), on utilise un \textbf{diagramme de pile} (ou notation conventionnelle) :

\begin{center}
\textbf{($-$) $\quad$ \ce{Zn(s)} $|$ \ce{Zn^2+(aq)} $||$ \ce{Cu^2+(aq)} $|$ \ce{Cu(s)} $\quad$ ($+$)}
\end{center}

\begin{propriete}[Règles de la notation conventionnelle]
\begin{enumerate}
\item On écrit l'\textbf{anode} (pôle $-$ ) à \textbf{gauche} et la \textbf{cathode} (pôle $+$) à \textbf{droite}.
\item Une \textbf{barre verticale simple} $|$ marque une \textbf{frontière de phase} (contact métal/solution).
\item Un \textbf{double trait vertical} $||$ représente le \textbf{pont salin} (ou la paroi poreuse) entre les deux demi-piles.
\item Les \textbf{états physiques} (s, aq, l, g) sont indiqués entre parenthèses.
\item Les \textbf{concentrations} (ou activités) peuvent être précisées entre parenthèses après l'espèce, ex. \ce{Zn^2+(aq, 1,0 mol/L)}.
\end{enumerate}
\end{propriete}

\courssub{Évolution des électrodes : anode et cathode}

Que se passe-t-il concrètement sur chaque lame pendant le fonctionnement ?

\textbf{Sur la lame de zinc (pôle $-$).} Les atomes de zinc de la lame cèdent deux électrons et passent en solution sous forme d'ions \ce{Zn^2+} : c'est une \textbf{oxydation}. La lame se ronge, s'amincit ; sa masse diminue au fil du temps. On dit que le zinc est une électrode « consommée » ou « soluble ».

\[ \ce{Zn(s) -> Zn^2+(aq) + 2e-} \]

\textbf{Sur la lame de cuivre (pôle $+$).} Les ions \ce{Cu^2+} de la solution captent les électrons arrivant par le circuit et se déposent sous forme de cuivre métallique rougeâtre : c'est une \textbf{réduction}. La lame s'épaissit ; sa masse augmente.

\[ \ce{Cu^2+(aq) + 2e- -> Cu(s)} \]

\begin{definition}[Anode et cathode]
L'\cle{anode} est l'électrode où se produit l'\textbf{oxydation} ; dans une pile (générateur), c'est le \textbf{pôle négatif} (elle fournit des électrons au circuit).

La \cle{cathode} est l'électrode où se produit la \textbf{réduction} ; dans une pile, c'est le \textbf{pôle positif}.

Moyen mnémotechnique : \textbf{A}node — oxyd\textbf{A}tion (voyelles) ; \textbf{C}athode — rédu\textbf{C}tion (consonnes). Dans la pile Daniell : anode = zinc (pôle $-$), cathode = cuivre (pôle $+$).
\end{definition}

\begin{exoresolu}[Masse des électrodes après fonctionnement]
\textbf{Énoncé.} Une pile Daniell débite un courant d'intensité constante $I = \SI{50}{mA}$ pendant une durée $\Delta t = \SI{2}{h}$. On donne $M(\ce{Zn}) = \SI{65,4}{g.mol^{-1}}$, $M(\ce{Cu}) = \SI{63,5}{g.mol^{-1}}$ et la charge d'une mole d'électrons (faraday) : $\mathcal{F} = \SI{96500}{C.mol^{-1}}$. Calculer la variation de masse de chaque électrode.

\tcblower
\textbf{Solution détaillée.}

\textbf{Étape 1 : quantité d'électricité échangée.}
\[ Q = I \times \Delta t = 0{,}050 \times (2 \times 3600) = \SI{360}{C} \]

\textbf{Étape 2 : quantité de matière d'électrons.}
\[ n(e^-) = \frac{Q}{\mathcal{F}} = \frac{360}{96500} \approx \SI{3,73e-3}{mol} \]

\textbf{Étape 3 : électrode de zinc (anode).} D'après $\ce{Zn -> Zn^2+ + 2e-}$, il faut 2 moles d'électrons par mole de zinc dissous :
\[ n(\ce{Zn}) = \frac{n(e^-)}{2} \approx \SI{1,86e-3}{mol} \]
\[ \Delta m(\ce{Zn}) = -n(\ce{Zn}) \times M(\ce{Zn}) = -1{,}86\times 10^{-3} \times 65{,}4 \approx \SI{-0,12}{g} \]
La lame de zinc a \textbf{perdu environ \SI{0,12}{g}} : elle s'est amincie.

\textbf{Étape 4 : électrode de cuivre (cathode).} D'après $\ce{Cu^2+ + 2e- -> Cu}$ :
\[ n(\ce{Cu}) = \frac{n(e^-)}{2} \approx \SI{1,86e-3}{mol} \]
\[ \Delta m(\ce{Cu}) = +n(\ce{Cu}) \times M(\ce{Cu}) = 1{,}86\times 10^{-3} \times 63{,}5 \approx \SI{+0,12}{g} \]
La lame de cuivre a \textbf{gagné environ \SI{0,12}{g}} de dépôt rougeâtre. Remarque : les deux variations de masse sont presque égales en valeur absolue car les masses molaires de \ce{Zn} et \ce{Cu} sont voisines, mais ce n'est qu'une coïncidence numérique !
\end{exoresolu}

\begin{savaistu}[John Daniell et les télégraphes]
Le chimiste britannique \textbf{John Frederic Daniell} (1790--1845) a inventé cette pile en \textbf{1836} pour remédier aux défauts de la pile de Volta, dont la tension chutait rapidement. Sa pile, robuste et durable, a alimenté les premières lignes de \textbf{télégraphe} électrique qui ont révolutionné les communications au XIX\textsuperscript{e} siècle : c'est grâce à elle que des messages ont pu traverser continents et océans en quelques minutes ! Le principe — deux demi-piles, un transfert d'électrons canalisé — est encore aujourd'hui celui des piles de nos lampes torches, de nos télécommandes et des batteries qui équipent les centres de santé en zone rurale non électrifiée. Quand ta pile est « usée », c'est tout simplement que l'un des réactifs (le réducteur de l'anode ou l'oxydant de la cathode) a été entièrement consommé : la réaction d'oxydoréduction ne peut plus avoir lieu.
\end{savaistu}

\begin{aretenir}[L'essentiel de la section]
\begin{itemize}
\item La réaction $\ce{Zn + Cu^2+ -> Zn^2+ + Cu}$ est spontanée ; en mélange direct, son énergie se perd en chaleur.
\item En séparant les deux couples en \textbf{deux demi-piles} reliées par un circuit extérieur et un \textbf{pont salin}, on force les électrons à circuler dans le circuit : c'est une \textbf{pile}.
\item Dans la pile Daniell : le \textbf{zinc} est l'\textbf{anode} (pôle $-$, oxydation, la lame s'amincit) ; le \textbf{cuivre} est la \textbf{cathode} (pôle $+$, réduction, dépôt rougeâtre).
\item Les électrons circulent du pôle $-$ vers le pôle $+$ dans le circuit extérieur ; le courant conventionnel circule en sens inverse ; les ions migrent dans le pont salin pour assurer l'électroneutralité.
\item Notation conventionnelle : \ce{(-) Zn(s) | Zn^2+(aq) || Cu^2+(aq) | Cu(s) (+)}.
\end{itemize}
\end{aretenir}

\coursec{Fonctionnement, équations et notation conventionnelle}

Dans la section précédente, nous avons réalisé la pile Daniell et observé ses effets : l'aiguille du galvanomètre dévie, la lame de zinc s'amincit, un dépôt rouge de cuivre apparaît sur l'autre lame. Nous avons constaté \emph{que} la pile fonctionne. Il nous reste maintenant à comprendre \emph{comment} elle fonctionne : quelles réactions chimiques se déroulent à chaque électrode, comment écrire le bilan global, et comment les chimistes notent une pile de façon universelle, sans dessin. C'est l'objet de cette section.

\courssub{Les deux demi-équations : le moteur chimique de la pile}

\textbf{L'idée intuitive.} Une pile, c'est une réaction d'oxydoréduction spontanée que l'on a \cle{séparée en deux} : l'oxydation d'un côté, la réduction de l'autre. Les électrons, au lieu de passer directement du réducteur à l'oxydant, sont obligés de voyager dans le circuit extérieur — et c'est ce courant que l'on récolte.

\begin{definition}[Anode et cathode]
Dans une pile en fonctionnement (mode décharge) :
\begin{itemize}
\item l'\cle{anode} est l'électrode où se produit l'\textbf{oxydation} (perte d'électrons) ; c'est le \textbf{pôle négatif} $(-)$ de la pile, car elle libère des électrons dans le circuit extérieur ;
\item la \cle{cathode} est l'électrode où se produit la \textbf{réduction} (gain d'électrons) ; c'est le \textbf{pôle positif} $(+)$.
\end{itemize}
Moyen mnémotechnique : \emph{anode} et \emph{oxydation} commencent par une voyelle ; \emph{cathode} et \emph{réduction} commencent par une consonne.
\end{definition}

\textbf{À l'anode : le zinc s'oxyde.} Le zinc métallique est le réducteur du couple \ce{Zn^2+}/\ce{Zn}. Il cède deux électrons et passe en solution sous forme d'ions \ce{Zn^2+} :

\[ \ce{Zn (s) -> Zn^2+ (aq) + 2 e-} \]

C'est une \textbf{oxydation} (perte d'électrons, augmentation du nombre d'oxydation de 0 à +II). Conséquence visible : la lame de zinc se consume progressivement, et la concentration en ions \ce{Zn^2+} augmente dans le bécher.

\textbf{À la cathode : les ions cuivre se réduisent.} L'ion cuivre(II) est l'oxydant du couple \ce{Cu^2+}/\ce{Cu}. Il capte les deux électrons arrivés par le circuit et se dépose sous forme de cuivre métallique :

\[ \ce{Cu^2+ (aq) + 2 e- -> Cu (s)} \]

C'est une \textbf{réduction} (gain d'électrons, diminution du nombre d'oxydation de +II à 0). Conséquence visible : un dépôt rougeâtre de cuivre épaissit la lame, et la couleur bleue de la solution de sulfate de cuivre pâlit.

\begin{attention}[Ne pas inverser les pôles]
Une erreur classique consiste à dire « le zinc est le pôle $+$ car il donne des électrons ». C'est exactement l'inverse : en \textbf{donnant} des électrons, le zinc devient la \textbf{source} d'électrons du circuit, donc le \textbf{pôle négatif} $(-)$. Les électrons circulent dans le circuit extérieur du pôle $(-)$ vers le pôle $(+)$, tandis que le courant conventionnel $I$ circule du $(+)$ vers le $(-)$.
\end{attention}

\begin{popfigure}
\begin{center}
\begin{tikzpicture}[scale=0.95, every node/.style={font=\small}]
% Bécher gauche (zinc)
\draw[thick, popdark] (0,0) -- (0,3) (0,0) -- (3,0) (3,0) -- (3,3);
\fill[popblue!20] (0.1,0.1) rectangle (2.9,1.9);
\draw[popblue!60] (0.1,1.9) -- (2.9,1.9);
\node[popblue!80!black, align=center] at (1.5,0.5) {\ce{ZnSO4 (aq)}\\ $\ce{Zn^2+}$, $\ce{SO4^2-}$};
% Lame de zinc
\fill[popmuted!70] (1.25,0.6) rectangle (1.75,3.8);
\node[rotate=90, white, font=\footnotesize\bfseries] at (1.5,2.2) {Zinc};
% Bécher droit (cuivre)
\draw[thick, popdark] (7,0) -- (7,3) (7,0) -- (10,0) (10,0) -- (10,3);
\fill[popgreen!20] (7.1,0.1) rectangle (9.9,1.9);
\draw[popgreen!60] (7.1,1.9) -- (9.9,1.9);
\node[popgreen!50!black, align=center] at (8.5,0.5) {\ce{CuSO4 (aq)}\\ $\ce{Cu^2+}$, $\ce{SO4^2-}$};
% Lame de cuivre
\fill[poporange!80] (8.25,0.6) rectangle (8.75,3.8);
\node[rotate=90, white, font=\footnotesize\bfseries] at (8.5,2.2) {Cuivre};
% Pont salin
\draw[line width=5pt, popgold!70] (2.6,2.5) -- (2.6,4.2) -- (7.4,4.2) -- (7.4,2.5);
\node[popgold!60!black, above] at (5,4.2) {pont salin (\ce{K+}, \ce{NO3-})};
% Fils et ampèremètre
\draw[thick, popdark] (1.5,3.8) -- (1.5,5.0) -- (4.3,5.0);
\draw[thick, popdark] (8.5,3.8) -- (8.5,5.0) -- (5.7,5.0);
\draw[thick, popdark, fill=white] (5,5.0) circle (0.45);
\node at (5,5.0) {\textbf{A}};
% Sens du courant
\draw[-{Stealth[length=3mm]}, very thick, poppink] (6.3,5.6) -- (4.1,5.6);
\node[poppink, above] at (5.2,5.6) {courant $I$};
\draw[-{Stealth[length=3mm]}, very thick, poppurple] (3.5,6.2) -- (5.9,6.2);
\node[poppurple, above] at (4.7,6.2) {électrons $e^-$};
% Pôles
\node[font=\Large\bfseries, poppink] at (1.5,5.4) {$-$};
\node[font=\Large\bfseries, poppink] at (8.5,5.4) {$+$};
% Demi-équations
\node[align=center, popdark] at (1.5,-0.8) {\textbf{Anode} $(-)$ : oxydation\\ $\ce{Zn -> Zn^2+ + 2 e-}$};
\node[align=center, popdark] at (8.5,-0.8) {\textbf{Cathode} $(+)$ : réduction\\ $\ce{Cu^2+ + 2 e- -> Cu}$};
% Bilan
\node[align=center, font=\small\bfseries, popink] at (5,-2.0) {Bilan : $\ce{Zn + Cu^2+ -> Zn^2+ + Cu}$};
\end{tikzpicture}
\end{center}
\legende{Pile Daniell en fonctionnement : à chaque électrode, une demi-équation ; dans le circuit extérieur, les électrons vont du zinc $(-)$ vers le cuivre $(+)$.}
\end{popfigure}

\courssub{L'équation-bilan de fonctionnement}

Pour obtenir l'équation-bilan, on \textbf{additionne} les deux demi-équations après s'être assuré que le nombre d'électrons perdus à l'anode égale le nombre d'électrons gagnés à la cathode. Ici, c'est déjà le cas (2 électrons de chaque côté) :

\begin{align*}
\text{Anode : } & \ce{Zn -> Zn^2+ + 2 e-}\\
\text{Cathode : } & \ce{Cu^2+ + 2 e- -> Cu}\\ \hline
\text{Bilan : } & \ce{Zn + Cu^2+ -> Zn^2+ + Cu}
\end{align*}

\begin{aretenir}[Règle d'or]
Dans l'\cle{équation-bilan} de fonctionnement d'une pile, les électrons \textbf{n'apparaissent jamais} : ils s'éliminent par addition des demi-équations. Les électrons ne circulent que dans les fils et les électrodes ; dans les solutions, ce sont les \textbf{ions} qui assurent le passage du courant (d'où le rôle indispensable du pont salin).
\end{aretenir}

\begin{methode}[Écrire l'équation-bilan de fonctionnement d'une pile]
\begin{enumerate}
\item \textbf{Identifier les deux couples} mis en jeu (ici \ce{Zn^2+}/\ce{Zn} et \ce{Cu^2+}/\ce{Cu}).
\item \textbf{Écrire chaque demi-équation dans le bon sens} : oxydation à l'anode (le réducteur cède des électrons), réduction à la cathode (l'oxydant capte des électrons).
\item \textbf{Équilibrer les électrons échangés} : multiplier si nécessaire une ou les deux demi-équations pour que le nombre d'électrons cédés égale le nombre d'électrons captés.
\item \textbf{Additionner} les demi-équations et \textbf{simplifier} les électrons, qui doivent disparaître du bilan.
\end{enumerate}
\end{methode}

\begin{attention}[Erreur fréquente au Probatoire]
Beaucoup d'élèves écrivent un bilan du type $\ce{Zn + Cu^2+ -> Zn^2+ + Cu + 2 e-}$ : c'est \textbf{faux}. Si des électrons subsistent dans ton bilan, c'est que les demi-équations n'ont pas été correctement multipliées avant l'addition. Vérifie toujours : le bilan ne doit contenir \textbf{aucun} électron, et il doit être équilibré en atomes \emph{et} en charges.
\end{attention}

\textbf{La règle générale de prévision.} Comment savoir, sans faire l'expérience, quel métal s'oxyde ? Nous verrons en section 4 une classification quantitative fondée sur les potentiels. Retenons dès maintenant la règle qualitative :

\begin{propriete}[Règle de prévision du sens de fonctionnement]
Quand on associe deux couples oxydant/réducteur dans une pile :
\begin{itemize}
\item le \textbf{réducteur du couple de potentiel standard le plus faible} (le métal le plus « réducteur », ici le zinc) \textbf{s'oxyde} : il constitue l'anode, pôle $(-)$ ;
\item l'\textbf{oxydant de l'autre couple} (ici \ce{Cu^2+}) \textbf{se réduit} : il intervient à la cathode, pôle $(+)$.
\end{itemize}
Le zinc est plus réducteur que le cuivre : c'est donc lui qui « sacrifie » ses électrons.
\end{propriete}

\courssub{La notation conventionnelle d'une pile}

Dessiner la pile à chaque fois serait fastidieux. Les chimistes ont adopté une écriture symbolique, universelle, qui résume toute la pile en une ligne.

\begin{definition}[Notation conventionnelle]
Une pile se note en écrivant, \textbf{de gauche à droite} : l'anode (pôle $-$), puis une barre verticale pour la séparation électrode/solution, puis l'ion de la demi-pile anodique, puis un \textbf{double trait} symbolisant le pont salin, puis l'ion de la demi-pile cathodique, une barre, et enfin la cathode (pôle $+$). Pour la pile Daniell :
\[ (-)\ \ce{Zn} \;|\; \ce{Zn^2+} \;\|\; \ce{Cu^2+} \;|\; \ce{Cu}\ (+) \]
On peut préciser les concentrations (conditions standards : \SI{1}{mol.L^{-1}}) :
\[ (-)\ \ce{Zn} \,|\, \ce{Zn^2+} (\SI{1}{mol.L^{-1}}) \,\|\, \ce{Cu^2+} (\SI{1}{mol.L^{-1}}) \,|\, \ce{Cu}\ (+) \]
\end{definition}

\begin{popfigure}
\begin{center}
\begin{tikzpicture}[every node/.style={font=\small}]
% Schéma réel simplifié (haut)
\fill[popmuted!70] (0.6,1.2) rectangle (1.0,3.0);
\fill[poporange!80] (6.0,1.2) rectangle (6.4,3.0);
\draw[thick, popdark] (0.2,0.4) -- (0.2,2.4) (0.2,0.4) -- (3.0,0.4) (3.0,0.4) -- (3.0,2.4);
\draw[thick, popdark] (4.0,0.4) -- (4.0,2.4) (4.0,0.4) -- (6.8,0.4) (6.8,0.4) -- (6.8,2.4);
\draw[line width=3pt, popgold!70] (2.6,2.0) -- (2.6,3.2) -- (4.4,3.2) -- (4.4,2.0);
\draw[thick, popdark] (0.8,3.0) -- (0.8,3.8) -- (6.2,3.8) -- (6.2,3.0);
\node[font=\bfseries, poppink] at (0.8,4.15) {$-$};
\node[font=\bfseries, poppink] at (6.2,4.15) {$+$};
% Flèches de correspondance
\draw[-{Stealth}, thick, popblue] (0.8,0.1) -- (0.8,-0.9);
\draw[-{Stealth}, thick, popblue] (1.6,0.1) -- (1.9,-0.9);
\draw[-{Stealth}, thick, popblue] (3.5,0.1) -- (3.4,-0.9);
\draw[-{Stealth}, thick, popblue] (5.0,0.1) -- (4.9,-0.9);
\draw[-{Stealth}, thick, popblue] (6.2,0.1) -- (6.2,-0.9);
% Notation (bas)
\node[font=\large] at (3.5,-1.4) {$(-)\ \ce{Zn} \;\Big|\; \ce{Zn^2+} \;\Big\|\; \ce{Cu^2+} \;\Big|\; \ce{Cu}\ (+)$};
\node[font=\footnotesize, popmuted, align=center] at (0.8,-2.15) {anode\\ (métal)};
\node[font=\footnotesize, popmuted, align=center] at (1.9,-2.15) {ion de la\\ demi-pile $(-)$};
\node[font=\footnotesize, popmuted, align=center] at (3.4,-2.15) {pont\\ salin};
\node[font=\footnotesize, popmuted, align=center] at (4.9,-2.15) {ion de la\\ demi-pile $(+)$};
\node[font=\footnotesize, popmuted, align=center] at (6.2,-2.15) {cathode\\ (métal)};
\end{tikzpicture}
\end{center}
\legende{Correspondance entre le schéma réel de la pile et sa notation conventionnelle : on « lit » la pile de gauche (anode) vers la droite (cathode), le double trait figurant le pont salin.}
\end{popfigure}

\textbf{Lire une notation conventionnelle.} Cette écriture se lit comme une phrase et donne immédiatement toutes les informations :

\begin{itemize}
\item \textbf{les couples} : à gauche \ce{Zn^2+}/\ce{Zn}, à droite \ce{Cu^2+}/\ce{Cu} ;
\item \textbf{les pôles} : par convention, l'anode (pôle $-$) est toujours écrite \textbf{à gauche}, la cathode (pôle $+$) \textbf{à droite} ;
\item \textbf{le sens du courant} : les électrons partent de l'électrode écrite à gauche et rejoignent celle de droite par le circuit extérieur ; le courant conventionnel circule donc de la droite vers la gauche dans le circuit extérieur ;
\item \textbf{les demi-équations} : oxydation du métal de gauche, réduction de l'ion de droite.
\end{itemize}

\begin{exemplebox}[Lecture d'une notation]
On donne la pile : $(-)\ \ce{Fe} \,|\, \ce{Fe^2+} \,\|\, \ce{Ag+} \,|\, \ce{Ag}\ (+)$. Sans aucun schéma, on en déduit :
\begin{itemize}
\item couples : \ce{Fe^2+}/\ce{Fe} et \ce{Ag+}/\ce{Ag} ;
\item anode $(-)$ : le fer, qui s'oxyde : $\ce{Fe -> Fe^2+ + 2 e-}$ ;
\item cathode $(+)$ : l'argent, où $\ce{Ag+}$ se réduit : $\ce{Ag+ + e- -> Ag}$ ;
\item bilan (après multiplication de la deuxième demi-équation par 2) : $\ce{Fe + 2 Ag+ -> Fe^2+ + 2 Ag}$.
\end{itemize}
\end{exemplebox}

\courssub{Caractéristiques d'une pile : f.é.m., polarité, capacité et usure}

\begin{definition}[Caractéristiques d'une pile]
Une pile est caractérisée par :
\begin{itemize}
\item sa \cle{force électromotrice} (f.é.m.) $E$, tension mesurée à circuit ouvert (courant quasi nul) entre ses bornes, exprimée en volts (V) ; pour la pile Daniell dans les conditions standards, $E^\circ \approx \SI{1,10}{V}$ ;
\item sa \cle{polarité} : la nature de ses pôles $(+)$ et $(-)$, déterminée par la nature des couples ;
\item sa \cle{capacité} : la quantité d'électricité maximale $Q = I \times t$ qu'elle peut débiter avant d'être usée, souvent exprimée en ampère-heures (A$\cdot$h).
\end{itemize}
\end{definition}

\begin{exemplebox}[Quantité d'électricité débitée par une pile Daniell]
Une pile Daniell débite un courant d'intensité $I = \SI{0,20}{A}$ pendant $t = \SI{10}{h}$. La quantité d'électricité mise en jeu est :
\[ Q = I \times t = 0,20 \times (10 \times 3600) = \SI{7200}{C} \]
soit $Q = 7,2 \times 10^3\,\text{C}$. En moles d'électrons (avec $F \approx \SI{96500}{C.mol^{-1}}$) :
\[ n(e^-) = \frac{Q}{F} = \frac{7200}{96500} \approx \SI{7,46e-2}{mol} \]
Ce type de calcul, très classique au Probatoire, relie le fonctionnement électrique de la pile à l'avancement de sa réaction chimique.
\end{exemplebox}

\textbf{L'usure de la pile.} Une pile n'est pas éternelle : sa réaction de fonctionnement \textbf{consomme des réactifs}. La réaction $\ce{Zn + Cu^2+ -> Zn^2+ + Cu}$ s'arrête dès que l'un des deux réactifs est épuisé :

\begin{itemize}
\item soit la \textbf{lame de zinc} est entièrement consumée (dissoute en ions \ce{Zn^2+}) ;
\item soit les \textbf{ions \ce{Cu^2+}} sont tous réduits (la solution bleue devient incolore).
\end{itemize}

La pile est alors dite \textbf{« à plat »} : sa f.é.m. s'effondre et elle ne peut plus débiter de courant. Le réactif épuisé en premier est le \textbf{réactif limitant} de la réaction de fonctionnement — tout se passe comme dans un dosage stœchiométrique.

\begin{exoresolu}[Durée de vie d'une pile Daniell]
\textbf{Énoncé.} Une pile Daniell contient une lame de zinc de masse $m = \SI{6,54}{g}$ et un large excès d'ions \ce{Cu^2+}. Elle débite un courant constant $I = \SI{0,10}{A}$. On donne $M(\ce{Zn}) = \SI{65,4}{g.mol^{-1}}$ et $F = \SI{96500}{C.mol^{-1}}$. Calculer la durée de fonctionnement de la pile.
\tcblower
\textbf{Solution.} Le zinc est le réactif limitant (les ions \ce{Cu^2+} sont en excès).

\textbf{Étape 1 — Quantité de zinc disponible :}
\[ n(\ce{Zn}) = \frac{m}{M} = \frac{6,54}{65,4} = \SI{0,100}{mol} \]

\textbf{Étape 2 — Quantité d'électrons libérée.} D'après la demi-équation $\ce{Zn -> Zn^2+ + 2 e-}$, chaque mole de zinc libère 2 moles d'électrons :
\[ n(e^-) = 2 \times n(\ce{Zn}) = \SI{0,200}{mol} \]

\textbf{Étape 3 — Quantité d'électricité totale :}
\[ Q = n(e^-) \times F = 0,200 \times 96500 = \SI{19300}{C} \]

\textbf{Étape 4 — Durée de fonctionnement :}
\[ t = \frac{Q}{I} = \frac{19300}{0,10} = \SI{193000}{s} \approx \SI{53,6}{h} \]

La pile peut fonctionner environ \textbf{54 heures} avant que sa lame de zinc ne soit entièrement consumée.
\end{exoresolu}

\begin{savaistu}[Pourquoi les vieilles piles fuient-elles ?]
Dans une pile saline usagée (les piles cylindriques classiques), c'est presque toujours l'\textbf{électrode de zinc} — qui sert aussi de boîtier extérieur — qui a été entièrement consommée par l'oxydation $\ce{Zn -> Zn^2+ + 2 e-}$. Le boîtier percé laisse alors s'échapper l'électrolyte, une pâte \textbf{corrosive} qui endommage les appareils (télécommandes, lampes torches). C'est pourquoi il faut retirer les piles d'un appareil qu'on n'utilise plus, et ne jamais les jeter n'importe où : elles contiennent des métaux lourds polluants.
\end{savaistu}

\begin{aretenir}[L'essentiel de la section]
\begin{itemize}
\item À l'\textbf{anode} $(-)$, oxydation du zinc : $\ce{Zn -> Zn^2+ + 2 e-}$ ; à la \textbf{cathode} $(+)$, réduction des ions cuivre : $\ce{Cu^2+ + 2 e- -> Cu}$.
\item Équation-bilan : $\ce{Zn + Cu^2+ -> Zn^2+ + Cu}$ — \textbf{sans électrons}.
\item Le réducteur du couple de potentiel le plus faible s'oxyde ; l'oxydant de l'autre couple se réduit.
\item Notation conventionnelle : $(-)\ \ce{Zn} \,|\, \ce{Zn^2+} \,\|\, \ce{Cu^2+} \,|\, \ce{Cu}\ (+)$, anode à gauche, cathode à droite, double trait pour le pont salin.
\item La pile Daniell a une f.é.m. $E^\circ \approx \SI{1,10}{V}$ ; elle est « à plat » quand un réactif est épuisé (réactif limitant).
\end{itemize}
\end{aretenir}

\coursec{La demi-pile de référence : l'électrode à hydrogène}

Dans la section précédente, nous avons appris à décrire le fonctionnement de la pile Daniell : le zinc s'oxyde au pôle négatif, le cuivre se réduit au pôle positif, et la f.é.m. mesurée vaut \SI{1,10}{V}. Nous savons donc dire \emph{qui} est le pôle $+$ et \emph{qui} est le pôle $-$ dans un couple donné. Mais une question plus ambitieuse se pose maintenant : peut-on attribuer à \emph{chaque couple} une valeur numérique qui traduise, à elle seule, son « pouvoir oxydant » ou son « pouvoir réducteur » ? Si oui, il suffirait de comparer deux nombres pour prévoir le sens de n'importe quelle réaction d'oxydoréduction, sans avoir à réaliser la pile. C'est exactement l'objectif de cette section : découvrir la référence universelle qui rend cette comparaison possible.

\courssub{Le problème de la mesure : pourquoi faut-il une référence ?}

Imagine que l'on te demande : « Quelle est l'altitude de Yaoundé ? » Tu réponds : environ \SI{750}{m}. Mais \SI{750}{m} \emph{par rapport à quoi} ? Par rapport au niveau de la mer, bien sûr. Or le niveau de la mer n'est pas une vérité absolue : c'est un \textbf{zéro choisi par convention}, un étalon accepté par tous. Sans ce zéro commun, chaque ville donnerait son altitude par rapport à son propre point de référence et aucune comparaison ne serait possible.

En électrochimie, le problème est identique. Un voltmètre ne mesure jamais le potentiel d'une électrode seule : il mesure une \emph{différence de potentiel} entre \emph{deux} électrodes. Dire « le potentiel du couple \ce{Cu^2+}/\ce{Cu} » n'a donc aucun sens en soi, exactement comme dire « l'altitude de Yaoundé » sans préciser le zéro. Pour construire une échelle de potentiels, il faut :

\begin{itemize}
\item choisir une \cle{demi-pile de référence}, jouant le rôle du « niveau de la mer » ;
\item lui attribuer \emph{par convention} la valeur zéro ;
\item mesurer tous les autres couples \emph{par rapport à elle}.
\end{itemize}

La communauté scientifique internationale a choisi pour ce rôle un couple un peu particulier : le couple de l'hydrogène, \ce{H+}/\ce{H2}. Voyons comment la demi-pile correspondante est construite.

\courssub{Constitution de l'électrode standard à hydrogène (ESH)}

\begin{definition}[Électrode standard à hydrogène]
L'\cle{électrode standard à hydrogène} (ESH) est la demi-pile de référence constituée d'un fil de \textbf{platine platiné} (recouvert de noir de platine) plongeant dans une solution d'ions \ce{H+} de concentration \SI{1}{mol.L^{-1}}, sur lequel barbote du dihydrogène \ce{H2} sous la pression de \SI{1}{bar}, à la température de \SI{25}{\celsius}. Par convention internationale, son potentiel est fixé à :
\[ E^\circ(\ce{H+}/\ce{H2}) = \SI{0,00}{V} \quad \text{à toute température.} \]
\end{definition}

Décortiquons cette constitution, car chaque détail a une raison d'être.

\textbf{Le couple mis en jeu.} C'est le couple \ce{H+}/\ce{H2}, dont la demi-équation électronique s'écrit :
\[ \ce{2H+ + 2e- <=> H2} \]
Dans le sens de la réduction (de gauche à droite), deux ions hydrogène captent deux électrons pour former une molécule de dihydrogène ; dans le sens de l'oxydation (de droite à gauche), le dihydrogène cède deux électrons et libère deux ions \ce{H+}. L'équilibre entre ces deux formes s'établit \emph{à la surface du platine}.

\textbf{Pourquoi le platine ?} Le platine joue un double rôle. D'abord, il est \emph{conducteur} : il assure le passage des électrons vers le circuit extérieur. Ensuite, il est \emph{chimiquement inerte} : il ne s'oxyde pas et ne participe pas à la réaction. Mais il possède une propriété remarquable : sa surface \emph{catalyse} l'équilibre $\ce{2H+ + 2e- <=> H2}$, c'est-à-dire qu'elle permet à cet équilibre de s'établir rapidement. Le noir de platine déposé sur le fil (« platine platiné ») augmente énormément la surface de contact, ce qui rend la catalyse encore plus efficace.

\textbf{Pourquoi ces conditions précises ?} Le potentiel d'une demi-pile dépend des concentrations et des pressions. Pour que la référence soit universelle et reproductible dans tous les laboratoires du monde, il faut fixer des \cle{conditions standard} : concentration en \ce{H+} égale à \SI{1}{mol.L^{-1}}, pression de \ce{H2} égale à \SI{1}{bar}, température de \SI{25}{\celsius}. C'est pourquoi on parle d'électrode \emph{standard} à hydrogène.

\begin{experience}[Description de l'ESH]
\textbf{Dispositif.} Un tube de verre contient une solution acide (par exemple de l'acide chlorhydrique) de concentration \SI{1}{mol.L^{-1}} en ions \ce{H+}. Un fil de platine platiné, soudé à un fil de connexion, plonge dans la solution. Un courant de gaz dihydrogène, maintenu à la pression de \SI{1}{bar}, arrive par un tube latéral et barbote à la surface du platine ; l'excès de gaz s'échappe par une ouverture. L'ensemble est thermostatisé à \SI{25}{\celsius}.

\textbf{Observations.} Des bulles de \ce{H2} se forment en continu à la surface du platine. Reliée à un voltmètre, l'électrode présente un potentiel stable et parfaitement reproductible.

\textbf{Interprétation.} À la surface du platine, l'équilibre $\ce{2H+ + 2e- <=> H2}$ s'établit : les molécules de \ce{H2} adsorbées sur le platine et les ions \ce{H+} de la solution échangent des électrons en permanence. Le platine sert de « réservoir » d'électrons et de siège de la réaction, sans jamais être consommé.
\end{experience}

\begin{popfigure}
\begin{center}
\begin{tikzpicture}[scale=0.95, every node/.style={font=\small}]
% Solution
\fill[popblueL] (-1.6,0) rectangle (1.6,3.2);
% Paroi du tube de verre
\draw[line width=1.2pt, popdark] (-1.6,-0.6) -- (-1.6,4.6);
\draw[line width=1.2pt, popdark] (1.6,-0.6) -- (1.6,4.6);
\draw[line width=1.2pt, popdark] (-1.6,-0.6) -- (1.6,-0.6);
% Niveau de la solution
\draw[popblue, dashed] (-1.6,3.2) -- (1.6,3.2);
\node[popblue, anchor=west] at (1.7,3.2) {solution de \ce{H+} à \SI{1}{mol.L^{-1}}};
% Tube d'arrivée de H2
\draw[line width=1pt, popdark] (-3.4,2.2) -- (-0.5,2.2) -- (-0.5,0.4);
\draw[line width=1pt, popdark] (-3.4,1.6) -- (-0.9,1.6) -- (-0.9,0.4);
\draw[-{Stealth}, popgreen, line width=1.2pt] (-4.4,1.9) -- (-3.5,1.9);
\node[popgreen, anchor=east] at (-4.5,1.9) {\ce{H2} (\SI{1}{bar})};
% Bulles
\foreach \x/\y/\r in {-0.7/0.7/0.09,-0.65/1.1/0.07,-0.75/1.5/0.08,-0.6/1.9/0.06,-0.7/2.3/0.08,-0.62/2.7/0.06,-0.68/3.0/0.05}{
  \draw[popblue!70!popdark] (\x,\y) circle (\r);}
% Plaque de platine
\fill[popdark] (0.35,0.2) rectangle (0.75,1.8);
\draw[poppink, line width=0.8pt, decorate, decoration={random steps, segment length=2pt, amplitude=0.6pt}] (0.35,0.2) rectangle (0.75,1.8);
\node[anchor=west, popdark] at (0.85,0.7) {platine platiné};
% Fil de connexion
\draw[line width=1pt, popdark] (0.55,1.8) -- (0.55,4.4) -- (2.6,4.4);
\fill[popgold] (2.6,4.4) circle (0.09);
\node[anchor=west] at (2.75,4.4) {vers le voltmètre};
% Sortie de gaz
\draw[line width=1pt, popdark] (1.0,4.6) -- (1.0,5.2) -- (1.8,5.2);
\draw[-{Stealth}, popgreen] (1.8,5.2) -- (2.6,5.2);
\node[popgreen, anchor=west] at (2.65,5.2) {échappement de \ce{H2}};
% Température
\node[popmuted, anchor=east] at (-1.7,-0.3) {$T = \SI{25}{\celsius}$};
\end{tikzpicture}
\end{center}
\legende{Schéma de l'électrode standard à hydrogène (ESH) : le dihydrogène barbote à la surface du platine platiné plongeant dans une solution d'ions \ce{H+} à \SI{1}{mol.L^{-1}}, à \SI{25}{\celsius}.}
\end{popfigure}

\begin{savaistu}[Pourquoi le platine, ce métal précieux ?]
Le platine est l'un des métaux les plus inertes qui soient : il résiste à presque tous les acides et ne s'oxyde pas, même à chaud. C'est cette inertie qui le rend irremplaçable ici : tout autre métal moins noble (fer, zinc, cuivre\dots) réagirait avec les ions \ce{H+} et fausserait la mesure. De plus, le platine est un excellent \emph{catalyseur} : sa surface adsorbe les molécules de \ce{H2} et facilite l'échange électronique $\ce{2H+ + 2e- <=> H2}$. C'est d'ailleurs ce même pouvoir catalytique qui fait du platine un composant clé des pots catalytiques des automobiles et des piles à combustible à hydrogène, technologie d'avenir pour la production d'électricité propre.
\end{savaistu}

\courssub{La convention du zéro : une référence arbitraire}

\begin{propriete}[Convention internationale]
Par convention internationale, le potentiel standard du couple \ce{H+}/\ce{H2} est fixé à :
\[ E^\circ(\ce{H+}/\ce{H2}) = \SI{0,00}{V} \quad \text{à toute température.} \]
Tous les potentiels d'oxydoréduction des autres couples sont mesurés \emph{par rapport à cette référence}.
\end{propriete}

\begin{attention}[Le zéro de l'ESH n'est pas une mesure !]
Erreur très fréquente : croire que le potentiel de l'ESH a été \emph{mesuré} et qu'on a « trouvé » \SI{0,00}{V}. C'est faux ! On ne peut pas mesurer le potentiel d'une électrode seule. La valeur \SI{0,00}{V} est une \textbf{convention arbitraire}, un choix collectif, exactement comme le niveau de la mer pour les altitudes ou le point de fusion de la glace pour l'échelle Celsius. Le Mont Cameroun culmine à environ \SI{4100}{m} \emph{par rapport au niveau de la mer} ; de même, le couple \ce{Cu^2+}/\ce{Cu} a un potentiel de \SI{+0,34}{V} \emph{par rapport à l'ESH}. Retenir : seules les \emph{différences} de potentiel ont un sens physique.
\end{attention}

\courssub{Principe de la mesure d'un potentiel standard}

Comment mesurer concrètement le potentiel d'un couple, par exemple \ce{Zn^2+}/\ce{Zn} ? La démarche est limpide :

\begin{methode}[Mesure du potentiel standard d'un couple]
\begin{enumerate}
\item Réaliser la \textbf{demi-pile standard} du couple étudié : par exemple une lame de zinc plongeant dans une solution de \ce{Zn^2+} à \SI{1}{mol.L^{-1}}, à \SI{25}{\celsius}.
\item L'\textbf{associer à l'ESH} par un pont salin et relier les deux électrodes à un voltmètre (circuit ouvert, courant quasi nul).
\item Mesurer la \textbf{f.é.m.} de la pile ainsi formée et identifier ses \textbf{pôles}.
\item En déduire le potentiel standard : $E^\circ(\text{couple}) = \pm\, e$, avec le signe $+$ si l'électrode étudiée est le pôle $+$ face à l'ESH, et le signe $-$ si elle est le pôle $-$.
\end{enumerate}
\end{methode}

En effet, la f.é.m. d'une pile est la différence entre les potentiels de ses deux électrodes :
\[ e = E_{+} - E_{-} \]
Comme l'un des deux pôles est l'ESH, dont le potentiel vaut exactement \SI{0,00}{V}, la f.é.m. mesurée donne \emph{directement}, en valeur absolue, le potentiel du couple étudié. Il ne reste qu'à déterminer son signe grâce à la polarité.

\textbf{Le signe du potentiel.} C'est le point capital :
\begin{itemize}
\item si l'électrode étudiée est le \textbf{pôle $-$} face à l'ESH (elle cède des électrons, son métal est plus réducteur que \ce{H2}), alors son potentiel standard est \textbf{négatif} ;
\item si l'électrode étudiée est le \textbf{pôle $+$} face à l'ESH (elle capte des électrons, son cation est plus oxydant que \ce{H+}), alors son potentiel standard est \textbf{positif}.
\end{itemize}

\begin{popfigure}
\begin{center}
\begin{tikzpicture}[scale=0.9, every node/.style={font=\small}]
% ===== Bécher ESH (gauche) =====
\fill[popblueL] (-5.6,0) rectangle (-3.2,2.6);
\draw[line width=1.1pt, popdark] (-5.6,-0.4) -- (-5.6,3.4);
\draw[line width=1.1pt, popdark] (-3.2,-0.4) -- (-3.2,3.4);
\draw[line width=1.1pt, popdark] (-5.6,-0.4) -- (-3.2,-0.4);
\draw[popblue, dashed] (-5.6,2.6) -- (-3.2,2.6);
\node[align=center] at (-4.4,-0.9) {ESH\\ \ce{H+} (\SI{1}{mol.L^{-1}}), \ce{H2} (\SI{1}{bar})};
% Platine
\fill[popdark] (-4.55,0.3) rectangle (-4.25,2.0);
\draw[line width=1pt, popdark] (-4.4,2.0) -- (-4.4,4.6);
% Bulles
\foreach \x/\y in {-4.9/0.6,-4.8/1.0,-4.95/1.4,-4.85/1.8,-4.9/2.2}{
  \draw[popblue!70!popdark] (\x,\y) circle (0.06);}
\draw[-{Stealth}, popgreen] (-6.6,1.2) -- (-5.7,1.2);
\node[popgreen, anchor=east] at (-6.7,1.2) {\ce{H2}};
% ===== Bécher Zn (droite) =====
\fill[popgreenL] (1.2,0) rectangle (3.6,2.6);
\draw[line width=1.1pt, popdark] (1.2,-0.4) -- (1.2,3.4);
\draw[line width=1.1pt, popdark] (3.6,-0.4) -- (3.6,3.4);
\draw[line width=1.1pt, popdark] (1.2,-0.4) -- (3.6,-0.4);
\draw[popgreen!60!popdark, dashed] (1.2,2.6) -- (3.6,2.6);
\node[align=center] at (2.4,-0.9) {lame de zinc dans\\ \ce{Zn^2+} (\SI{1}{mol.L^{-1}})};
% Lame de zinc
\fill[popmuted] (2.25,0.3) rectangle (2.55,2.4);
\draw[line width=1pt, popdark] (2.4,2.4) -- (2.4,4.6);
% ===== Pont salin =====
\draw[line width=3pt, poporange] (-3.9,2.3) .. controls (-2.5,4.0) and (0.5,4.0) .. (1.9,2.3);
\draw[line width=1pt, poporange!60!popdark] (-3.9,2.3) .. controls (-2.5,4.0) and (0.5,4.0) .. (1.9,2.3);
\node[poporange!70!popdark] at (-1.0,4.15) {pont salin};
% ===== Voltmètre =====
\draw[line width=1pt, popdark] (-4.4,4.6) -- (-1.0,4.6);
\draw[line width=1pt, popdark] (2.4,4.6) -- (-0.2,4.6);
\draw[line width=1.2pt, popdark] (-0.6,4.6) circle (0.45);
\node at (-0.6,4.6) {\textbf{V}};
\node[popdark, anchor=south] at (-0.6,5.15) {$e = \SI{0,76}{V}$};
% Polarités
\node[popink, font=\large\bfseries] at (-4.4,3.7) {$+$};
\node[popink, font=\large\bfseries] at (2.4,3.7) {$-$};
% Résultat
\node[align=center, poppurple] at (4.5,2.0) {$E^\circ(\ce{Zn^2+}/\ce{Zn})$\\ $= \SI{-0,76}{V}$};
\end{tikzpicture}
\end{center}
\legende{Mesure du potentiel standard du couple \ce{Zn^2+}/\ce{Zn} : la demi-pile au zinc, associée à l'ESH, forme une pile de f.é.m. \SI{0,76}{V} dont le zinc est le pôle $-$. On en déduit $E^\circ(\ce{Zn^2+}/\ce{Zn}) = \SI{-0,76}{V}$.}
\end{popfigure}

\begin{exoresolu}[Détermination du potentiel standard du couple \ce{Zn^2+}/\ce{Zn}]
On réalise, à \SI{25}{\celsius}, une pile en associant une demi-pile standard au zinc (lame de zinc plongeant dans une solution de \ce{Zn^2+} à \SI{1}{mol.L^{-1}}) et une électrode standard à hydrogène. Le voltmètre, branché entre les deux électrodes, indique une f.é.m. $e = \SI{0,76}{V}$. L'étude de la polarité montre que la lame de zinc constitue le pôle négatif de la pile.
\begin{enumerate}
\item Écrire les demi-équations aux électrodes et l'équation-bilan de fonctionnement.
\item En déduire le potentiel standard du couple \ce{Zn^2+}/\ce{Zn}.
\end{enumerate}
\tcblower
\textbf{1. Équations aux électrodes.}

Au pôle $-$ (zinc), il se produit une \emph{oxydation} : le métal zinc cède des électrons et passe en solution sous forme d'ions \ce{Zn^2+} :
\[ \ce{Zn -> Zn^2+ + 2e-} \]
Au pôle $+$ (ESH), il se produit une \emph{réduction} : les ions \ce{H+} captent les électrons à la surface du platine et dégagent du dihydrogène :
\[ \ce{2H+ + 2e- -> H2} \]
L'équation-bilan de fonctionnement s'obtient en additionnant les deux demi-équations (les deux électrons échangés se simplifient) :
\[ \ce{Zn + 2H+ -> Zn^2+ + H2} \]
On observe donc que la lame de zinc se consomme progressivement et qu'un dégagement gazeux de \ce{H2} se produit à la surface du platine.

\textbf{2. Potentiel standard du couple \ce{Zn^2+}/\ce{Zn}.}

La f.é.m. est la différence entre le potentiel du pôle $+$ et celui du pôle $-$ :
\[ e = E^\circ(\ce{H+}/\ce{H2}) - E^\circ(\ce{Zn^2+}/\ce{Zn}) \]
\[ \SI{0,76}{V} = \SI{0,00}{V} - E^\circ(\ce{Zn^2+}/\ce{Zn}) \]
d'où :
\[ \boxed{E^\circ(\ce{Zn^2+}/\ce{Zn}) = \SI{-0,76}{V}} \]
Le signe négatif traduit le fait que le zinc est le pôle $-$ face à l'ESH : le zinc est un réducteur \emph{plus fort} que le dihydrogène.
\end{exoresolu}

De la même manière, si l'on associe une demi-pile standard au cuivre à l'ESH, on mesure une f.é.m. de \SI{0,34}{V}, mais cette fois c'est le \emph{cuivre} qui constitue le pôle $+$ : les ions \ce{Cu^2+} se réduisent ($\ce{Cu^2+ + 2e- -> Cu}$) tandis que le dihydrogène s'oxyde ($\ce{H2 -> 2H+ + 2e-}$). Le potentiel standard du couple \ce{Cu^2+}/\ce{Cu} est donc \emph{positif} :
\[ E^\circ(\ce{Cu^2+}/\ce{Cu}) = \SI{+0,34}{V} \]

\begin{exemplebox}[Deux mesures de référence à connaître]
\begin{center}
\begin{tabularx}{\linewidth}{|l|X|X|X|}
\hline
\rowcolor{popblueL} \textbf{Couple étudié} & \textbf{Pôle face à l'ESH} & \textbf{f.é.m. mesurée} & \textbf{Potentiel standard} \\
\hline
\ce{Zn^2+}/\ce{Zn} & pôle $-$ & \SI{0,76}{V} & $E^\circ = \SI{-0,76}{V}$ \\
\hline
\ce{Cu^2+}/\ce{Cu} & pôle $+$ & \SI{0,34}{V} & $E^\circ = \SI{+0,34}{V}$ \\
\hline
\end{tabularx}
\end{center}
Remarque au passage : $\SI{+0,34}{V} - (\SI{-0,76}{V}) = \SI{1,10}{V}$, qui est exactement la f.é.m. de la pile Daniell étudiée dans les sections précédentes. Tout se recoupe !
\end{exemplebox}

\begin{aretenir}[L'essentiel de la section]
\begin{itemize}
\item On ne peut mesurer que des \emph{différences} de potentiel : il faut une référence universelle, l'\cle{électrode standard à hydrogène} (ESH).
\item L'ESH met en jeu le couple \ce{H+}/\ce{H2} ($\ce{2H+ + 2e- <=> H2}$) sur du platine platiné inerte et catalyseur, dans les conditions standard : $[\ce{H+}] = \SI{1}{mol.L^{-1}}$, $P(\ce{H2}) = \SI{1}{bar}$, $T = \SI{25}{\celsius}$.
\item Par convention : $E^\circ(\ce{H+}/\ce{H2}) = \SI{0,00}{V}$ à toute température. Ce zéro est \emph{arbitraire}, comme le niveau de la mer.
\item Pour mesurer le potentiel standard d'un couple, on l'associe à l'ESH : la f.é.m. donne la valeur absolue, la polarité donne le signe (pôle $-$ face à l'ESH $\Rightarrow$ $E^\circ < 0$ ; pôle $+$ $\Rightarrow$ $E^\circ > 0$).
\end{itemize}
\end{aretenir}

\textbf{Remarque pratique.} L'ESH est un instrument délicat : il faut du dihydrogène très pur, une pression rigoureusement contrôlée, un platine parfaitement propre. Dans les laboratoires, on lui préfère des électrodes de référence plus commodes et plus sûres, comme l'électrode au calomel ou l'électrode argent/chlorure d'argent (\ce{Ag}/\ce{AgCl}), dont les potentiels, mesurés une fois pour toutes par rapport à l'ESH, sont bien connus. Au programme de Première, cependant, tous les raisonnements et toutes les tables de potentiels sont exprimés \emph{par rapport à l'ESH}.

Nous disposons désormais d'une échelle de mesure universelle : chaque couple oxydant/réducteur possède un potentiel standard, positif ou négatif, qui le situe par rapport à l'hydrogène. Dans la section suivante, nous exploiterons cette \emph{classification quantitative} pour prévoir, par un simple calcul, la polarité et la f.é.m. de n'importe quelle pile.

\coursec{Potentiel d'oxydoréduction et classification quantitative}

Dans la section précédente, nous avons construit un instrument remarquable : l'\cle{électrode standard à hydrogène} (ESH), dont le potentiel est fixé par convention à \SI{0,00}{V}. Grâce à elle, nous disposons d'un « zéro » sur notre échelle, comme le niveau de la mer sert de zéro pour mesurer les altitudes. Mais à quoi sert un zéro si l'on ne mesure rien ? C'est exactement l'objet de cette section : mesurer, pour chaque couple oxydant/réducteur, une grandeur unique — le \cle{potentiel d'oxydoréduction} — qui chiffre sa « force », puis ranger tous les couples sur une même échelle. Tu vas voir qu'avec cette échelle en main, prévoir le sens d'une réaction d'oxydoréduction devient un jeu d'enfant.

\courssub{Le potentiel d'oxydoréduction : une intuition d'abord}

Imagine deux boxeurs sur un ring. Pour savoir qui va gagner, tu aimerais connaître la force de chacun. En oxydoréduction, c'est pareil : quand on met en présence deux couples, par exemple \ce{Cu^2+}/\ce{Cu} et \ce{Zn^2+}/\ce{Zn}, on aimerait savoir \emph{à l'avance} qui arrachera les électrons à qui.

Dans la pile Daniell étudiée précédemment, nous avons constaté expérimentalement que les électrons circulent du zinc vers le cuivre : le zinc « pousse » ses électrons, le cuivre les « aspire ». On dit que le zinc est un réducteur \emph{plus fort} que le cuivre, ou de manière équivalente que l'ion \ce{Cu^2+} est un oxydant \emph{plus fort} que l'ion \ce{Zn^2+}. Jusqu'ici, ce classement était \emph{qualitatif} : nous savions qui gagne, mais pas « de combien ».

L'idée géniale consiste à attribuer à chaque couple un nombre, mesuré en volts, qui quantifie cette force. Ce nombre s'obtient très simplement : on fabrique une pile en associant la demi-pile du couple étudié à l'ESH, et on mesure sa f.é.m. ainsi que sa polarité.

\begin{definition}[Potentiel d'oxydoréduction d'un couple]
Le \cle{potentiel d'oxydoréduction} d'un couple Ox/Red, noté $E(\text{Ox}/\text{Red})$, est une grandeur exprimée en \cle{volts} (V) qui mesure le pouvoir oxydant ou réducteur du couple. Il se mesure comme la f.é.m. de la pile formée par la demi-pile du couple et l'électrode de référence (l'ESH), affectée du signe de la polarité de l'électrode du couple.
\end{definition}

\courssub{Le potentiel standard $E^{\circ}$ : définition rigoureuse}

Le potentiel d'un couple dépend des conditions expérimentales : concentrations des espèces dissoutes, pression des gaz, température. Pour pouvoir comparer les couples entre eux, il faut donc se placer tous dans les \emph{mêmes} conditions, dites \cle{conditions standard} :

\begin{itemize}
\item concentrations des espèces dissoutes égales à \SI{1}{mol.L^{-1}} ;
\item pression des gaz égale à \SI{1}{bar} ;
\item température de \SI{25}{\celsius}.
\end{itemize}

\begin{definition}[Potentiel standard d'oxydoréduction]
Le \cle{potentiel standard d'oxydoréduction} d'un couple Ox/Red, noté $E^{\circ}(\text{Ox}/\text{Red})$, est la \cle{f.é.m.} de la pile formée par la demi-pile standard du couple Ox/Red et l'électrode standard à hydrogène (ESH), \emph{affectée du signe de la polarité de l'électrode du couple} :
\begin{itemize}
\item si l'électrode du couple est le pôle $\oplus$ de la pile, $E^{\circ} > 0$ ;
\item si l'électrode du couple est le pôle $\ominus$ de la pile, $E^{\circ} < 0$.
\end{itemize}
Par convention, $E^{\circ}(\ce{H+}/\ce{H2}) = \SI{0,00}{V}$.
\end{definition}

\begin{exemplebox}[Mesure du potentiel standard du couple \ce{Cu^2+}/\ce{Cu}]
On réalise la pile suivante : une lame de cuivre plongeant dans une solution de \ce{CuSO4} à \SI{1}{mol.L^{-1}} est reliée par un pont salin à une ESH. Le voltmètre indique une f.é.m. $e = \SI{0,34}{V}$, et l'électrode de cuivre est le pôle $\oplus$ (les électrons circulent de l'ESH vers le cuivre : \ce{Cu^2+} se réduit). On en déduit :
\[ E^{\circ}(\ce{Cu^2+}/\ce{Cu}) = +\SI{0,34}{V}. \]
Avec une lame de zinc à la place du cuivre, on mesure $e = \SI{0,76}{V}$, mais cette fois c'est le zinc qui est le pôle $\ominus$ (le zinc s'oxyde). Donc :
\[ E^{\circ}(\ce{Zn^2+}/\ce{Zn}) = -\SI{0,76}{V}. \]
Le signe n'est pas un détail : il porte toute l'information sur la polarité !
\end{exemplebox}

\courssub{Interprétation : lire la force des oxydants et des réducteurs}

Voici la règle de lecture de l'échelle des potentiels, qu'il faut absolument maîtriser :

\begin{propriete}[Signification du potentiel standard]
\begin{itemize}
\item Plus $E^{\circ}(\text{Ox}/\text{Red})$ est \textbf{élevé} (grand, positif), plus l'\cle{oxydant} du couple est \textbf{fort} (il capte facilement les électrons) — et plus le réducteur conjugué est faible.
\item Plus $E^{\circ}(\text{Ox}/\text{Red})$ est \textbf{faible} (négatif), plus le \cle{réducteur} du couple est \textbf{fort} (il cède facilement ses électrons) — et plus l'oxydant conjugué est faible.
\end{itemize}
\end{propriete}

Autrement dit : en haut de l'échelle vivent les oxydants « voraces » (comme \ce{Cl2}, qui arrache des électrons à presque tout le monde) ; en bas de l'échelle vivent les réducteurs « généreux » (comme le lithium, qui donne ses électrons à tout venant). L'ion \ce{H+}, avec son potentiel nul, marque la frontière : au-dessus de lui, les oxydants plus forts que \ce{H+} ; en dessous, les métaux réducteurs capables de dégager du dihydrogène au contact d'un acide.

\begin{attention}[Un potentiel négatif ne veut pas dire « pas de réaction » !]
Erreur très fréquente au Probatoire : croire qu'un couple de potentiel négatif, comme \ce{Zn^2+}/\ce{Zn} ($E^{\circ} = -\SI{0,76}{V}$), est « inerte » ou « ne réagit pas ». C'est faux ! Un potentiel négatif signifie simplement que le réducteur du couple est \emph{plus fort que} \ce{H2} : le zinc cède ses électrons plus facilement que le dihydrogène. C'est d'ailleurs pour cela que le zinc est attaqué par les acides avec dégagement de \ce{H2}, alors que le cuivre ($E^{\circ} > 0$) ne l'est pas. Retiens : le signe compare le couple à \ce{H+}/\ce{H2}, rien de plus.
\end{attention}

\courssub{La classification quantitative des couples}

En répétant la mesure pour tous les couples usuels, on obtient une table de potentiels standard. Voici les valeurs à connaître pour le programme de Première :

\begin{center}
\begin{tabularx}{\linewidth}{|c|>{\centering\arraybackslash}X|c|}
\hline
\rowcolor{popblueL} \textbf{Couple Ox/Red} & \textbf{Demi-équation (réduction)} & $E^{\circ}$ (V) \\
\hline
\ce{Li+}/\ce{Li} & \ce{Li+ + e- <=> Li} & $-3{,}04$ \\
\hline
\ce{Zn^2+}/\ce{Zn} & \ce{Zn^2+ + 2e- <=> Zn} & $-0{,}76$ \\
\hline
\ce{Fe^2+}/\ce{Fe} & \ce{Fe^2+ + 2e- <=> Fe} & $-0{,}44$ \\
\hline
\ce{H+}/\ce{H2} & \ce{2H+ + 2e- <=> H2} & $0{,}00$ \\
\hline
\ce{Cu^2+}/\ce{Cu} & \ce{Cu^2+ + 2e- <=> Cu} & $+0{,}34$ \\
\hline
\ce{Ag+}/\ce{Ag} & \ce{Ag+ + e- <=> Ag} & $+0{,}80$ \\
\hline
\ce{Cl2}/\ce{Cl-} & \ce{Cl2 + 2e- <=> 2Cl-} & $+1{,}36$ \\
\hline
\end{tabularx}
\end{center}

Visualisons cette hiérarchie sur une échelle verticale, qui est la représentation à savoir reproduire :

\begin{popfigure}
\begin{center}
\begin{tikzpicture}[scale=1.0]
% Axe vertical
\draw[->, thick, popdark] (0,-4.6) -- (0,4.9) node[above, font=\small\bfseries] {$E^{\circ}$ (V)};
% Graduations
\foreach \y/\lab in {-4/-4, -3/-3, -2/-2, -1/-1, 0/0, 1/1, 2/2, 3/3, 4/4}
  \draw[popmuted] (-0.12,\y) -- (0.12,\y) node[right=2pt, font=\footnotesize, popmuted] {\lab};
% Couples (1 unité = 1 V)
\fill[popink] (0,-3.04) circle (2.2pt) node[left=4pt, font=\small] {\ce{Li+}/\ce{Li} ($-3{,}04$~V)};
\fill[popink] (0,-0.76) circle (2.2pt) node[left=4pt, font=\small] {\ce{Zn^2+}/\ce{Zn} ($-0{,}76$~V)};
\fill[popink] (0,-0.44) circle (2.2pt) node[left=4pt, font=\small] {\ce{Fe^2+}/\ce{Fe} ($-0{,}44$~V)};
\fill[poporange] (0,0) circle (2.6pt) node[left=4pt, font=\small\bfseries] {\ce{H+}/\ce{H2} ($0{,}00$~V)};
\fill[popink] (0,0.34) circle (2.2pt) node[left=4pt, font=\small] {\ce{Cu^2+}/\ce{Cu} ($+0{,}34$~V)};
\fill[popink] (0,0.80) circle (2.2pt) node[left=4pt, font=\small] {\ce{Ag+}/\ce{Ag} ($+0{,}80$~V)};
\fill[popink] (0,1.36) circle (2.2pt) node[left=4pt, font=\small] {\ce{Cl2}/\ce{Cl-} ($+1{,}36$~V)};
% Flèches de lecture
\draw[->, very thick, popgreen] (1.6,-1.2) -- (1.6,2.6) node[midway, right, font=\small, align=left, text=popgreen] {oxydant de\\plus en plus fort};
\draw[->, very thick, poppurple] (-4.6,2.6) -- (-4.6,-1.2) node[midway, left, font=\small, align=right, text=poppurple] {réducteur de\\plus en plus fort};
\end{tikzpicture}
\end{center}
\legende{Échelle des potentiels standard : vers le haut, les oxydants sont de plus en plus forts ; vers le bas, les réducteurs sont de plus en plus forts. Le couple \ce{H+}/\ce{H2} sert de référence ($E^{\circ} = \SI{0,00}{V}$).}
\end{popfigure}

\begin{savaistu}[Pourquoi le lithium est-il le roi des batteries ?]
Avec $E^{\circ}(\ce{Li+}/\ce{Li}) = -\SI{3,04}{V}$, le lithium est le réducteur le plus fort de toute la classification : aucun métal ne cède ses électrons plus facilement. Associé à un couple de potentiel élevé, il permet de fabriquer des piles de très forte f.é.m. et de grande densité d'énergie. C'est pourquoi les batteries au lithium équipent téléphones portables et véhicules électriques — y compris ceux qu'on voit de plus en plus à Douala et Yaoundé. En contrepartie, le lithium réagit violemment avec l'eau : ces batteries doivent être parfaitement étanches.
\end{savaistu}

\courssub{Lien avec la règle du gamma : du qualitatif au quantitatif}

En début de module, tu as appris la \cle{règle du gamma} : en plaçant deux couples l'un au-dessus de l'autre dans une classification qualitative, la réaction spontanée suit les branches du gamma ($\gamma$) : l'oxydant du couple placé \emph{au-dessus} réagit avec le réducteur du couple placé \emph{en dessous}.

La classification quantitative par les potentiels standard \emph{précise et généralise} cette règle : le classement vertical n'est plus approximatif, il est chiffré. « Au-dessus » signifie désormais « de potentiel $E^{\circ}$ plus élevé ». La règle du gamma devient ainsi un outil de \emph{prévision quantitative} : non seulement on sait qui réagit avec qui, mais on peut calculer la f.é.m. de la pile correspondante, comme nous le ferons dans la section suivante.

\begin{popfigure}
\begin{center}
\begin{tikzpicture}[scale=1.0]
% Échelle
\draw[->, thick, popdark] (0,-1.6) -- (0,1.8) node[above, font=\small\bfseries] {$E^{\circ}$};
% Couple Cu (en haut)
\node[font=\small\bfseries, text=popblue] (cuox) at (-1.4,0.8) {\ce{Cu^2+}};
\node[font=\small\bfseries, text=popblue] (cured) at (1.4,0.8) {\ce{Cu}};
\node[font=\footnotesize, text=popmuted] at (2.6,0.8) {$E^{\circ} = +\SI{0,34}{V}$};
\draw[popblue] (-0.7,0.8) -- (0.7,0.8);
% Couple Zn (en bas)
\node[font=\small\bfseries, text=poppurple] (znox) at (-1.4,-0.8) {\ce{Zn^2+}};
\node[font=\small\bfseries, text=poppurple] (znred) at (1.4,-0.8) {\ce{Zn}};
\node[font=\footnotesize, text=popmuted] at (2.6,-0.8) {$E^{\circ} = -\SI{0,76}{V}$};
\draw[poppurple] (-0.7,-0.8) -- (0.7,-0.8);
% Branches du gamma
\draw[->, very thick, poporange] (cuox.south east) .. controls (0.2,0.1) .. (znred.north west);
% Labels
\node[font=\footnotesize, text=poporange, align=center] at (-2.9,-0.1) {sens de la\\réaction\\spontanée};
\fill[popblue] (0,0.8) circle (1.8pt);
\fill[poppurple] (0,-0.8) circle (1.8pt);
\end{tikzpicture}
\end{center}
\legende{Règle du gamma appliquée aux couples \ce{Cu^2+}/\ce{Cu} et \ce{Zn^2+}/\ce{Zn} : l'oxydant du couple de plus haut potentiel (\ce{Cu^2+}) réagit avec le réducteur du couple de plus bas potentiel (\ce{Zn}). C'est exactement la réaction de la pile Daniell.}
\end{popfigure}

\courssub{Prévision d'une réaction spontanée}

\begin{propriete}[Critère de prévision]
Deux couples Ox$_1$/Red$_1$ et Ox$_2$/Red$_2$ étant mis en présence, la réaction \cle{spontanée} est celle qui oppose :
\begin{itemize}
\item l'\textbf{oxydant du couple de plus haut potentiel} ($E^{\circ}$ le plus élevé) ;
\item au \textbf{réducteur du couple de plus bas potentiel} ($E^{\circ}$ le plus faible).
\end{itemize}
Elle s'écrit : Ox$_1$ (haut) $+$ Red$_2$ (bas) $\ce{->}$ Red$_1$ $+$ Ox$_2$.
\end{propriete}

\begin{methode}[Prévoir le sens d'une réaction d'oxydoréduction]
\begin{enumerate}
\item \textbf{Identifier} les deux couples Ox/Red présents dans le milieu et écrire leurs demi-équations de réduction.
\item \textbf{Placer} les deux couples sur l'échelle des potentiels standard (ou comparer leurs valeurs de $E^{\circ}$).
\item \textbf{Appliquer} la règle : l'oxydant du couple le plus haut (plus grand $E^{\circ}$) oxyde le réducteur du couple le plus bas (plus petit $E^{\circ}$).
\item \textbf{Écrire} l'équation-bilan en combinant les demi-équations de manière à éliminer les électrons.
\item \textbf{Vérifier} que les électrons n'apparaissent plus dans le bilan et que les charges sont équilibrées.
\end{enumerate}
\end{methode}

\begin{exemplebox}[L'ion argent oxyde le cuivre]
Comparons : $E^{\circ}(\ce{Ag+}/\ce{Ag}) = +\SI{0,80}{V} > E^{\circ}(\ce{Cu^2+}/\ce{Cu}) = +\SI{0,34}{V}$. L'oxydant du couple le plus haut est \ce{Ag+} ; le réducteur du couple le plus bas est \ce{Cu}. La réaction spontanée est donc :
\begin{align*}
\text{réduction : } & \ce{Ag+ + e- -> Ag} \quad (\times 2)\\
\text{oxydation : } & \ce{Cu -> Cu^2+ + 2e-}\\[2pt]
\text{bilan : } & \ce{2Ag+ + Cu -> 2Ag + Cu^2+}
\end{align*}
C'est la réaction de « l'arbre de Diane » : une lame de cuivre plongée dans une solution de nitrate d'argent se recouvre de cristaux d'argent métallique, tandis que la solution bleuit (formation d'ions \ce{Cu^2+}). Réaction totale, rapide, sans catalyseur ni chauffage.
\end{exemplebox}

\begin{exoresolu}[Le fer est-il attaqué par les ions cuivre ?]
On plonge un clou en fer dans une solution bleue de sulfate de cuivre (II). Prévoir s'il se produit une réaction et, le cas échéant, écrire son équation-bilan. On donne : $E^{\circ}(\ce{Fe^2+}/\ce{Fe}) = -\SI{0,44}{V}$ et $E^{\circ}(\ce{Cu^2+}/\ce{Cu}) = +\SI{0,34}{V}$.
\tcblower
\textbf{Étape 1 — Couples en présence :} \ce{Cu^2+}/\ce{Cu} et \ce{Fe^2+}/\ce{Fe}.

\textbf{Étape 2 — Comparaison des potentiels :}
\[ E^{\circ}(\ce{Cu^2+}/\ce{Cu}) = +\SI{0,34}{V} > E^{\circ}(\ce{Fe^2+}/\ce{Fe}) = -\SI{0,44}{V}. \]

\textbf{Étape 3 — Règle du gamma :} l'oxydant du couple le plus haut, \ce{Cu^2+}, réagit avec le réducteur du couple le plus bas, \ce{Fe}. Une réaction spontanée a donc lieu.

\textbf{Étape 4 — Équation-bilan :}
\begin{align*}
\text{réduction : } & \ce{Cu^2+ + 2e- -> Cu}\\
\text{oxydation : } & \ce{Fe -> Fe^2+ + 2e-}\\
\text{bilan : } & \ce{Cu^2+ + Fe -> Cu + Fe^2+}
\end{align*}

\textbf{Observations attendues :} le clou se recouvre d'un dépôt rougeâtre de cuivre métallique et la solution se décolore progressivement (les ions bleus \ce{Cu^2+} sont remplacés par les ions \ce{Fe^2+}, vert très pâle). C'est d'ailleurs le principe utilisé pour protéger ou cuivrer des objets en fer.
\end{exoresolu}

\courssub{Cas limite : des potentiels très proches}

Que se passe-t-il lorsque les deux couples ont des potentiels standard \emph{très proches} ? La « différence de force » entre l'oxydant et le réducteur est alors faible : la réaction a bien lieu dans le sens prévu par la règle du gamma, mais elle est \cle{peu avancée} — elle conduit à un état d'\cle{équilibre} où les quatre espèces coexistent en quantités notables.

\begin{propriete}[Écart des potentiels et avancement de la réaction]
\begin{itemize}
\item Plus l'écart $\Delta E^{\circ} = E^{\circ}_{\text{haut}} - E^{\circ}_{\text{bas}}$ est \textbf{grand}, plus la réaction est \textbf{totale} (forte constante d'équilibre).
\item Si $\Delta E^{\circ}$ est \textbf{très faible} (voisin de zéro), la réaction est \textbf{limitée} : on obtient un équilibre chimique, et la réaction inverse n'est plus négligeable.
\end{itemize}
\end{propriete}

Ce cas limite n'est pas une exception théorique : c'est précisément ce qui rend possibles les \emph{accumulateurs} (batteries rechargeables). Dans un accumulateur au plomb de voiture, les couples mis en jeu ont des potentiels tels que la réaction peut être inversée en forçant le courant dans l'autre sens lors de la charge. La notion d'équilibre en oxydoréduction, que tu approfondiras en Terminale, est déjà en germe ici.

\begin{aretenir}[L'essentiel de la section]
\begin{itemize}
\item Le potentiel standard $E^{\circ}(\text{Ox}/\text{Red})$ est la f.é.m. de la pile « couple étudié // ESH », signée selon la polarité de l'électrode du couple ; $E^{\circ}(\ce{H+}/\ce{H2}) = \SI{0,00}{V}$ par convention.
\item Plus $E^{\circ}$ est élevé, plus l'oxydant est fort ; plus $E^{\circ}$ est faible, plus le réducteur est fort.
\item Un potentiel négatif ne signifie pas « inerte » : il signifie « réducteur plus fort que \ce{H2} ».
\item Réaction spontanée : oxydant du couple de plus haut $E^{\circ}$ $+$ réducteur du couple de plus bas $E^{\circ}$ (règle du gamma quantifiée).
\item Potentiels très proches $\Rightarrow$ réaction limitée conduisant à un équilibre.
\end{itemize}
\end{aretenir}

Nous avons maintenant tous les outils en main : une échelle chiffrée et une règle de prévision. Dans la section suivante, nous les mettrons au service de la pile elle-même : prévoir sa polarité, calculer sa f.é.m. et écrire ses équations de fonctionnement sans même avoir besoin de la construire.

\coursec{Prévision de la polarité et calcul de la f.é.m.}

Dans la section précédente, nous avons construit l'outil fondamental : l'échelle des potentiels standards d'oxydoréduction. Cette échelle classe les couples \ce{Ox/Red} par ordre croissant de leur aptitude à capter des électrons (pouvoir oxydant). Nous savons désormais que \ce{F2} est le plus fort oxydant (\cle{$E^\circ = \SI{+2,87}{V}$}) et que \ce{Li+} est le plus faible (\cle{$E^\circ = \SI{-3,05}{V}$}). Mais une échelle ne sert à rien si l'on ne sait pas l'utiliser pour \emph{prédire} ce qui se passe quand on met deux couples en contact via un circuit externe. C'est tout l'objet de cette section : transformer deux demi-piles en une pile fonctionnelle, en déterminer les pôles, calculer sa tension (f.é.m.) et écrire son équation-bilan.

\courssub{Problématique : associer deux demi-piles standard}

Imaginons que nous disposions de deux béchers :
\begin{itemize}
    \item Bécher 1 : une lame de zinc \ce{Zn(s)} plongée dans une solution de \ce{ZnSO4} \SI{1,0}{mol.L^{-1}} (couple \ce{Zn^2+/Zn}, \cle{$E^\circ_1 = \SI{-0,76}{V}$}).
    \item Bécher 2 : une lame de cuivre \ce{Cu(s)} plongée dans une solution de \ce{CuSO4} \SI{1,0}{mol.L^{-1}} (couple \ce{Cu^2+/Cu}, \cle{$E^\circ_2 = \SI{+0,34}{V}$}).
\end{itemize}
Nous les relions par un pont salin (circuit interne) et un fil conducteur muni d'un voltmètre (circuit externe). \textbf{Que se passe-t-il ?} Dans quel sens les électrons circulent-ils ? Quelle tension le voltmètre affiche-t-il ? Quelle est l'équation de la réaction chimique spontanée ?

\courssub{Règle de polarité : identifier l'anode et la cathode}

La réponse repose sur une propriété thermodynamique fondamentale : \emph{la réaction spontanée est celle qui fait baisser l'énergie libre du système}. En électrochimie, cela se traduit par une règle simple, conséquence directe de la définition du potentiel standard :

\begin{aretenir}[Règle de polarité d'une pile]
Lorsque l'on associe deux demi-piles standard de couples \ce{Ox1/Red1} (\cle{$E^\circ_1$}) et \ce{Ox2/Red2} (\cle{$E^\circ_2$}) :
\begin{itemize}
    \item Le couple de \textbf{plus haut potentiel} (le plus fort oxydant) impose sa réaction de \textbf{réduction}. Son électrode est la \cle{cathode} (pôle \textbf{+}, borne rouge du voltmètre).
    \item Le couple de \textbf{plus bas potentiel} (le plus fort réducteur) impose sa réaction d'\textbf{oxydation}. Son électrode est l'\cle{anode} (pôle \textbf{−}, borne noire/COM du voltmètre).
\end{itemize}
Les électrons circulent dans le fil externe \textbf{de l'anode (–) vers la cathode (+)}.
\end{aretenir}

\emph{Intuition :} Le couple qui a le plus « envie » de prendre des électrons (potentiel élevé) les arrache à l'autre. L'électrode où les électrons sont « arrachés » (oxydation) se charge négativement ; celle où ils arrivent (réduction) se charge positivement.

\courssub{Calcul de la force électromotrice (f.é.m.) standard}

La force électromotrice standard \cle{$E^\circ_{\text{pile}}$} (ou \cle{$\Delta E^\circ$}) est la différence de potentiel entre les deux électrodes \emph{mesurée circuit ouvert} (courant nul). Elle s'exprime en volts (V).

\begin{propriete}[Formule de la f.é.m. standard]
\[
\boxed{E^\circ_{\text{pile}} = E^\circ_{\text{cathode}} - E^\circ_{\text{anode}} = E^\circ_{(+)} - E^\circ_{(-)}}
\]
\begin{itemize}
    \item \cle{$E^\circ_{\text{cathode}}$} : potentiel standard du couple réduit à la cathode (réduction).
    \item \cle{$E^\circ_{\text{anode}}$} : potentiel standard du couple oxydé à l'anode (oxydation).
    \item Le résultat est \textbf{toujours positif} pour une pile fonctionnant spontanément. Si vous trouvez une valeur négative, vous avez inversé l'anode et la cathode.
\end{itemize}
\end{propriete}

\courssub{Application à la pile Daniell}

Reprenons notre exemple : \ce{Zn^2+/Zn} (\cle{$\SI{-0,76}{V}$}) et \ce{Cu^2+/Cu} (\cle{$\SI{+0,34}{V}$}).
\begin{enumerate}
    \item \textbf{Classement :} \cle{$\SI{-0,76}{V} < \SI{+0,34}{V}$}.
    \item \textbf{Pôles :} \ce{Cu^2+/Cu} a le potentiel le plus haut $\rightarrow$ \textbf{Cathode (+)}. \ce{Zn^2+/Zn} a le potentiel le plus bas $\rightarrow$ \textbf{Anode (–)}.
    \item \textbf{F.é.m. :}
    \[
    E^\circ_{\text{pile}} = E^\circ(\ce{Cu^2+/Cu}) - E^\circ(\ce{Zn^2+/Zn}) = 0,34 - (-0,76) = \mathbf{\SI{1,10}{V}}
    \]
\end{enumerate}
C'est la valeur de référence de la pile Daniell, celle que tu mesureras (à quelques hundredièmes près) avec un voltmètre de bonne qualité sur une pile fraîche.

\courssub{Méthode complète : de la classification à l'équation-bilan}

Voici la démarche rigoureuse, étape par étape, pour résoudre \emph{tout} problème de prévision sur une pile standard.

\begin{methode}[Détermination de la polarité, de la f.é.m. et de l'équation-bilan d'une pile]
\textbf{Étape 1 : Classer les potentiels.} \\
Écrire les deux couples avec leur \cle{$E^\circ$} (tableau de référence). Identifier le plus grand (\cle{$E^\circ_{\max}$}) et le plus petit (\cle{$E^\circ_{\min}$}).

\textbf{Étape 2 : Identifier les pôles et la polarité.} \\
\begin{itemize}
    \item Cathode (pôle +) = couple de \cle{$E^\circ_{\max}$} $\rightarrow$ Réduction : $\ce{Ox_{\max} + n e- -> Red_{\max}}$.
    \item Anode (pôle –) = couple de \cle{$E^\circ_{\min}$} $\rightarrow$ Oxydation : $\ce{Red_{\min} -> Ox_{\min} + n e-}$.
\end{itemize}

\textbf{Étape 3 : Calculer la f.é.m. standard.} \\
\cle{$E^\circ_{\text{pile}} = E^\circ_{\max} - E^\circ_{\min}$} (résultat positif en V).

\textbf{Étape 4 : Écrire les équations aux électrodes et l'équation-bilan.} \\
\begin{itemize}
    \item Écrire la demi-équation de réduction (cathode) et d'oxydation (anode).
    \item Équilibrer les charges par les électrons \ce{e-}.
    \item Multiplier les demi-équations par les coefficients nécessaires pour avoir \textbf{le même nombre d'électrons} des deux côtés.
    \item Additionner les deux demi-équations et simplifier les électrons (et les espèces communes le cas échéant) pour obtenir l'équation-bilan.
\end{itemize}

\textbf{Étape 5 (Bonus) : Écrire la notation conventionnelle.} \\
\cle{Anode | Electrolyte anode || Electrolyte cathode | Cathode} (phases séparées par \cle{|}, pont salin par \cle{||}).
\end{methode}

\begin{popfigure}
\begin{tikzpicture}[
    node distance=1.2cm and 1.5cm,
    box/.style={draw=popdark, thick, rounded corners, fill=popcream, minimum width=3.5cm, minimum height=1.2cm, align=center, font=\small},
    arrow/.style={-Stealth, thick, popblue},
    decision/.style={diamond, draw=poporange, thick, fill=poporangeL, aspect=2, inner sep=4pt, font=\small, align=center}
]
\node[box] (start) {Données : 2 couples \\ Ox1/Red1 (E°1) et Ox2/Red2 (E°2)};
\node[box, below=of start] (step1) {\textbf{Étape 1} : Classer E°1 et E°2 \\ (croissant)};
\node[box, below=of step1] (step2) {\textbf{Étape 2} : Identifier pôles \\ Cathode = E° max (+) \\ Anode = E° min (–)};
\node[box, below=of step2] (step3) {\textbf{Étape 3} : Calculer f.é.m. \\ E°pile = E°(+) – E°(–) > 0};
\node[box, below=of step3] (step4) {\textbf{Étape 4} : Écrire équations \\ Oxydation (Anode) + Réduction (Cathode) $\rightarrow$ Équation-bilan};
\node[box, below=of step4, fill=popgreenL, draw=popgreen] (step5) {\textbf{Étape 5} : Notation conventionnelle \\ Anode | ... || ... | Cathode};

\draw[arrow] (start) -- (step1);
\draw[arrow] (step1) -- (step2);
\draw[arrow] (step2) -- (step3);
\draw[arrow] (step3) -- (step4);
\draw[arrow] (step4) -- (step5);
\end{tikzpicture}
\legende{Organigramme de la méthode en 5 étapes pour prévoir le fonctionnement d'une pile standard.}
\end{popfigure}

\courssub{Écriture des équations aux électrodes et équation-bilan}

Appliquons l'étape 4 à la pile Daniell.
\begin{itemize}
    \item \textbf{Anode (–) :} Oxydation du zinc. \ce{Zn(s) -> Zn^2+(aq) + 2 e-}
    \item \textbf{Cathode (+) :} Réduction du cuivre. \ce{Cu^2+(aq) + 2 e- -> Cu(s)}
    \item \textbf{Bilan :} Les 2 électrons s'annulent.
    \[
    \ce{Zn(s) + Cu^2+(aq) -> Zn^2+(aq) + Cu(s)}
    \]
    \emph{Observation :} La lame de zinc s'use (masse diminue), la solution se décolore (disparition des ions \ce{Cu^2+} bleus), et du cuivre rouge se dépose sur la lame de cuivre (masse augmente).
\end{itemize}

\courssub{Notation conventionnelle de la pile}

Elle résume la pile en une ligne, \textbf{toujours de l'anode (gauche) vers la cathode (droite)}.

\begin{propriete}[Notation conventionnelle (diagramme de pile)]
\begin{itemize}
    \item On sépare les phases différentes par une barre verticale \cle{|}.
    \item Le pont salin (interface entre les deux solutions) est noté \cle{||} (double barre).
    \item Les états physiques sont indiqués en indice : \cle{(s)}, \cle{(aq)}, \cle{(l)}, \cle{(g)}.
    \item Les concentrations (ou pressions) standards \SI{1}{mol.L^{-1}} (\SI{1}{bar}) sont souvent omises ou notées \cle{(1 M)}.
\end{itemize}
Pour la pile Daniell :
\[
\ce{Zn(s) | Zn^2+(aq, 1 M) || Cu^2+(aq, 1 M) | Cu(s)}
\]
\end{propriete}

\courssub{Exemples guidés variés}

Appliquons la méthode à trois piles classiques du programme.

\begin{exoresolu}[Pile \ce{Fe^2+/Fe} / \ce{Ag+/Ag} : polarité, f.é.m. et équation-bilan]
\textbf{Données :} \cle{$E^\circ(\ce{Fe^2+/Fe}) = \SI{-0,44}{V}$} ; \cle{$E^\circ(\ce{Ag+/Ag}) = \SI{+0,80}{V}$}. Solutions \SI{1}{M}, T = \SI{25}{\celsius}.

\textbf{1. Classement :} \cle{$\SI{-0,44}{V} < \SI{+0,80}{V}$}.
\textbf{2. Pôles :}
\begin{itemize}
    \item Cathode (+) : couple \ce{Ag+/Ag} (E° max). Réduction : \ce{Ag+ + e- -> Ag(s)}.
    \item Anode (–) : couple \ce{Fe^2+/Fe} (E° min). Oxydation : \ce{Fe(s) -> Fe^2+ + 2 e-}.
\end{itemize}
\textbf{3. F.é.m. :}
\[
E^\circ_{\text{pile}} = 0,80 - (-0,44) = \mathbf{\SI{1,24}{V}}
\]
\textbf{4. Équations :}
\begin{itemize}
    \item Anode : \ce{Fe(s) -> Fe^2+(aq) + 2 e-} \quad (x1)
    \item Cathode : \ce{Ag+(aq) + e- -> Ag(s)} \quad (x2 pour 2 e-)
\end{itemize}
\textbf{Bilan :}
\[
\ce{Fe(s) + 2 Ag+(aq) -> Fe^2+(aq) + 2 Ag(s)}
\]
\textbf{5. Notation :}
\[
\ce{Fe(s) | Fe^2+(aq, 1 M) || Ag+(aq, 1 M) | Ag(s)}
\]
\tcblower
\textbf{Analyse :} Le fer métallique se dissout (corrosion anodique), les ions argent se déposent en métal argenté brillant sur l'électrode d'argent. La f.é.m. de \SI{1,24}{V} est supérieure à celle de Daniell (\SI{1,10}{V}) car l'écart de potentiel est plus grand.
\end{exoresolu}

\begin{exemplebox}[Pile \ce{Zn^2+/Zn} / \ce{Ag+/Ag} (Pile Zn/Ag)]
\textbf{Données :} \cle{$E^\circ(\ce{Zn^2+/Zn}) = \SI{-0,76}{V}$} ; \cle{$E^\circ(\ce{Ag+/Ag}) = \SI{+0,80}{V}$}.
\begin{enumerate}
    \item \textbf{Pôles :} Cathode (+) = Ag ; Anode (–) = Zn.
    \item \textbf{F.é.m. :} \cle{$E^\circ = 0,80 - (-0,76) = \SI{1,56}{V}$}.
    \item \textbf{Équations :}
    \begin{align*}
    \text{Anode : } & \ce{Zn(s) -> Zn^2+ + 2 e-} \\
    \text{Cathode : } & \ce{2 Ag+ + 2 e- -> 2 Ag(s)} \\
    \text{Bilan : } & \ce{Zn(s) + 2 Ag+(aq) -> Zn^2+(aq) + 2 Ag(s)}
    \end{align*}
    \item \textbf{Notation :} \ce{Zn(s) | Zn^2+(aq) || Ag+(aq) | Ag(s)}.
\end{enumerate}
Cette pile délivre une tension élevée (\SI{1,56}{V}), proche des piles alcalines du commerce (\SI{1,5}{V}).
\end{exemplebox}

\begin{exemplebox}[Pile \ce{Fe^2+/Fe} / \ce{Cu^2+/Cu} (Pile Fe/Cu)]
\textbf{Données :} \cle{$E^\circ(\ce{Fe^2+/Fe}) = \SI{-0,44}{V}$} ; \cle{$E^\circ(\ce{Cu^2+/Cu}) = \SI{+0,34}{V}$}.
\begin{enumerate}
    \item \textbf{Pôles :} Cathode (+) = Cu ; Anode (–) = Fe.
    \item \textbf{F.é.m. :} \cle{$E^\circ = 0,34 - (-0,44) = \SI{0,78}{V}$}.
    \item \textbf{Bilan :} \ce{Fe(s) + Cu^2+(aq) -> Fe^2+(aq) + Cu(s)}.
    \item \textbf{Notation :} \ce{Fe(s) | Fe^2+(aq) || Cu^2+(aq) | Cu(s)}.
\end{enumerate}
\end{exemplebox}

\begin{exemplebox}[Pile \ce{Zn^2+/Zn} / \ce{Fe^2+/Fe} (Pile Zn/Fe)]
\textbf{Données :} \cle{$E^\circ(\ce{Zn^2+/Zn}) = \SI{-0,76}{V}$} ; \cle{$E^\circ(\ce{Fe^2+/Fe}) = \SI{-0,44}{V}$}.
\begin{enumerate}
    \item \textbf{Pôles :} Cathode (+) = Fe (potentiel \emph{moins négatif}) ; Anode (–) = Zn.
    \item \textbf{F.é.m. :} \cle{$E^\circ = (-0,44) - (-0,76) = \SI{0,32}{V}$}.
    \item \textbf{Bilan :} \ce{Zn(s) + Fe^2+(aq) -> Zn^2+(aq) + Fe(s)}.
    \item \textbf{Notation :} \ce{Zn(s) | Zn^2+(aq) || Fe^2+(aq) | Fe(s)}.
\end{enumerate}
\emph{Remarque :} Même si les deux potentiels sont négatifs, la pile fonctionne car il y a une \emph{différence} de potentiel. Le zinc est un réducteur plus fort que le fer.
\end{exemplebox}

\begin{popfigure}
\begin{tikzpicture}[
    scale=0.9,
    every node/.style={font=\small},
    beaker/.style={draw=popdark, thick, fill=popblueL!30, minimum width=2.5cm, minimum height=3cm, rounded corners=2pt},
    electrode/.style={draw=popdark, very thick, fill=popmuted, minimum width=0.4cm, minimum height=2.5cm},
    solution/.style={fill=popblue!20, draw=none},
    bridge/.style={draw=popdark, thick, fill=poporangeL, minimum width=0.8cm, minimum height=1.5cm, rounded corners=2pt},
    wire/.style={popdark, very thick},
    voltmeter/.style={draw=popdark, thick, fill=popcream, minimum width=1.5cm, minimum height=1cm, rounded corners=2pt},
    arrow/.style={-Stealth, poporange, thick},
    earrow/.style={-Stealth, poppink, very thick, dashed}
]
% Becher Gauche (Anode Zn)
\node[beaker] (beakL) at (0,0) {};
\node[solution, minimum width=2.1cm, minimum height=2.2cm] at (0,-0.2) {};
\node[electrode] (Zn) at (0,1.2) {};
\node[above=2pt of Zn, popdark] {\ce{Zn(s)}};
\node[below=2pt of Zn, popmuted] {Anode (–)};
\node at (0,-1.8) {\ce{Zn^2+(aq) + SO4^2-}};

% Becher Droit (Cathode Ag)
\node[beaker] (beakR) at (7,0) {};
\node[solution, minimum width=2.1cm, minimum height=2.2cm] at (7,-0.2) {};
\node[electrode] (Ag) at (7,1.2) {};
\node[above=2pt of Ag, popdark] {\ce{Ag(s)}};
\node[below=2pt of Ag, popmuted] {Cathode (+)};
\node at (7,-1.8) {\ce{Ag+(aq) + NO3-}};

% Pont salin
\node[bridge] (bridge) at (3.5, -1.5) {};
\node[above=2pt of bridge] {\ce{KNO3 (sat)}};
\draw[wire] (beakL.south east) -- (bridge.north west);
\draw[wire] (bridge.north east) -- (beakR.south west);

% Circuit externe
\draw[wire] (Zn.north) -- ++(0,1) -| (4.5, 2.2) node[voltmeter, pos=0.5] (V) {} -| (Ag.north);
\node[above=2pt of V] {\textbf{Voltmètre}};
\node[right=2pt of V, poppink, font=\bfseries\large] {\SI{1,56}{V}};

% Flèches électrons (fil)
\draw[earrow] (Zn.north) -- ++(0,0.5) node[midway, left] {$e^-$};
\draw[earrow] (4.5, 2.2) -- (Ag.north) node[midway, above] {$e^-$};

% Flèches ions (solutions)
\draw[arrow] (0,-0.5) -- ++(-1,0) node[left] {\ce{Zn^2+} \textbf{produit}};
\draw[arrow] (7,-0.5) -- ++(1,0) node[right] {\ce{Ag+} \textbf{consommé}};
\draw[arrow, popgreen] (bridge.north) -- ++(0,0.8) -| (0,-0.5) node[midway, left] {\ce{NO3-}};
\draw[arrow, popgreen] (bridge.north) -- ++(0,0.8) -| (7,-0.5) node[midway, right] {\ce{K+}};

% Légendes réactions
\node[align=left, popdark, font=\scriptsize] at (-3.5, 1.5) {Anode (Oxydation) : \\ \ce{Zn -> Zn^2+ + 2e-}};
\node[align=right, popdark, font=\scriptsize] at (10.5, 1.5) {Cathode (Réduction) : \\ \ce{2Ag+ + 2e- -> 2Ag}};
\end{tikzpicture}
\legende{Schéma de la pile Zn/Ag : polarité, circulation des électrons (flèches roses) et des ions, f.é.m. standard \SI{1,56}{V}.}
\end{popfigure}

\courssub{Pièges à éviter et compléments pour le Probatoire}

\begin{attention}[Erreur fréquente : multiplier le potentiel par le coefficient stœchiométrique]
\textbf{Jamais !} Le potentiel standard \cle{$E^\circ$} est une \cle{grandeur intensive} (comme la température, la pression). Il ne dépend \textbf{pas} de la quantité de matière.
\begin{itemize}
    \item Pour la pile Fe/Ag, on a multiplié la demi-équation de l'argent par 2 pour équilibrer les électrons.
    \item \textbf{Erreur :} \cle{$E^\circ = 0,80 \times 2 - (-0,44) = \SI{2,04}{V}$} $\leftarrow$ \textbf{FAUX}.
    \item \textbf{Correct :} \cle{$E^\circ = 0,80 - (-0,44) = \SI{1,24}{V}$}.
\end{itemize}
L'équilibrage des électrons ne sert qu'à écrire l'équation-bilan correcte, \textbf{pas} à calculer la tension.
\end{attention}

\begin{attention}[F.é.m. mesurée vs f.é.m. théorique : la résistance interne]
La formule \cle{$E^\circ_{\text{pile}} = E^\circ_{(+)} - E^\circ_{(-)}$} donne la tension \textbf{circuit ouvert} (courant \cle{$I = 0$}). Si tu fermes le circuit (ou si le voltmètre a une résistance interne trop faible), un courant circule. La pile a une \cle{résistance interne $r$} (électrolytes, pont salin, interfaces). La tension aux bornes \cle{$U$} chute alors :
\[
U = E^\circ_{\text{pile}} - r I \quad (\text{loi d'Ohm pour le générateur})
\]
\begin{itemize}
    \item Un bon voltmètre numérique a une résistance d'entrée \cle{$R_{\text{entrée}} \approx \SI{10}{\mega\ohm}$} : \cle{$I \approx 0$}, \cle{$U \approx E^\circ$}.
    \item Un voltmètre analogique (galvanomètre) ou une pile « usée » (résistance interne augmentée) donnera une tension \textbf{inférieure} à la f.é.m. théorique.
    \item Au Probatoire, on te demandera souvent : « La tension mesurée est-elle égale à la f.é.m. calculée ? Justifie. » $\rightarrow$ Réponse : « Non (ou oui si voltmètre idéal), car la pile a une résistance interne qui provoque une chute de tension \cle{$rI$} dès que le circuit est fermé. »
\end{itemize}
\end{attention}

\courssub{Lien avec la mesure expérimentale}

En TP, tu réalises la pile (ex: Daniell ou Zn/Ag) et tu mesures sa tension.
\begin{enumerate}
    \item \textbf{Branchement du voltmètre :} Borne \textbf{COM (noire)} sur l'électrode prévue comme \textbf{anode (–)} ; Borne \textbf{V/$\Omega$ (rouge)} sur l'électrode prévue comme \textbf{cathode (+)}.
    \item \textbf{Lecture :} Si l'affichage est \textbf{positif} (ex: \SI{1,10}{V}), ta prévision de polarité est \textbf{correcte}.
    \item \textbf{Lecture négative :} Si l'affichage est \textbf{négatif} (ex: \SI{-1,10}{V}), tu as inversé les bornes (ou ta prévision de polarité était fausse). Inverse les fils : la valeur absolue est la f.é.m.
    \item \textbf{Précision :} La valeur mesurée diffère légèrement de la valeur théorique (ex: \SI{1,08}{V} au lieu de \SI{1,10}{V}) à cause de :
    \begin{itemize}
        \item Non-standardité des solutions (concentration $\neq$ \SI{1}{M}, jonction liquide au pont salin).
        \item Polarisation des électrodes (produits de réaction qui adhèrent mal).
        \item Résistance interne de la pile et du voltmètre (même si faible).
        \item Température différente de \SI{25}{\celsius}.
    \end{itemize}
\end{enumerate}

\begin{savaistu}[La pile au lithium : la championne de la f.é.m.]
En associant le couple le plus réducteur du tableau (\ce{Li+/Li}, \cle{$E^\circ = \SI{-3,05}{V}$}) à un couple oxydant fort comme \ce{Co^3+/Co^2+} (\cle{$\SI{+1,82}{V}$}) ou \ce{MnO2/Mn2O3} (\cle{$\SI{+0,95}{V}$} en milieu basique), on obtient des f.é.m. de \textbf{3 à 3,7 V}. C'est le principe des \textbf{piles lithium-ion} (téléphones, ordinateurs, voitures électriques). Au Cameroun, l'essor du solaire photovoltaïque (panneaux sur les toits de Yaoundé ou Douala) repose sur le stockage par batteries lithium (LiFePO4, \SI{3,2}{V} par élément) pour l'éclairage nocturne. La chimie des potentiels standards guide directement le choix des matériaux pour l'énergie du futur !
\end{savaistu}

\begin{exercice}[Entraînement : Pile \ce{Cd^2+/Cd} / \ce{Ni^2+/Ni}]
Données : \cle{$E^\circ(\ce{Cd^2+/Cd}) = \SI{-0,40}{V}$} ; \cle{$E^\circ(\ce{Ni^2+/Ni}) = \SI{-0,25}{V}$}.
\begin{enumerate}
    \item Déterminer la polarité de la pile (anode, cathode, pôles + et –).
    \item Calculer sa f.é.m. standard.
    \item Écrire les deux demi-équations aux électrodes et l'équation-bilan.
    \item Donner sa notation conventionnelle.
    \item Si on branche un voltmètre (borne COM sur l'anode), quelle valeur affiche-t-il ?
\end{enumerate}
\end{exercice}

\begin{corrige}[Exercice : Pile Cd/Ni]
\begin{enumerate}
    \item \textbf{Classement :} \cle{$\SI{-0,40}{V} < \SI{-0,25}{V}$}. \\
    Anode (–) = \ce{Cd} (oxydation : \ce{Cd -> Cd^2+ + 2e-}). \\
    Cathode (+) = \ce{Ni} (réduction : \ce{Ni^2+ + 2e- -> Ni}).
    \item \textbf{F.é.m. :} \cle{$E^\circ = (-0,25) - (-0,40) = \SI{+0,15}{V}$}.
    \item \textbf{Bilan :} \ce{Cd(s) + Ni^2+(aq) -> Cd^2+(aq) + Ni(s)}.
    \item \textbf{Notation :} \ce{Cd(s) | Cd^2+(aq) || Ni^2+(aq) | Ni(s)}.
    \item \textbf{Voltmètre :} Affiche \textbf{\SI{+0,15}{V}} (valeur positive car borne COM sur anode).
\end{enumerate}
\end{corrige}

\coursec{Généralisation, piles usuelles et synthèse}

Dans la section précédente, tu as appris à prévoir la polarité d'une pile et à calculer sa f.é.m. à partir des potentiels standard des couples mis en jeu. Tu sais désormais que la pile Daniell $\ce{Zn}/\ce{Zn^2+} \parallel \ce{Cu^2+}/\ce{Cu}$ n'est qu'un cas particulier : en associant le couple $\ce{Ag+}/\ce{Ag}$ ($E^\circ = \SI{+0,80}{V}$) au couple $\ce{Zn^2+}/\ce{Zn}$ ($E^\circ = \SI{-0,76}{V}$), on obtient une pile de f.é.m. $E^\circ = 0{,}80 - (-0{,}76) = \SI{1,56}{V}$. La question qui clôt cette leçon est donc naturelle : tout cela se généralise-t-il ? Où sont les piles de ta vie quotidienne — celle de ta lampe torche, de ta radio, de la moto-taxi du quartier ? C'est l'objet de cette synthèse.

\courssub{Généralisation : qu'est-ce qu'une pile ?}

\textbf{Généralisation : qu'est-ce qu'une pile ?}\par
Revenons à l'intuition. Une réaction d'oxydoréduction spontanée, comme celle du zinc plongé dans une solution de sulfate de cuivre, libère de l'énergie : le mélange chauffe. Cette énergie est « gaspillée » en chaleur parce que les électrons passent directement du réducteur à l'oxydant, par contact. L'idée géniale de la pile consiste à \cle{séparer physiquement} les deux couples et à \cle{canaliser} les électrons dans un circuit extérieur : l'énergie de la réaction devient alors travail électrique.

\begin{definition}[Pile électrochimique]
Une \cle{pile électrochimique} est un générateur qui convertit l'énergie chimique d'une réaction d'oxydoréduction \cle{spontanée} en énergie électrique, grâce à un transfert d'électrons \cle{canalisé} dans un circuit extérieur. Elle résulte de l'association de \cle{deux demi-piles} mettant en jeu deux couples oxydant/réducteur \cle{différents}, reliées par un pont salin (ou une paroi poreuse) qui assure la neutralité électrique des solutions.
\end{definition}

Le fonctionnement de toute pile obéit aux mêmes lois, quels que soient les couples :

\begin{itemize}
\item à l'\cle{anode} (pôle $-$) : \textbf{oxydation} du réducteur du couple de plus petit potentiel ; c'est l'électrode qui \emph{fournit} les électrons au circuit ;
\item à la \cle{cathode} (pôle $+$) : \textbf{réduction} de l'oxydant du couple de plus grand potentiel ; c'est l'électrode qui \emph{consomme} les électrons ;
\item dans le circuit extérieur, les \cle{électrons} circulent du pôle $-$ vers le pôle $+$ ; le courant conventionnel circule en sens inverse ;
\item dans le pont salin, les \cle{cations} migrent vers le compartiment de la cathode et les \cle{anions} vers celui de l'anode, compensant les déséquilibres de charges.
\end{itemize}

\begin{attention}[Anode = oxydation, dans tous les cas]
Le mot « anode » désigne \emph{toujours} le siège de l'oxydation, et « cathode » celui de la réduction. Dans une pile (générateur), l'anode est le pôle $-$ et la cathode le pôle $+$. Moyen mnémotechnique : \textbf{A}node-\textbf{O}xydation (voyelles), \textbf{C}athode-\textbf{R}éduction (consonnes). Ne confonds jamais le sens du courant (du $+$ vers le $-$ à l'extérieur) avec celui des électrons (du $-$ vers le $+$).
\end{attention}

\begin{popfigure}
\begin{center}
\begin{tikzpicture}[scale=1.0, every node/.style={font=\small}]
% Bécher anode
\draw[thick, popblue] (0,0) -- (0,2.6) (2.6,0) -- (2.6,2.6);
\draw[thick, popblue] (0,0) -- (2.6,0);
\draw[fill=popblueL, draw=none] (0.06,0.06) rectangle (2.54,1.7);
\node[popblue] at (1.3,1.0) {$\ce{M1^{n+}}$};
% Bécher cathode
\draw[thick, poporange] (5.4,0) -- (5.4,2.6) (8,0) -- (8,2.6);
\draw[thick, poporange] (5.4,0) -- (8,0);
\draw[fill=poporangeL, draw=none] (5.46,0.06) rectangle (7.94,1.7);
\node[poporange] at (6.7,1.0) {$\ce{M2^{n+}}$};
% Électrodes
\draw[fill=popmuted!70, draw=popdark] (1.1,0.4) rectangle (1.5,3.4);
\node[above, popdark] at (1.3,3.4) {$\ce{M1}$};
\draw[fill=popmuted!40, draw=popdark] (6.5,0.4) rectangle (6.9,3.4);
\node[above, popdark] at (6.7,3.4) {$\ce{M2}$};
% Pont salin
\draw[thick, popgreen, fill=popgreenL] (1.9,1.5) -- (1.9,2.3) .. controls (4,3.1) .. (6.1,2.3) -- (6.1,1.5) -- (5.9,1.5) -- (5.9,2.15) .. controls (4,2.85) .. (2.1,2.15) -- (2.1,1.5) -- cycle;
\node[popgreen!60!black] at (4,3.35) {pont salin (\ce{K+}, \ce{NO3-})};
\draw[-{Stealth}, popgreen!60!black] (4.35,2.7) -- (5.0,2.55) node[midway, above, font=\scriptsize] {\ce{K+}};
\draw[-{Stealth}, popgreen!60!black] (3.65,2.7) -- (3.0,2.55) node[midway, above, font=\scriptsize] {\ce{NO3-}};
% Circuit extérieur
\draw[thick, popdark] (1.3,3.4) -- (1.3,4.4) -- (6.7,4.4) -- (6.7,3.4);
\draw[thick, popdark, fill=popgoldL] (3.7,4.25) rectangle (4.3,4.55);
\node[popdark] at (4,4.4) {\scriptsize $R$};
% Électrons
\draw[-{Stealth}, very thick, poppink] (1.7,4.75) -- (3.4,4.75) node[midway, above, poppink] {$e^-$};
% Polarités
\node[popdark, font=\large\bfseries] at (0.7,3.9) {$-$};
\node[popdark, font=\large\bfseries] at (7.3,3.9) {$+$};
% Annotations
\node[align=center, popblue!60!black] at (1.3,-0.9) {\textbf{Anode} (pôle $-$)\\ oxydation : $\ce{M1 -> M1^{n+} + n e-}$};
\node[align=center, poporange!70!black] at (6.7,-0.9) {\textbf{Cathode} (pôle $+$)\\ réduction : $\ce{M2^{n+} + n e- -> M2}$};
\node[align=center] at (4,-2.1) {Bilan : $\ce{M1 + M2^{n+} -> M1^{n+} + M2}$ \quad;\quad $E^\circ_{\text{pile}} = E^\circ(\ce{M2^{n+}}/\ce{M2}) - E^\circ(\ce{M1^{n+}}/\ce{M1}) > 0$};
\end{tikzpicture}
\end{center}
\legende{Schéma récapitulatif d'une pile quelconque $\ce{M1} \mid \ce{M1^{n+}} \parallel \ce{M2^{n+}} \mid \ce{M2}$, avec $E^\circ(\ce{M2^{n+}}/\ce{M2}) > E^\circ(\ce{M1^{n+}}/\ce{M1})$. Toutes les grandeurs essentielles y figurent : polarités, demi-équations, sens des électrons, migration des ions du pont salin, f.é.m.}
\end{popfigure}

\courssub{Diversité des piles usuelles}

Le principe de Daniell (deux couples, deux compartiments) est peu pratique pour la vie courante : qui transporterait deux béchers de solutions dans sa poche ? Les piles commerciales condensent le même principe dans un boîtier étanche, avec des électrolytes gélifiés ou pâteux.

\textbf{La pile saline (pile Leclanché).} C'est la pile « classique » des lampes torches. L'anode est un cylindre de \cle{zinc} (qui sert aussi de boîtier), la cathode une tige de \cle{carbone} (graphite) entourée de dioxyde de manganèse $\ce{MnO2}$ ; l'électrolyte est une pâte de chlorure d'ammonium $\ce{NH4Cl}$. Les couples mis en jeu sont $\ce{Zn^2+}/\ce{Zn}$ et, en milieu ammoniacal, un couple du manganèse ($\ce{MnO2}/\ce{Mn2O3}$ simplifié). Sa f.é.m. vaut environ \SI{1,5}{V}. Le zinc est consommé au cours du fonctionnement : la pile « s'use » littéralement par sa paroi, d'où les fuites parfois observées dans les vieilles lampes.

\textbf{La pile alcaline.} Même principe (zinc et $\ce{MnO2}$), mais l'électrolyte est une solution concentrée d'hydroxyde de potassium $\ce{KOH}$ (d'où le nom « alcaline »). Elle délivre aussi environ \SI{1,5}{V}, mais avec une capacité bien supérieure et une meilleure tenue en courant : c'est la pile des jouets, des télécommandes et des radios.

\textbf{La pile bouton.} Miniaturisée (montres, calculatrices, aides auditives), elle utilise par exemple les couples $\ce{Zn(OH)4^2-}/\ce{Zn}$ et $\ce{Ag2O}/\ce{Ag}$ (pile à l'argent, f.é.m. $\approx \SI{1,55}{V}$) ou, autrefois, un couple au mercure — aujourd'hui abandonné pour sa toxicité.

\textbf{Les accumulateurs.} Un \cle{accumulateur} est une pile \cle{rechargeable} : la réaction d'oxydoréduction y est réversible. En décharge, il fonctionne comme une pile (il fournit du courant) ; en charge, un générateur extérieur impose le courant en sens inverse et force la réaction inverse (électrolyse). L'accumulateur au plomb des voitures et motos associe les couples $\ce{PbO2}/\ce{PbSO4}$ et $\ce{PbSO4}/\ce{Pb}$ dans l'acide sulfurique ; chaque élément délivre environ \SI{2}{V}, et six éléments en série donnent la batterie de \SI{12}{V}.

\begin{savaistu}[Accumulateurs et énergie solaire au Cameroun]
Les batteries des motos-taxi de Douala ou de Maroua, et surtout celles des installations solaires domestiques qui éclairent de nombreux villages non raccordés au réseau ENEO, fonctionnent exactement sur le principe de la pile Daniell : conversion d'énergie chimique en énergie électrique par une réaction d'oxydoréduction. La différence est que la réaction y est \emph{réversible} : le jour, le panneau solaire « recharge » la batterie en imposant la réaction inverse ; la nuit, la batterie se décharge comme une pile. Comprendre les piles, c'est comprendre le stockage de l'énergie solaire !
\end{savaistu}

\courssub{Intérêt énergétique et applications}

\begin{aretenir}[Conversion directe d'énergie]
Une pile réalise la conversion \cle{directe} de l'énergie chimique en énergie électrique, sans passer par la chaleur ni par des pièces en mouvement (contrairement à une centrale thermique ou à un groupe électrogène). Le rendement de cette conversion est élevé, et le dispositif est silencieux, portable et immédiatement disponible.
\end{aretenir}

Au Cameroun, les applications sont quotidiennes : lampes torches et lampes tempêtes rechargeables pendant les délestages, radios à piles dans les zones rurales, téléphones portables (batteries lithium-ion, accumulateurs modernes), onduleurs et batteries solaires pour l'alimentation de secours des homes et des centres de santé. La pile est l'un des rares générateurs d'électricité vraiment individuels et transportables.

\courssub{Tableau de synthèse des piles étudiées}

Le tableau ci-dessous rassemble les piles rencontrées dans cette leçon. Vérifie chaque f.é.m. par le calcul $E^\circ_{\text{pile}} = E^\circ(+) - E^\circ(-)$.

\begin{center}
\begin{tabularx}{\linewidth}{|l|X|X|c|c|}
\hline
\rowcolor{popblueL}
\textbf{Pile} & \textbf{Couple au pôle $-$} & \textbf{Couple au pôle $+$} & \textbf{Calcul} & $\boldsymbol{E^\circ_{\text{pile}}}$ \\
\hline
Daniell & $\ce{Zn^2+}/\ce{Zn}$ ($\SI{-0,76}{V}$) & $\ce{Cu^2+}/\ce{Cu}$ ($\SI{+0,34}{V}$) & $0{,}34-(-0{,}76)$ & \SI{1,10}{V} \\
\hline
Zinc--argent & $\ce{Zn^2+}/\ce{Zn}$ ($\SI{-0,76}{V}$) & $\ce{Ag+}/\ce{Ag}$ ($\SI{+0,80}{V}$) & $0{,}80-(-0{,}76)$ & \SI{1,56}{V} \\
\hline
Fer--cuivre & $\ce{Fe^2+}/\ce{Fe}$ ($\SI{-0,44}{V}$) & $\ce{Cu^2+}/\ce{Cu}$ ($\SI{+0,34}{V}$) & $0{,}34-(-0{,}44)$ & \SI{0,78}{V} \\
\hline
Zinc--fer & $\ce{Zn^2+}/\ce{Zn}$ ($\SI{-0,76}{V}$) & $\ce{Fe^2+}/\ce{Fe}$ ($\SI{-0,44}{V}$) & $-0{,}44-(-0{,}76)$ & \SI{0,32}{V} \\
\hline
Saline (Leclanché) & zinc & $\ce{MnO2}$/carbone & --- & $\approx \SI{1,5}{V}$ \\
\hline
Alcaline & zinc & $\ce{MnO2}$ & --- & $\approx \SI{1,5}{V}$ \\
\hline
Bouton (argent) & zinc & $\ce{Ag2O}/\ce{Ag}$ & --- & $\approx \SI{1,55}{V}$ \\
\hline
Accumulateur au plomb (1 élément) & $\ce{PbSO4}/\ce{Pb}$ & $\ce{PbO2}/\ce{PbSO4}$ & --- & $\approx \SI{2}{V}$ \\
\hline
\end{tabularx}
\end{center}

\begin{attention}[Deux pièges classiques]
\begin{itemize}
\item Deux couples de potentiels \emph{proches} donnent une f.é.m. \emph{faible} (pile zinc--fer : \SI{0,32}{V} seulement). La f.é.m. dépend de l'\emph{écart} des potentiels, pas de leur valeur absolue.
\item Une f.é.m. est \emph{positive par construction} : on écrit toujours $E^\circ_{\text{pile}} = E^\circ(+) - E^\circ(-) > 0$. Si ton calcul donne une valeur négative, c'est que tu as inversé les pôles.
\end{itemize}
\end{attention}

\begin{exoresolu}[Conception d'une pile de f.é.m. imposée]
On dispose des couples suivants : $\ce{Mg^2+}/\ce{Mg}$ ($E^\circ = \SI{-2,37}{V}$), $\ce{Zn^2+}/\ce{Zn}$ ($E^\circ = \SI{-0,76}{V}$), $\ce{Pb^2+}/\ce{Pb}$ ($E^\circ = \SI{-0,13}{V}$), $\ce{Cu^2+}/\ce{Cu}$ ($E^\circ = \SI{+0,34}{V}$). On veut réaliser une pile de f.é.m. la plus grande possible, puis donner sa notation conventionnelle, ses demi-équations et son bilan.
\tcblower
\textbf{1) Choix des couples.} La f.é.m. est maximale pour l'écart de potentiel maximal : on associe le couple de plus petit potentiel ($\ce{Mg^2+}/\ce{Mg}$) à celui de plus grand potentiel ($\ce{Cu^2+}/\ce{Cu}$).

\textbf{2) Polarité.} $E^\circ(\ce{Cu^2+}/\ce{Cu}) > E^\circ(\ce{Mg^2+}/\ce{Mg})$ : le cuivre est le pôle $+$ (cathode), le magnésium le pôle $-$ (anode).

\textbf{3) Demi-équations.}
\begin{align*}
\text{Anode (oxydation) : } & \ce{Mg -> Mg^2+ + 2 e-}\\
\text{Cathode (réduction) : } & \ce{Cu^2+ + 2 e- -> Cu}
\end{align*}

\textbf{4) Bilan de fonctionnement} (les électrons s'éliminent, 2 contre 2) :
\[ \ce{Mg + Cu^2+ -> Mg^2+ + Cu} \]
Réaction spontanée et quasi totale, car le cuivre est un oxydant nettement plus fort que $\ce{Mg^2+}$.

\textbf{5) F.é.m. standard :}
\[ E^\circ_{\text{pile}} = E^\circ(+) - E^\circ(-) = 0{,}34 - (-2{,}37) = \SI{2,71}{V}. \]

\textbf{6) Notation conventionnelle} (pôle $-$ à gauche) :
\[ \ce{Mg} \mid \ce{Mg^2+} \parallel \ce{Cu^2+} \mid \ce{Cu} \]
\end{exoresolu}

\courssub{Fiche réflexe Probatoire et bilan des savoir-faire}

\begin{methode}[Fiche réflexe : traiter n'importe quel exercice de pile]
\begin{enumerate}
\item \textbf{Écrire} les deux couples mis en jeu avec leurs potentiels standard $E^\circ$.
\item \textbf{Identifier les pôles} : le pôle $+$ (cathode) correspond au couple de $E^\circ$ \emph{le plus grand} ; le pôle $-$ (anode) à l'autre.
\item \textbf{Écrire les demi-équations} : oxydation du réducteur à l'anode, réduction de l'oxydant à la cathode.
\item \textbf{Écrire le bilan} en égalant les nombres d'électrons échangés.
\item \textbf{Calculer la f.é.m.} : $E^\circ_{\text{pile}} = E^\circ(+) - E^\circ(-)$.
\item \textbf{Écrire la notation conventionnelle} : pôle $-$ à gauche, pôle $+$ à droite, double barre $\parallel$ pour le pont salin.
\end{enumerate}
\end{methode}

À l'issue de cette leçon, tu dois être capable de : \textbf{prévoir} la réaction spontanée entre deux couples (règle du gamma ou comparaison des $E^\circ$), \textbf{écrire} les demi-équations et le bilan de fonctionnement, \textbf{noter} une pile conventionnellement, \textbf{calculer} sa f.é.m. standard, et \textbf{réaliser} expérimentalement une pile en mesurant ses pôles (voltmètre) et sa f.é.m. (circuit ouvert).

\begin{exercice}[Pile fer--argent]
On réalise une pile en associant les demi-piles $\ce{Fe^2+}/\ce{Fe}$ ($E^\circ = \SI{-0,44}{V}$) et $\ce{Ag+}/\ce{Ag}$ ($E^\circ = \SI{+0,80}{V}$).
\begin{enumerate}
\item Déterminer la polarité de la pile et écrire sa notation conventionnelle.
\item Écrire les demi-équations aux électrodes et le bilan de fonctionnement.
\item Calculer la f.é.m. standard de cette pile.
\end{enumerate}
\end{exercice}

\begin{corrige}[Exercice : pile fer--argent]
\textbf{1.} $E^\circ(\ce{Ag+}/\ce{Ag}) > E^\circ(\ce{Fe^2+}/\ce{Fe})$ : l'argent est le pôle $+$ (cathode), le fer le pôle $-$ (anode). Notation : $\ce{Fe} \mid \ce{Fe^2+} \parallel \ce{Ag+} \mid \ce{Ag}$.

\textbf{2.} Anode : $\ce{Fe -> Fe^2+ + 2 e-}$. Cathode : $\ce{Ag+ + e- -> Ag}$. En multipliant la seconde par 2 puis en sommant : $\ce{Fe + 2 Ag+ -> Fe^2+ + 2 Ag}$.

\textbf{3.} $E^\circ_{\text{pile}} = 0{,}80 - (-0{,}44) = \SI{1,24}{V}$.
\end{corrige}

\courssub{Piles usagées et environnement}

Une pile « morte » n'est pas un déchet ordinaire. Les piles salines et alcalines contiennent du \cle{zinc} et du \cle{manganèse} ; certaines piles bouton anciennes contiennent du \cle{mercure}, et les accumulateurs rechargeables au nickel-cadmium du \cle{cadmium} — des métaux lourds toxiques qui, jetés dans la nature, contaminent les sols et les nappes phréatiques, puis remontent la chaîne alimentaire jusqu'à l'homme. L'accumulateur au plomb, lui, contient du plomb et de l'acide sulfurique.

\begin{attention}[Que faire des piles usagées ?]
Ne jette \textbf{jamais} une pile usagée dans la nature, dans les ordures ménagères brûlées à l'air libre, ni au feu : la chaleur fait éclater le boîtier et libère des vapeurs toxiques. Conserve les piles usagées dans un bocal fermé et dépose-les dans un point de collecte spécifique (certains commerces, ateliers de recyclage de batteries). Les batteries de motos et de voitures se rapportent chez le vendeur ou le garagiste, qui les reprend pour recyclage du plomb. Protéger l'environnement, c'est aussi protéger ta santé et celle de ton quartier.
\end{attention}

\begin{aretenir}[L'essentiel de la leçon]
\begin{itemize}
\item Toute pile associe \textbf{deux demi-piles de couples différents} ; son fonctionnement est une réaction d'oxydoréduction \textbf{spontanée} dont les électrons sont \textbf{canalisés} dans un circuit extérieur.
\item Anode $=$ oxydation $=$ pôle $-$ (couple de plus petit $E^\circ$) ; cathode $=$ réduction $=$ pôle $+$ (couple de plus grand $E^\circ$).
\item $E^\circ_{\text{pile}} = E^\circ(+) - E^\circ(-) > 0$ ; notation : pôle $-$ à gauche, pôle $+$ à droite.
\item Piles usuelles : saline et alcaline ($\approx \SI{1,5}{V}$), bouton ($\approx \SI{1,55}{V}$), accumulateur au plomb ($\approx \SI{2}{V}$ par élément, rechargeable).
\item Une pile convertit directement l'énergie chimique en énergie électrique ; les piles usagées se recyclent, jamais ne se jettent ni ne se brûlent.
\end{itemize}
\end{aretenir}

Tu maîtrises maintenant l'ensemble du chapitre : de l'expérience de Daniell à la classification quantitative des couples, en passant par l'électrode à hydrogène de référence. Ces savoirs te serviront directement dans la suite du programme : l'électrolyse, qui est en quelque sorte la « pile à l'envers », et la corrosion des métaux, qui n'est autre qu'une pile naturelle non désirée.

\newpage

\coursec{Travaux pratiques}

\courssub{Objectifs du TP}
\begin{itemize}
    \item Réaliser expérimentalement la pile Daniell (\ce{Zn}|\ce{Zn^2+}||\ce{Cu^2+}|\ce{Cu}) et d'autres piles à partir de couples métalliques disponibles (\ce{Zn}|\ce{Fe^2+}, \ce{Fe}|\ce{Cu^2+}, \ce{Cu}|\ce{Ag^+}).
    \item Identifier les pôles (anode et cathode) à l'aide d'un voltmètre et déterminer le sens du courant électrique dans le circuit externe.
    \item Mesurer la force électromotrice (f.é.m.) de chaque pile réalisée.
    \item Calculer les f.é.m. théoriques à partir des potentiels standards d'oxydoréduction (\cle{$E^\circ$}) et comparer aux valeurs expérimentales en expliquant les écarts (concentrations non standard, résistance interne, jonction liquide, pureté des électrodes).
    \item Écrire les équations des demi-réactions aux électrodes et l'équation-bilan de fonctionnement pour chaque pile.
\end{itemize}

\courssub{Matériel et produits}
\begin{center}
\setlength{\tabcolsep}{4pt}
\begin{tabular}{|l|p{5.5cm}|p{5.5cm}|}\hline
\rowcolor{popblueL}
\textbf{Matériel verrier / consommables} & \textbf{Produits chimiques} & \textbf{Appareillage / Divers} \\\hline
\begin{minipage}{\linewidth}
\begin{itemize}\itemsep2pt
    \item 4 béchers \SI{100}{mL} (ou godets en plastique récupérés, type pots de yaourt lavés)
    \item 4 pinces crocodiles avec fils de connexion (ou fils électriques dénudés)
    \item Bandes de papier filtre (largeur \SI{1}{cm}, longueur \SI{10}{cm}) pour pont salin
    \item Tube en U (optionnel) + gel d'agar-agar \SI{2}{\%} saturé en \ce{KNO3}
    \item Lames de papier émeri (grain fin) ou lime métallique
    \item Baguettes de verre pour agitation
    \item Pipettes graduées \SI{10}{mL} ou seringues \SI{10}{mL}
    \item Bécher \SI{250}{mL} pour récupération des déchets
\end{itemize}
\end{minipage} &
\begin{itemize}\itemsep2pt
    \item Lame de zinc (Zn) : \SI{3}{cm} $\times$ \SI{1}{cm} (récupérée sur pile usagée ou toit en tôle neuve décapée)
    \item Lame de cuivre (Cu) : \SI{3}{cm} $\times$ \SI{1}{cm} (fil électrique rigide dénudé ou plaque de cuivre)
    \item Clou en fer (Fe) ou lame de fer doux : \SI{3}{cm} (décapé à l'acide chlorhydrique dilué puis rincé)
    \item Fil d'argent (Ag) ou pièce de monnaie en alliage Ag (optionnel, pour pile Cu/Ag)
    \item Solution \ce{ZnSO4} \SI{0,5}{mol.L^{-1}} (\approx \SI{72}{g.L^{-1}})
    \item Solution \ce{CuSO4,5H2O} \SI{0,5}{mol.L^{-1}} (\approx \SI{125}{g.L^{-1}})
    \item Solution \ce{FeSO4,7H2O} \SI{0,5}{mol.L^{-1}} (\approx \SI{140}{g.L^{-1}})
    \item Solution \ce{AgNO3} \SI{0,1}{mol.L^{-1}} (\approx \SI{17}{g.L^{-1}}, optionnel)
    \item Solution \ce{KNO3} saturée (\approx \SI{300}{g.L^{-1}}) ou \ce{NaCl} \SI{1}{mol.L^{-1}} pour pont salin
    \item Eau distillée (ou eau de pluie filtrée et bouillie)
\end{itemize} &
\begin{itemize}\itemsep2pt
    \item Voltmètre numérique (multimètre) calibre \SI{2}{V} ou \SI{20}{V} (impédance d'entrée $\ge$ \SI{10}{M\Omega})
    \item Lunettes de protection obligatoires
    \item Gants en nitrile (recommandés)
    \item Chronomètre (ou montre)
    \item Support de TP pour notes
\end{itemize} \\\hline
\end{tabular}
\end{center}

\courssub{Sécurité}
\begin{attention}[Consignes de sécurité indispensables]
\textbf{Pictogrammes et risques identifiés :}
\begin{itemize}
    \item \ce{CuSO4,5H2O} : \textbf{SGH07} (Irritant -- \emph{Attention}), \textbf{SGH09} (Dangereux pour l'environnement -- \emph{Attention}). \textit{Risque :} Irritation yeux/peau, toxique pour les organismes aquatiques.
    \item \ce{FeSO4,7H2O} : \textbf{SGH07} (Irritant). \textit{Risque :} Irritation yeux/peau.
    \item \ce{ZnSO4} : \textbf{SGH07} (Irritant), \textbf{SGH09} (Environnement).
    \item \ce{AgNO3} (si utilisé) : \textbf{SGH05} (Corrosif -- \emph{Danger}), \textbf{SGH09} (Environnement). \textit{Risque :} Brûlures graves, taches noires indélébiles sur peau/vêtements (réduction en Ag métallique).
    \item \ce{KNO3} / \ce{NaCl} : Faible toxicité, mais éviter ingestion.
\end{itemize}
\textbf{Gestes à adopter (contexte lycée camerounais) :}
\begin{enumerate}
    \item \textbf{Port obligatoire des lunettes} de protection (ou lunettes de soleil enveloppantes si pas de lunettes de labo) pendant toute la manipulation.
    \item \textbf{Manipulation soigneuse} des solutions de sulfate de cuivre (bleu) et de nitrate d'argent (incolore mais tachant) : ne pas renverser, essuyer toute projection immédiatement à l'eau.
    \item \textbf{Interdiction formelle} de porter les mains à la bouche/yeux ; se laver les mains soigneusement au savon après le TP.
    \item \textbf{Récupération systématique} des solutions usées dans le bécher « Déchets métaux lourds » (ne pas jeter à l'évier : \ce{Cu^2+}, \ce{Zn^2+}, \ce{Ag+} polluent l'environnement et les fosses septiques). L'enseignant gérera l'élimination selon le protocole établissement (précipitation par \ce{NaOH} / \ce{Na2S} puis filtration).
    \item En cas de projection dans les yeux : rinçage immédiat $\ge$ \SI{15}{min} au point d'eau le plus proche, avis médical.
\end{enumerate}
\end{attention}

\courssub{Schéma du dispositif}
\begin{popfigure}
\begin{tikzpicture}[
    scale=0.9,
    beaker/.style={draw=popdark, thick, fill=popblueL!20, minimum width=3.5cm, minimum height=5cm, rounded corners=3pt},
    electrode/.style={draw=popdark, very thick, fill=popmuted!30, minimum width=0.6cm, minimum height=4cm},
    wire/.style={draw=poporange, very thick, -{Stealth[length=3mm]}},
    saltbridge/.style={draw=poppurple, very thick, fill=poppurpleL!50, rounded corners=2pt},
    voltmeter/.style={draw=popdark, thick, fill=popcream, minimum width=2cm, minimum height=1.5cm, rounded corners=2pt},
    label/.style={font=\small, text=popdark},
    ion/.style={font=\tiny, text=poporange, ->, >=Stealth}
]
% Bécher Gauche (Anode)
\node[beaker] (beakL) at (0,0) {};
\node[electrode, anchor=north] (Zn) at (beakL.north) {Zn};
\node[label] at (Zn.north west) {Anode (-)};
\node[label, align=center] at (-1.75, -1.5) {Solution\\ \ce{ZnSO4} \SI{0,5}{M}};
\draw[ion] (-1.5, -0.5) -- ++(0.5,0) node[right] {\ce{Zn^2+}};
\draw[ion] (-1.5, -1.0) -- ++(-0.5,0) node[left] {\ce{SO4^2-}};

% Bécher Droit (Cathode)
\node[beaker] (beakR) at (9,0) {};
\node[electrode, anchor=north] (Cu) at (beakR.north) {Cu};
\node[label] at (Cu.north east) {Cathode (+)};
\node[label, align=center] at (10.75, -1.5) {Solution\\ \ce{CuSO4} \SI{0,5}{M}};
\draw[ion] (10.5, -0.5) -- ++(-0.5,0) node[left] {\ce{Cu^2+}};
\draw[ion] (10.5, -1.0) -- ++(0.5,0) node[right] {\ce{SO4^2-}};

% Pont salin
\draw[saltbridge] (beakL.north east) -- ++(0, 1.5) -| (beakR.north west) -- cycle;
\node[label, rotate=90, align=center] at (4.5, 2.2) {Pont salin\\ (\ce{KNO3} sat.)};

% Voltmètre
\node[voltmeter] (volt) at (4.5, 4.5) {};
\node[label] at (volt.center) {V};
\node[label] at (volt.north) {Voltmètre};

% Connexions
\draw[wire] (Zn.north west) -- ++(-0.5, 0.5) -| (volt.west) node[left, midway, label] {COM (-)};
\draw[wire] (Cu.north east) -- ++(0.5, 0.5) -| (volt.east) node[right, midway, label] {V (+)};

% Légendes réactions
\node[label, align=left, text width=3.5cm] at (-1.75, -3.5) {\textbf{Demi-réaction anode :}\\[2pt] \ce{Zn -> Zn^2+ + 2e-}};
\node[label, align=right, text width=3.5cm] at (10.75, -3.5) {\textbf{Demi-réaction cathode :}\\[2pt] \ce{Cu^2+ + 2e- -> Cu}};
\node[label, align=center, font=\bfseries] at (4.5, -4.5) {\textbf{Pile Daniell : \ce{Zn(s) + Cu^2+(aq) -> Zn^2+(aq) + Cu(s)}}};
\end{tikzpicture}
\legende{Schéma de la pile Daniell réalisée au laboratoire. L'électrode de zinc (anode, pôle -) se dissout, l'électrode de cuivre (cathode, pôle +) grossit. Le pont salin assure la neutralité électrique. Le voltmètre (borne COM sur l'anode) affiche une tension positive.}
\end{popfigure}

\courssub{Protocole}
\begin{enumerate}
    \item \textbf{Préparation des électrodes :} Nettoyer soigneusement les lames de Zn, Cu, Fe (et Ag le cas échéant) avec le papier émeri pour éliminer la couche d'oxyde/rouille. Rincer à l'eau distillée, sécher avec papier absorbant sans toucher la surface active avec les doigts. Identifier chaque lame (marqueur indélébile sur la partie non immergée).
    \item \textbf{Préparation des demi-piles :} Verser \SI{50}{mL} de chaque solution (\ce{ZnSO4}, \ce{CuSO4}, \ce{FeSO4}, \ce{AgNO3}) dans 4 béchers distincts étiquetés. Plonger l'électrode correspondante dans chaque solution (Zn dans \ce{ZnSO4}, Cu dans \ce{CuSO4}, etc.) de façon à ce qu'environ \SI{2}{cm} soient immergés. Fixer les pinces crocodiles sur la partie émergée.
    \item \textbf{Préparation du pont salin :} Imbiber une bande de papier filtre dans la solution de \ce{KNO3} saturée (ou \ce{NaCl} \SI{1}{M}) pendant \SI{2}{min}. L'égoutter légèrement. \textit{Alternative sans tube en U :} Plier la bande en U et plonger chaque extrémité dans deux béchers adjacents (ex: Zn et Cu). Vérifier qu'il n'y a pas de bulle d'air bloquant le contact liquide.
    \item \textbf{Montage de la pile Daniell (Référence) :}
    \begin{enumerate}
        \item Relier la pince de l'électrode \textbf{Zn} à la borne \textbf{COM (noire / -)} du voltmètre.
        \item Relier la pince de l'électrode \textbf{Cu} à la borne \textbf{V / $\Omega$ / mA (rouge / +)} du voltmètre.
        \item Mettre le pont salin entre le bécher Zn et le bécher Cu.
        \item Régler le voltmètre sur le calibre \textbf{DC \SI{2}{V}} (ou \SI{20}{V} si incertain).
        \item Relever la tension $U$ affichée dès stabilisation ($\approx$ \SI{10}{s}). Noter le signe ($+$ ou $-$).
    \end{enumerate}
    \item \textbf{Identification des pôles :}
    \begin{itemize}
        \item Si $U > 0$ : La borne rouge (+) est à un potentiel plus élevé que la borne noire (-). L'électrode Cu est le pôle \textbf{+ (cathode)}, l'électrode Zn est le pôle \textbf{- (anode)}. Le courant conventionnel sort par Cu, entre par Zn.
        \item Si $U < 0$ : Inverser les fils sur le voltmètre ou inverser mentalement les pôles.
    \end{itemize}
    \item \textbf{Autres piles :} Reproduire l'opération (étapes 3 et 4) pour les couples suivants en changeant les béchers et le pont salin (utiliser un pont salin neuf à chaque fois pour éviter la contamination croisée) :
    \begin{itemize}
        \item Pile \ce{Zn}|\ce{Fe^2+} (Zn dans \ce{ZnSO4}, Fe dans \ce{FeSO4})
        \item Pile \ce{Fe}|\ce{Cu^2+} (Fe dans \ce{FeSO4}, Cu dans \ce{CuSO4})
        \item Pile \ce{Cu}|\ce{Ag^+} (Cu dans \ce{CuSO4}, Ag dans \ce{AgNO3}) -- \textit{si Ag disponible}
    \end{itemize}
    \item \textbf{Nettoyage :} Démonter. Rincer les électrodes à l'eau distillée. Verser toutes les solutions usées dans le bécher « Déchets ». Ranger le matériel.
\end{enumerate}

\courssub{Observations et résultats attendus}
\begin{center}
\setlength{\tabcolsep}{3pt}
\begin{tabular}{|c|c|c|c|c|c|}\hline
\rowcolor{popblueL}
\textbf{Pile réalisée} & \textbf{Branchement V/mètre} & \textbf{Tension mesurée $U_{\text{exp}}$ (V)} & \textbf{Pôle + (Cathode)} & \textbf{Pôle - (Anode)} & \textbf{Sens courant externe} \\\hline
Daniell \ce{Zn}|\ce{Cu^2+} & COM $\leftarrow$ Zn ; V/$\Omega$ $\leftarrow$ Cu & \textbf{+1,05 à +1,10} & Cu & Zn & Cu $\to$ Zn \\\hline
\ce{Zn}|\ce{Fe^2+} & COM $\leftarrow$ Zn ; V/$\Omega$ $\leftarrow$ Fe & \textbf{+0,70 à +0,78} & Fe & Zn & Fe $\to$ Zn \\\hline
\ce{Fe}|\ce{Cu^2+} & COM $\leftarrow$ Fe ; V/$\Omega$ $\leftarrow$ Cu & \textbf{+0,75 à +0,82} & Cu & Fe & Cu $\to$ Fe \\\hline
\ce{Cu}|\ce{Ag^+} & COM $\leftarrow$ Cu ; V/$\Omega$ $\leftarrow$ Ag & \textbf{+0,42 à +0,48} & Ag & Cu & Ag $\to$ Cu \\\hline
\end{tabular}
\end{center}
\vspace{0,5em}
\begin{aretenir}[Observations visuelles qualitatives]
\begin{itemize}
    \item \textbf{Pile Daniell :} La solution incolore \ce{ZnSO4} le reste. La solution bleue \ce{CuSO4} se décolore progressivement (disparition \ce{Cu^2+}). L'électrode de Zn s'amincit/rugosit ; l'électrode de Cu se recouvre d'un dépôt rouge-orangé brillant (Cu métallique).
    \item \textbf{Pile Zn/Fe :} Dépôt grisâtre (Fe) sur l'électrode de Fe ; Zn s'use.
    \item \textbf{Pile Fe/Cu :} Dépôt rouge (Cu) sur l'électrode de Cu ; le clou Fe se corrode (peut noircir par formation de \ce{Fe^2+} puis oxydation air en \ce{Fe^3+}).
    \item \textbf{Pile Cu/Ag :} Dépôt blanc brillant/cristallin (Ag) sur le fil Cu ; solution \ce{AgNO3} incolore, Cu se teinte en bleu (\ce{Cu^2+} formé).
\end{itemize}
\end{aretenir}

\courssub{Exploitation}
\begin{exoresolu}[Questions guidées et réponses]
\textbf{1. Écrire les demi-équations électroniques aux électrodes et l'équation-bilan pour la pile Daniell.}
\tcblower
\textbf{Anode (pôle - , oxydation) :} \ce{Zn(s) -> Zn^2+(aq) + 2e-}
\textbf{Cathode (pôle + , réduction) :} \ce{Cu^2+(aq) + 2e- -> Cu(s)}
\textbf{Équation-bilan :} \ce{Zn(s) + Cu^2+(aq) -> Zn^2+(aq) + Cu(s)} \quad (Transfert de 2 électrons)

\textbf{2. En utilisant le tableau des potentiels standards ci-dessous, calculer la f.é.m. standard $E^\circ_{\text{pile}}$ pour chaque pile réalisée.}
\begin{center}
\begin{tabular}{|c|c|}\hline
\rowcolor{popblueL} \textbf{Couple} & \textbf{$E^\circ$ (V/ENH)} \\\hline
\ce{Zn^2+/Zn} & -0,76 \\\hline
\ce{Fe^2+/Fe} & -0,44 \\\hline
\ce{Cu^2+/Cu} & +0,34 \\\hline
\ce{Ag+/Ag} & +0,80 \\\hline
\end{tabular}
\end{center}
\tcblower
Rappel : $E^\circ_{\text{pile}} = E^\circ_{\text{cathode}} - E^\circ_{\text{anode}}$ (cathode = couple au potentiel le plus élevé).
\begin{align*}
\text{Daniell (Zn/Cu)} &: E^\circ = 0,34 - (-0,76) = \mathbf{+1,10\ V} \\
\text{Zn/Fe} &: E^\circ = -0,44 - (-0,76) = \mathbf{+0,32\ V} \\
\text{Fe/Cu} &: E^\circ = 0,34 - (-0,44) = \mathbf{+0,78\ V} \\
\text{Cu/Ag} &: E^\circ = 0,80 - 0,34 = \mathbf{+0,46\ V}
\end{align*}

\textbf{3. Comparer les valeurs expérimentales ($U_{\text{exp}}$) aux valeurs théoriques ($E^\circ_{\text{pile}}$). Commenter les écarts.}
\tcblower
\begin{center}
\setlength{\tabcolsep}{4pt}
\begin{tabular}{|c|c|c|c|}\hline
\rowcolor{popblueL} \textbf{Pile} & $E^\circ_{\text{théo}}$ (V) & $U_{\text{exp}}$ (V) & $\Delta U$ (V) \\\hline
Daniell & 1,10 & $\approx$ 1,07 & -0,03 \\\hline
Zn/Fe & 0,32 & $\approx$ 0,75 & \textbf{+0,43 !} \\\hline
Fe/Cu & 0,78 & $\approx$ 0,78 & $\approx$ 0 \\\hline
Cu/Ag & 0,46 & $\approx$ 0,45 & -0,01 \\\hline
\end{tabular}
\end{center}
\textbf{Analyse :}
\begin{itemize}
    \item \textbf{Pile Daniell, Fe/Cu, Cu/Ag :} Écarts faibles ($< \SI{0,05}{V}$). Dus principalement aux \textbf{concentrations non standard} (\SI{0,5}{M} au lieu de \SI{1}{M}), à la \textbf{résistance interne} de la pile (pont salin papier filtre = résistance ohmique non négligeable, chute de tension aux bornes du voltmètre si courant non nul, bien que $R_{\text{voltmètre}}$ soit grand) et à la \textbf{potentiel de jonction liquide} au niveau du pont salin.
    \item \textbf{Pile Zn/Fe :} \textbf{Écart énorme et signe potentiellement inversé selon l'état de surface du fer.} Le potentiel standard \ce{Fe^2+/Fe} (-0,44 V) suppose du fer pur. Nos clous/lames sont en \textbf{acier} (alliage Fe-C) avec impuretés, et surtout recouverts d'oxydes/hydroxydes (\ce{Fe2O3}, \ce{Fe(OH)2}) même après décapage. La corrosion du fer en milieu neutre/acide (\ce{Fe -> Fe^2+ + 2e-} couplée à \ce{O2 + 2H2O + 4e- -> 4OH-}) crée des potentials mixtes. Le fer « passivé » ou corrodé peut avoir un potentiel effectif bien plus noble (proche de 0 V ou +), expliquant une f.é.m. mesurée proche de \ce{Zn/Cu}. \textbf{C'est un excellent exemple des limites des potentiels standards sur matériaux réels.}
\end{itemize}

\textbf{4. Pourquoi utilise-t-on un pont salin en \ce{KNO3} (ou \ce{NaCl}) et non en \ce{CuSO4} ou \ce{AgNO3} ?}
\tcblower
Les ions du pont salin (\ce{K+}, \ce{NO3-}, \ce{Na+}, \ce{Cl-}) doivent être \textbf{inertes} (ne pas réagir aux électrodes, ne pas précipiter avec les ions en solution). \ce{KNO3} est idéal : fortes mobilités ioniques voisines (limite potentiel de jonction), solubilité élevée, pas de précipitation avec \ce{SO4^2-}, \ce{Zn^2+}, \ce{Cu^2+}, \ce{Fe^2+}, \ce{Ag+}. \ce{AgNO3} ferait précipiter \ce{AgCl} (si \ce{Cl-}) ou réagirait au Zn. \ce{CuSO4} modifierait la composition de l'électrolyte.

\textbf{5. Écrire l'équation-bilan pour la pile Fe/Cu. Prédire l'évolution de la couleur des solutions.}
\tcblower
Anode (Fe) : \ce{Fe(s) -> Fe^2+(aq) + 2e-}
Cathode (Cu) : \ce{Cu^2+(aq) + 2e- -> Cu(s)}
Bilan : \ce{Fe(s) + Cu^2+(aq) -> Fe^2+(aq) + Cu(s)}
\textbf{Évolution :} La solution bleue (\ce{Cu^2+}) se décolore. La solution verte pâle (\ce{Fe^2+}) s'intensifie en vert (formation de \ce{Fe^2+} supplémentaire). Dépôt rouge sur l'électrode Cu.

\textbf{6. (Extension) Si on remplace le pont salin papier par un tube en U avec gel d'agar-agar \ce{KNO3}, que change-t-il ?}
\tcblower
Meilleure stabilité mécanique, moins de diffusion rapide des électrolytes dans le pont (meilleure séparation des demi-piles), résistance ohmique souvent plus faible et plus constante $\to$ mesures plus reproductibles, f.é.m. légèrement plus proche de la valeur théorique (moins de chute de tension interne).
\end{exoresolu}

\courssub{Conclusion}
\begin{aretenir}[Bilan du TP Piles]
À l'issue de ce TP, tu as acquis les compétences expérimentales et théoriques suivantes :
\begin{enumerate}
    \item \textbf{Savoir-faire technique :} Monter une pile électrochimique fonctionnelle avec du matériel de laboratoire simple (béchers, lames métalliques, pont salin papier filtre), brancher correctement un voltmètre (polarité COM / V) et relever une f.é.m.
    \item \textbf{Analyse physique :} Identifier l'anode (oxydation, pôle -) et la cathode (réduction, pôle +) grâce au signe de la tension mesurée. Déduire le sens du courant conventionnel dans le circuit externe (du + vers le -).
    \item \textbf{Modélisation chimique :} Écrire les demi-équations aux électrodes et l'équation-bilan globale pour chaque pile testée (Daniell, Zn/Fe, Fe/Cu, Cu/Ag).
    \item \textbf{Exploitation quantitative :} Calculer la f.é.m. théorique standard $E^\circ_{\text{pile}} = E^\circ_{\text{cathode}} - E^\circ_{\text{anode}}$ à partir de la classification des couples (\cle{potentiels standards}).
    \item \textbf{Esprit critique (Savoir-être) :} Comparer théorie et expérience. Comprendre que les écarts proviennent des conditions non standard (concentration \SI{0,5}{M}, température ambiante, potentiel de jonction liquide, résistance interne du pont salin) et surtout de la \textbf{nature réelle des électrodes} (impuretés, passivation, oxydes de surface) -- cas frappant de la pile Zn/Fe avec du fer commercial.
    \item \textbf{Sécurité \& Environnement :} Appliquer les consignes (lunettes, gants, récupération déchets métaux lourds \ce{Cu^2+}, \ce{Zn^2+}, \ce{Ag+}) indispensables en chimie pratique.
\end{enumerate}
Ce TP ancre la notion abstraite de \cle{potentiel d'oxydoréduction} dans une réalité mesurable : la tension aux bornes d'une pile est la manifestation macroscopique de la différence d'affinité électronique entre deux couples redox. C'est le principe de toute batterie (pile alcaline, accumulateur plomb, Li-ion) et de la corrosion métallique.
\end{aretenir}

\newpage

\coursec{Méthodes et exercices résolus}

\begin{propriete}[Propriétés physiques des alcools usuels]
Le tableau ci-dessous récapitule les propriétés physiques des premiers alcools aliphatiques, pertinents pour comprendre la fabrication artisanale du \emph{vin de palme} et du \emph{bili-bili} au Cameroun.
\begin{center}
\begin{tabularx}{\linewidth}{|>{\bfseries}l|X|X|X|}
\hline
\rowcolor{popblueL} \textbf{Alcool} & \textbf{Formule semi-développée} & \textbf{Point d'ébullition (\SI{}{\celsius})} & \textbf{Miscibilité dans l'eau} \\
\hline
Méthanol & \ce{CH3OH} & 65 & Totale \\
\hline
Éthanol & \ce{CH3CH2OH} & 78 & Totale \\
\hline
Propan-1-ol & \ce{CH3CH2CH2OH} & 97 & Totale \\
\hline
Butan-1-ol & \ce{CH3CH2CH2CH2OH} & 118 & Limitée (\SI{73}{g.L^{-1}}) \\
\hline
\end{tabularx}
\end{center}
\legende{Évolution du point d'ébullition et de la miscibilité avec la chaîne carbonée.}
\end{propriete}

\coursec{Les piles électrochimiques : fonctionnement, caractérisation et aspects quantitatifs}

\courssub{Rappels et définitions fondamentales}

\begin{definition}[Potentiel standard d'oxydoréduction]
Le \cle{potentiel standard d'oxydoréduction} $E^{\circ}$ d'un couple $\ce{Ox}/\ce{Red}$ est la tension (en volts) mesurée à \SI{25}{\celsius} pour une pile où ce couple est associé à l'\textbf{électrode normale à hydrogène} (ENH), toutes les espèces étant à l'état standard (solutions \SI{1}{mol.L^{-1}}, gaz à \SI{1}{bar}). Il traduit la force relative de l'oxydant : plus $E^{\circ}$ est \textbf{élevé}, plus l'oxydant est \textbf{fort} ; plus $E^{\circ}$ est \textbf{bas} (négatif), plus le réducteur est \textbf{fort}.
\end{definition}

\begin{definition}[Électrode normale à hydrogène (ENH)]
L'\cle{ENH} est la demi-pile de référence, de potentiel nul par convention : $E^{\circ}(\ce{H+}/\ce{H2}) = \SI{0,00}{V}$. Elle est constituée d'un électrode de platine inerte baignée par une solution \SI{1}{mol.L^{-1}} en $\ce{H+}$ et parcourue par du dihydrogène gazeux à \SI{1}{bar} : $\ce{Pt} | \ce{H2}(g, \SI{1}{bar}) | \ce{H+}(\SI{1}{mol.L^{-1}})$. Le couple est $\ce{H+}/\ce{H2}$.
\end{definition}

\begin{propriete}[Classification quantitative des couples]
On range les couples sur un \textbf{axe vertical des potentiels} croissants vers le haut.
\begin{itemize}
\item En \textbf{haut} : les oxydants les plus forts (ex. $\ce{Ag+}$, $\ce{Cu^2+}$, $\ce{O2}$).
\item En \textbf{bas} : les réducteurs les plus forts (ex. $\ce{Zn}$, $\ce{Fe}$, $\ce{Al}$, $\ce{Mg}$).
\item La \cle{règle du gamma ($\gamma$)} : une réaction spontanée met en jeu l'\textbf{oxydant du couple le plus haut} et le \textbf{réducteur du couple le plus bas}. Les deux espèces forment un $\gamma$ sur l'axe.
\end{itemize}
\end{propriete}

\courssub{Méthode 1 : Prévoir la réaction spontanée et la polarité d'une pile}

\begin{methode}[Prévoir le sens de la réaction et trouver les pôles]
\begin{enumerate}
\item \textbf{Identifier les deux couples} mis en jeu et repérer leurs potentiels standards $E^{\circ}$ (donnés dans l'énoncé ou dans la classification).
\item \textbf{Comparer les potentiels} : le couple dont le potentiel est le \cle{plus élevé} subit une \cle{réduction} (son oxydant capte les électrons) ; c'est le pôle \cle{positif} ($+$), appelé \cle{cathode}. Le couple de potentiel le plus faible subit une \cle{oxydation} ; c'est le pôle \cle{négatif} ($-$), appelé \cle{anode}.
\item \textbf{Écrire les demi-équations} aux électrodes : réduction à la cathode, oxydation à l'anode, en équilibrant les atomes et les charges (ajout d'$e^-$).
\item \textbf{Combiner} les deux demi-équations (en multipliant si besoin pour égaliser les électrons échangés) pour obtenir l'équation-bilan de fonctionnement : les électrons ne doivent \emph{jamais} apparaître dans le bilan.
\item \textbf{Vérifier} : la réaction spontanée est toujours celle entre l'oxydant du couple de plus haut potentiel et le réducteur du couple de plus bas potentiel (règle du \og gamma \fg{}).
\end{enumerate}
\textbf{Réflexe} : dans une pile, les électrons circulent dans le circuit extérieur du pôle $-$ vers le pôle $+$ ; le courant conventionnel circule en sens inverse.
\end{methode}

\begin{exoresolu}[Pile cuivre--argent]
On réalise une pile en associant une demi-pile $\ce{Ag+}/\ce{Ag}$ et une demi-pile $\ce{Cu^2+}/\ce{Cu}$, reliées par un pont salin. On donne : $E^{\circ}(\ce{Ag+}/\ce{Ag}) = \SI{0,80}{V}$ et $E^{\circ}(\ce{Cu^2+}/\ce{Cu}) = \SI{0,34}{V}$.
\begin{enumerate}
\item Indiquer la polarité de chaque électrode et le nom de chaque électrode (anode ou cathode).
\item Écrire les équations aux électrodes et l'équation-bilan de fonctionnement.
\item Préciser le sens de circulation des électrons dans le circuit extérieur.
\end{enumerate}
\tcblower
\textbf{1. Polarité.} On compare les potentiels standards :
\[ E^{\circ}(\ce{Ag+}/\ce{Ag}) = \SI{0,80}{V} > E^{\circ}(\ce{Cu^2+}/\ce{Cu}) = \SI{0,34}{V} \]
Le couple $\ce{Ag+}/\ce{Ag}$ a le potentiel le plus élevé : $\ce{Ag+}$ est l'oxydant le plus fort, il se \textbf{réduit}. L'électrode d'argent est donc le \textbf{pôle positif} ($+$), c'est la \textbf{cathode}.\\
Le couple $\ce{Cu^2+}/\ce{Cu}$ a le potentiel le plus faible : le cuivre métallique s'\textbf{oxyde}. L'électrode de cuivre est le \textbf{pôle négatif} ($-$), c'est l'\textbf{anode}.

\textbf{2. Équations.}
\begin{itemize}
\item À la cathode ($+$, argent), réduction : $\ce{Ag+ + e- -> Ag}$
\item À l'anode ($-$, cuivre), oxydation : $\ce{Cu -> Cu^2+ + 2e-}$
\end{itemize}
La demi-équation de l'argent met en jeu 1 électron, celle du cuivre 2 électrons : on multiplie la première par 2 avant d'additionner :
\[ \ce{2Ag+ + Cu -> 2Ag + Cu^2+} \]
C'est l'équation-bilan de fonctionnement de la pile : le cuivre se dissout, l'argent se dépose.

\textbf{3. Sens des électrons.} Les électrons sont libérés à l'anode de cuivre et captés à la cathode d'argent : ils circulent dans le circuit extérieur \textbf{de l'électrode de cuivre (pôle $-$) vers l'électrode d'argent (pôle $+$)}.

\textbf{Conclusion} : pôle $+$ = électrode d'argent (cathode), pôle $-$ = électrode de cuivre (anode) ; bilan : $\ce{2Ag+ + Cu -> 2Ag + Cu^2+}$.
\end{exoresolu}

\courssub{Méthode 2 : Calculer la f.é.m. d'une pile}

\begin{methode}[Détermination de la force électromotrice]
\begin{enumerate}
\item Repérer les potentiels standards $E^{\circ}$ des deux couples.
\item Identifier le pôle $+$ (couple de potentiel le plus élevé) et le pôle $-$ (couple de potentiel le plus faible).
\item Appliquer la relation, dans les conditions standard :
\[ E = E^{\circ}_{(+)} - E^{\circ}_{(-)} \]
où $E^{\circ}_{(+)}$ est le potentiel du couple du pôle positif et $E^{\circ}_{(-)}$ celui du pôle négatif.
\item La f.é.m. est \textbf{toujours positive}. Si le calcul donne une valeur négative, c'est que les pôles ont été inversés : reprendre l'étape 2.
\item Exprimer le résultat en volts (V) avec le bon nombre de chiffres significatifs (en général 2 ou 3).
\end{enumerate}
\textbf{Réflexe} : la f.é.m. ne dépend \emph{pas} des coefficients stœchiométriques : on ne multiplie jamais un potentiel par un coefficient !
\end{methode}

\begin{exoresolu}[F.é.m. de la pile Daniell]
On considère la pile Daniell constituée des couples $\ce{Zn^2+}/\ce{Zn}$ et $\ce{Cu^2+}/\ce{Cu}$, avec des solutions de concentration \SI{1,0}{mol.L^{-1}}. On donne : $E^{\circ}(\ce{Zn^2+}/\ce{Zn}) = \SI{-0,76}{V}$ et $E^{\circ}(\ce{Cu^2+}/\ce{Cu}) = \SI{0,34}{V}$.
\begin{enumerate}
\item Déterminer la polarité de la pile.
\item Calculer sa f.é.m. standard.
\item Écrire l'équation-bilan de fonctionnement.
\end{enumerate}
\tcblower
\textbf{1. Polarité.} On compare : $E^{\circ}(\ce{Cu^2+}/\ce{Cu}) = \SI{0,34}{V} > E^{\circ}(\ce{Zn^2+}/\ce{Zn}) = \SI{-0,76}{V}$.\\
Le couple du cuivre a le potentiel le plus élevé : l'électrode de \textbf{cuivre est le pôle $+$} (cathode). L'électrode de \textbf{zinc est le pôle $-$} (anode).

\textbf{2. F.é.m. standard.}
\[ E = E^{\circ}_{(+)} - E^{\circ}_{(-)} = E^{\circ}(\ce{Cu^2+}/\ce{Cu}) - E^{\circ}(\ce{Zn^2+}/\ce{Zn}) \]
\[ E = 0,34 - (-0,76) = 0,34 + 0,76 = \SI{1,10}{V} \]
Attention au signe : soustraire un nombre négatif revient à additionner.

\textbf{3. Équation-bilan.}
\begin{itemize}
\item Cathode ($+$, cuivre), réduction : $\ce{Cu^2+ + 2e- -> Cu}$
\item Anode ($-$, zinc), oxydation : $\ce{Zn -> Zn^2+ + 2e-}$
\end{itemize}
Les deux demi-équations échangent 2 électrons : on additionne directement :
\[ \ce{Cu^2+ + Zn -> Cu + Zn^2+} \]

\textbf{Conclusion} : la pile Daniell a pour f.é.m. standard $E = \SI{1,10}{V}$, le pôle positif étant l'électrode de cuivre. C'est bien la valeur mesurée expérimentalement au voltmètre.
\end{exoresolu}

\courssub{Méthode 3 : Utiliser la notation conventionnelle d'une pile}

\begin{methode}[Écrire et interpréter le schéma conventionnel]
\begin{enumerate}
\item La notation conventionnelle s'écrit toujours en plaçant le \textbf{pôle $-$ à gauche} et le \textbf{pôle $+$ à droite} :
\[ \ominus \; \text{réducteur}_1 \; | \; \text{oxydant}_1 \; || \; \text{oxydant}_2 \; | \; \text{réducteur}_2 \; \oplus \]
\item Le trait simple $|$ symbolise la \textbf{frontière} entre le métal (électrode) et la solution ionique.
\item Le double trait $||$ symbolise le \textbf{pont salin} (ou la paroi poreuse) reliant les deux demi-piles.
\item On peut préciser les concentrations entre parenthèses, ex. $\ce{Zn^2+}(\SI{1,0}{mol.L^{-1}})$.
\item Pour retrouver la polarité à partir du schéma : à gauche, l'électrode subit l'\textbf{oxydation} (anode, pôle $-$) ; à droite, la \textbf{réduction} (cathode, pôle $+$).
\end{enumerate}
\textbf{Réflexe} : l'ordre d'écriture dans chaque demi-pile suit le sens parcouru depuis l'électrode : métal $|$ ions en solution.
\end{methode}

\begin{exoresolu}[Notation de la pile fer--plomb]
On réalise une pile en plongeant une lame de fer dans une solution de $\ce{Fe^2+}$ (\SI{0,10}{mol.L^{-1}}) et une lame de plomb dans une solution de $\ce{Pb^2+}$ (\SI{0,10}{mol.L^{-1}}). On donne : $E^{\circ}(\ce{Fe^2+}/\ce{Fe}) = \SI{-0,44}{V}$ et $E^{\circ}(\ce{Pb^2+}/\ce{Pb}) = \SI{-0,13}{V}$.
\begin{enumerate}
\item Donner la notation conventionnelle de cette pile.
\item Calculer sa f.é.m.
\item Écrire l'équation-bilan de fonctionnement.
\end{enumerate}
\tcblower
\textbf{1. Notation conventionnelle.} Comparaison des potentiels : $-0,13 > -0,44$, donc le couple $\ce{Pb^2+}/\ce{Pb}$ a le potentiel le plus élevé : le plomb est le pôle $+$ (cathode) ; le fer est le pôle $-$ (anode).\\
On place le pôle $-$ à gauche :
\[ \ominus \; \ce{Fe} \; | \; \ce{Fe^2+}(\SI{0,10}{mol.L^{-1}}) \; || \; \ce{Pb^2+}(\SI{0,10}{mol.L^{-1}}) \; | \; \ce{Pb} \; \oplus \]

\textbf{2. F.é.m.}
\[ E = E^{\circ}_{(+)} - E^{\circ}_{(-)} = (-0,13) - (-0,44) = \SI{0,31}{V} \]

\textbf{3. Équation-bilan.}
\begin{itemize}
\item Cathode ($+$, plomb) : $\ce{Pb^2+ + 2e- -> Pb}$
\item Anode ($-$, fer) : $\ce{Fe -> Fe^2+ + 2e-}$
\end{itemize}
\[ \ce{Pb^2+ + Fe -> Pb + Fe^2+} \]

\textbf{Conclusion} : la pile s'écrit $\ominus \; \ce{Fe} | \ce{Fe^2+} || \ce{Pb^2+} | \ce{Pb} \; \oplus$ et sa f.é.m. vaut \SI{0,31}{V}. La lame de fer s'amincit pendant le fonctionnement tandis que du plomb se dépose.
\end{exoresolu}

\courssub{Méthode 4 : Déterminer expérimentalement les pôles et la f.é.m. d'une pile}

\begin{methode}[Mesure au voltmètre]
\begin{enumerate}
\item \textbf{Réaliser la pile} : deux béchers contenant chacun la solution ionique et la lame métallique correspondante, reliés par un pont salin (papier-filtre imbibé de solution de $\ce{KNO3}$ ou de $\ce{KCl}$, ou tube en U).
\item \textbf{Brancher un voltmètre} (ou un multimètre en position voltmètre, calibre continu) aux bornes de la pile : borne \og V \fg{} sur une électrode, borne \og COM \fg{} sur l'autre.
\item \textbf{Lire la tension} :
\begin{itemize}
\item si l'affichage est \textbf{positif}, l'électrode reliée à la borne \og V \fg{} est le pôle $+$ ;
\item si l'affichage est \textbf{négatif}, les bornes sont inversées : l'électrode reliée à \og V \fg{} est le pôle $-$.
\end{itemize}
\item La valeur absolue lue est la \textbf{f.é.m.} $E$ de la pile (la pile débitant un courant quasi nul dans le voltmètre, la tension mesurée est bien la f.é.m.).
\item \textbf{Identifier les pôles sans appareil} : on peut aussi repérer les transformations visibles (dépôt métallique à la cathode, amincissement de l'anode) ou utiliser un électrolyte révélateur.
\end{enumerate}
\textbf{Réflexe} : le pont salin assure la neutralité électrique des solutions ; sans lui, la pile ne débite pas. Ses ions migrent : les cations vers la cathode ($+$), les anions vers l'anode ($-$).
\end{methode}

\begin{exoresolu}[Étude expérimentale d'une pile zinc--fer]
Un élève du lycée de Ngaoundéré réalise une pile en associant les demi-piles $\ce{Zn^2+}/\ce{Zn}$ et $\ce{Fe^2+}/\ce{Fe}$. Il branche la borne \og V \fg{} du voltmètre sur l'électrode de fer et la borne \og COM \fg{} sur l'électrode de zinc : l'appareil affiche $U = \SI{+0,32}{V}$. On donne $E^{\circ}(\ce{Zn^2+}/\ce{Zn}) = \SI{-0,76}{V}$.
\begin{enumerate}
\item Déterminer la polarité de la pile.
\item En déduire le potentiel standard du couple $\ce{Fe^2+}/\ce{Fe}$.
\item Indiquer le sens de migration des ions $\ce{K+}$ et $\ce{NO3-}$ du pont salin.
\end{enumerate}
\tcblower
\textbf{1. Polarité.} L'affichage est \textbf{positif} alors que la borne \og V \fg{} est reliée au fer : l'électrode de \textbf{fer est le pôle $+$} (cathode) et l'électrode de \textbf{zinc est le pôle $-$} (anode). La tension lue est la f.é.m. : $E = \SI{0,32}{V}$.

\textbf{2. Potentiel standard du couple $\ce{Fe^2+}/\ce{Fe}$.} On applique la relation de la f.é.m. :
\[ E = E^{\circ}_{(+)} - E^{\circ}_{(-)} = E^{\circ}(\ce{Fe^2+}/\ce{Fe}) - E^{\circ}(\ce{Zn^2+}/\ce{Zn}) \]
\[ E^{\circ}(\ce{Fe^2+}/\ce{Fe}) = E + E^{\circ}(\ce{Zn^2+}/\ce{Zn}) = 0,32 + (-0,76) = \SI{-0,44}{V} \]
On retrouve la valeur tabulée du couple $\ce{Fe^2+}/\ce{Fe}$, ce qui valide la mesure.

\textbf{3. Migration des ions du pont salin.} Dans le pont salin, les ions migrent de façon à compenser les charges qui apparaissent dans chaque compartiment :
\begin{itemize}
\item les cations $\ce{K+}$ migrent vers le compartiment de la \textbf{cathode} (électrode de fer, pôle $+$), où les ions $\ce{Fe^2+}$ disparaissent par réduction ;
\item les anions $\ce{NO3-}$ migrent vers le compartiment de l'\textbf{anode} (électrode de zinc, pôle $-$), où des ions $\ce{Zn^2+}$ apparaissent par oxydation.
\end{itemize}

\textbf{Conclusion} : pôle $+$ = fer, pôle $-$ = zinc, $E = \SI{0,32}{V}$ et $E^{\circ}(\ce{Fe^2+}/\ce{Fe}) = \SI{-0,44}{V}$.
\end{exoresolu}

\courssub{Méthode 5 : Étude quantitative du fonctionnement d'une pile}

\begin{methode}[Quantité d'électricité, avancement et variation de masse]
\begin{enumerate}
\item Écrire l'équation-bilan de fonctionnement et dresser un \textbf{tableau d'avancement} si les quantités initiales sont connues.
\item La \textbf{quantité d'électricité} débitée se calcule par :
\[ Q = I \times \Delta t \quad \text{avec } Q \text{ en coulombs (C), } I \text{ en ampères (A), } \Delta t \text{ en secondes (s)} \]
\item La quantité d'électrons échangés est : $n(e^-) = \dfrac{Q}{F}$ avec $F = \SI{96500}{C.mol^{-1}}$ (constante de Faraday).
\item Relier $n(e^-)$ à l'avancement $x$ grâce aux coefficients des demi-équations : si la demi-équation échange $n$ électrons, alors $n(e^-) = n \times x$.
\item En déduire les \textbf{variations de masse} des électrodes : $\Delta m = x \times M$ (perte à l'anode, gain à la cathode).
\end{enumerate}
\textbf{Réflexe} : une pile s'use quand l'un des réactifs est épuisé ; en général, c'est le métal de l'anode qui se consomme.
\end{methode}

\begin{exoresolu}[Pile Daniell en fonctionnement]
Une pile Daniell $\ominus \; \ce{Zn} | \ce{Zn^2+} || \ce{Cu^2+} | \ce{Cu} \; \oplus$ débite un courant d'intensité constante $I = \SI{50}{mA}$ pendant $\Delta t = \SI{2}{h}$. On donne : $M(\ce{Zn}) = \SI{65,4}{g.mol^{-1}}$, $M(\ce{Cu}) = \SI{63,5}{g.mol^{-1}}$, $F = \SI{96500}{C.mol^{-1}}$.
\begin{enumerate}
\item Calculer la quantité d'électricité débitée.
\item En déduire la variation de masse de chaque électrode.
\end{enumerate}
\tcblower
\textbf{1. Quantité d'électricité.} Conversion des données : $I = \SI{50}{mA} = \SI{0,050}{A}$ et $\Delta t = \SI{2}{h} = \SI{7200}{s}$.
\[ Q = I \times \Delta t = 0,050 \times 7200 = \SI{360}{C} \]

\textbf{2. Variations de masse.} Les demi-équations sont :
\[ \ce{Zn -> Zn^2+ + 2e-} \qquad \ce{Cu^2+ + 2e- -> Cu} \]
Chaque demi-équation échange 2 électrons par atome de métal. La quantité d'électrons échangés :
\[ n(e^-) = \frac{Q}{F} = \frac{360}{96500} = \SI{3,73e-3}{mol} \]
L'avancement de la réaction : $x = \dfrac{n(e^-)}{2} = \dfrac{3,73 \times 10^{-3}}{2} = \SI{1,87e-3}{mol}$.

\begin{center}
\begin{tabularx}{\linewidth}{|l|X|X|X|}
\hline
\rowcolor{popblueL} Espèce & $\ce{Zn}$ & $\ce{Cu^2+}$ & $n(e^-)$ \\
\hline
Quantités échangées & $-x$ & $-x$ & $2x$ \\
\hline
Valeur (mol) & $-1,87\times 10^{-3}$ & $-1,87\times 10^{-3}$ & $3,73\times 10^{-3}$ \\
\hline
\end{tabularx}
\end{center}

\begin{itemize}
\item \textbf{Anode de zinc} (elle se consomme) :
\[ \Delta m(\ce{Zn}) = -x \times M(\ce{Zn}) = -1,87 \times 10^{-3} \times 65,4 = \SI{-0,122}{g} \]
La lame de zinc a perdu environ \SI{0,12}{g}.
\item \textbf{Cathode de cuivre} (du cuivre se dépose) :
\[ \Delta m(\ce{Cu}) = +x \times M(\ce{Cu}) = 1,87 \times 10^{-3} \times 63,5 = \SI{+0,119}{g} \]
L'électrode de cuivre a gagné environ \SI{0,12}{g}.
\end{itemize}

\textbf{Conclusion} : en 2 heures, la pile a débité \SI{360}{C} ; l'anode de zinc s'est amincie de \SI{0,122}{g} et la cathode de cuivre s'est épaissie de \SI{0,119}{g}. C'est ce principe qui explique l'usure des piles usuelles.
\end{exoresolu}

\courssub{Le tableau d'avancement et la règle du gamma (Synthèse visuelle)}

\begin{popfigure}
\begin{center}
\begin{tikzpicture}[
    scale=0.9,
    every node/.style={font=\small},
    table cell/.style={draw=popline, minimum height=1.1cm, minimum width=2.2cm, align=center, inner sep=2pt},
    header cell/.style={table cell, fill=popblueL, font=\bfseries, text=popdark},
    title cell/.style={table cell, fill=popgreenL, font=\bfseries\itshape, text=popdark},
    reactant/.style={table cell, fill=poporange!15},
    product/.style={table cell, fill=popgreen!15},
    arrow style/.style={->, thick, poporange, shorten >=2pt, shorten <=2pt},
    gamma label/.style={midway, above, font=\small\bfseries, text=poporange}
]
% Coordonnées
\def\colA{0}   % Espèce
\def\colB{2.4} % État initial
\def\colC{4.8} % Variation
\def\colD{7.2} % État final
\def\rowH{1.2} % Hauteur ligne
\def\rowTitle{0}
\def\rowInit{-\rowH}
\def\rowVar{-2*\rowH}
\def\rowFinal{-3*\rowH}
\def\rowGamma{-4*\rowH}

% Titre du tableau
\node[title cell, minimum width=9.6cm] at (4.8, \rowH*0.5) {Tableau d'avancement pour \ce{Zn + Cu^2+ -> Zn^2+ + Cu}};

% En-têtes
\node[header cell] at (\colA, \rowTitle) {Espèce chimique};
\node[header cell] at (\colB, \rowTitle) {État initial\\$n_i^0$ (mol)};
\node[header cell] at (\colC, \rowTitle) {Variation\\$\Delta n_i$ (mol)};
\node[header cell] at (\colD, \rowTitle) {État final\\$n_i^f$ (mol)};

% Lignes Réactifs
\node[reactant] at (\colA, \rowInit) {\ce{Zn(s)}};
\node[reactant] at (\colB, \rowInit) {$n_{\ce{Zn}}^0$};
\node[reactant] at (\colC, \rowInit) {$-1 \times x$};
\node[reactant] at (\colD, \rowInit) {$n_{\ce{Zn}}^0 - x$};

\node[reactant] at (\colA, \rowInit-\rowH) {\ce{Cu^2+(aq)}};
\node[reactant] at (\colB, \rowInit-\rowH) {$n_{\ce{Cu^2+}}^0$};
\node[reactant] at (\colC, \rowInit-\rowH) {$-1 \times x$};
\node[reactant] at (\colD, \rowInit-\rowH) {$n_{\ce{Cu^2+}}^0 - x$};

% Lignes Produits
\node[product] at (\colA, \rowInit-2*\rowH) {\ce{Zn^2+(aq)}};
\node[product] at (\colB, \rowInit-2*\rowH) {$n_{\ce{Zn^2+}}^0$};
\node[product] at (\colC, \rowInit-2*\rowH) {$+1 \times x$};
\node[product] at (\colD, \rowInit-2*\rowH) {$n_{\ce{Zn^2+}}^0 + x$};

\node[product] at (\colA, \rowInit-3*\rowH) {\ce{Cu(s)}};
\node[product] at (\colB, \rowInit-3*\rowH) {$n_{\ce{Cu}}^0$};
\node[product] at (\colC, \rowInit-3*\rowH) {$+1 \times x$};
\node[product] at (\colD, \rowInit-3*\rowH) {$n_{\ce{Cu}}^0 + x$};

% Flèches et annotations Gamma (règle du gamma)
\draw[arrow style] (\colA-1.3, \rowInit) -- node[gamma label, left] {$\gamma_{\ce{Zn}} = 1$} (\colA-1.3, \rowInit-3*\rowH);
\draw[arrow style] (\colC+1.3, \rowInit) -- node[gamma label, right] {$\gamma_i \times x$} (\colC+1.3, \rowInit-3*\rowH);

% Légende Gamma
\node[align=center, text=popdark, font=\small\bfseries] at (4.8, \rowInit-4*\rowH) {
    \textbf{Règle du gamma :} Variation $= \gamma_i \times x$ \quad ($\gamma_i > 0$ pour produits, $< 0$ pour réactifs)
};
\end{tikzpicture}
\end{center}
\legende{Illustration de la règle du gamma dans un tableau d'avancement pour la pile Daniell. Les coefficients stœchiométriques ($\gamma_i$) lient l'avancement $x$ aux variations de quantités de matière.}
\end{popfigure}

\courssub{Méthode 6 : Classer des couples et prévoir une réaction spontanée}

\begin{methode}[Utiliser la classification quantitative (règle du gamma)]
\begin{enumerate}
\item Placer les couples sur un \textbf{axe vertical des potentiels} croissants vers le haut : plus $E^{\circ}$ est élevé, plus l'oxydant est fort ; plus $E^{\circ}$ est bas, plus le réducteur est fort.
\item Une réaction \textbf{spontanée} a lieu entre l'\textbf{oxydant du couple de plus haut potentiel} et le \textbf{réducteur du couple de plus bas potentiel} : les deux espèces forment un \og gamma \fg{} ($\gamma$) sur l'axe.
\item Écrire les demi-équations dans le bon sens (réduction en haut, oxydation en bas), équilibrer les électrons, puis combiner.
\item Si les deux espèces mises en présence ne forment pas un gamma (oxydant faible face à réducteur faible), \textbf{aucune réaction} n'a lieu.
\end{enumerate}
\textbf{Réflexe} : l'électrode normale à hydrogène (E.N.H.) sert de référence : $E^{\circ}(\ce{H+}/\ce{H2}) = \SI{0,00}{V}$ par convention. Un couple de potentiel négatif a un réducteur plus fort que $\ce{H2}$ : son métal est attaqué par les acides.
\end{methode}

\begin{exoresolu}[Réactions prévues grâce à la classification]
On donne les potentiels standards : $E^{\circ}(\ce{Zn^2+}/\ce{Zn}) = \SI{-0,76}{V}$ ; $E^{\circ}(\ce{H+}/\ce{H2}) = \SI{0,00}{V}$ ; $E^{\circ}(\ce{Cu^2+}/\ce{Cu}) = \SI{0,34}{V}$.
\begin{enumerate}
\item Classer les trois couples sur un axe de potentiels.
\item Prévoir si une solution d'acide chlorhydrique attaque le zinc ; le cuivre. Écrire l'équation-bilan quand la réaction a lieu.
\item Un bijoutier veut nettoyer un objet en cuivre avec de l'acide dilué : est-ce efficace pour dissoudre le cuivre ?
\end{enumerate}
\tcblower
\textbf{1. Classification.} On place les couples sur un axe vertical, potentiels croissants vers le haut :

\begin{popfigure}
\begin{tikzpicture}
\draw[->, thick, popdark] (0,0) -- (0,6.2) node[above] {$E^{\circ}$ (V)};
\foreach \y/\val in {0.8/{-0,76}, 3.4/{0,00}, 5.2/{0,34}} {
  \draw[popdark] (-0.1,\y) -- (0.1,\y) node[left=2pt, xshift=-2pt] {\val};
}
\node[right, popblue] at (0.15,5.2) {\ce{Cu^2+}/\ce{Cu}};
\node[right, popgreen] at (0.15,3.4) {\ce{H+}/\ce{H2} (référence)};
\node[right, poporange] at (0.15,0.8) {\ce{Zn^2+}/\ce{Zn}};
\draw[->, very thick, poppink] (2.9,5.2) .. controls (4.6,4.4) and (4.6,1.6) .. (2.9,0.8);
\node[poppink] at (5.4,3) {$\gamma$};
\end{tikzpicture}
\legende{Classification des couples et règle du gamma : $\ce{Cu^2+}$ (oxydant le plus fort) réagit avec $\ce{Zn}$ (réducteur le plus fort).}
\end{popfigure}

\textbf{2. Attaque par l'acide chlorhydrique.} L'acide apporte l'oxydant $\ce{H+}$ (couple $\ce{H+}/\ce{H2}$, $E^{\circ} = \SI{0,00}{V}$).
\begin{itemize}
\item \textbf{Zinc} : $E^{\circ}(\ce{H+}/\ce{H2}) > E^{\circ}(\ce{Zn^2+}/\ce{Zn})$, donc $\ce{H+}$ (oxydant du couple le plus haut) et $\ce{Zn}$ (réducteur du couple le plus bas) forment un gamma : la réaction \textbf{a lieu}.
\[ \ce{2H+ + 2e- -> H2} \qquad \ce{Zn -> Zn^2+ + 2e-} \]
\[ \ce{2H+ + Zn -> H2 + Zn^2+} \]
On observe un dégagement de dihydrogène : le zinc est attaqué.
\item \textbf{Cuivre} : $E^{\circ}(\ce{H+}/\ce{H2}) < E^{\circ}(\ce{Cu^2+}/\ce{Cu})$, donc $\ce{H+}$ et $\ce{Cu}$ ne forment \textbf{pas} un gamma : \textbf{aucune réaction} n'a lieu.
\end{itemize}

\textbf{3. Nettoyage du cuivre.} L'acide dilué ne peut pas dissoudre le cuivre métallique (pas de réaction). Il ne sert qu'à éliminer les dépôts calcaires ou les oxydes de surface, pas le métal lui-même.

\textbf{Conclusion} : seuls les métaux des couples de potentiel négatif (comme $\ce{Zn}$, $\ce{Fe}$, $\ce{Al}$) sont attaqués par les acides à ions $\ce{H+}$ avec dégagement de $\ce{H2}$ ; les métaux \og nobles \fg{} ($\ce{Cu}$, $\ce{Ag}$, $\ce{Au}$) résistent.
\end{exoresolu}

\begin{savaistu}[Piles au Cameroun : de la pile saline à la voiture électrique]
Au Cameroun, les piles alcalines (type LR6, LR03) alimentent lampes torches, radios et télécommandes dans les zones non électrifiées. Leur couple est $\ce{Zn}/\ce{MnO2}$ (f.é.m. $\approx \SI{1,5}{V}$). Les batteries plomb-acide (accumulateurs) des véhicules (taxis, cars, camions SONARA/CIMENCAM) utilisent le couple $\ce{Pb}/\ce{PbO2}$ dans l'acide sulfurique ($E \approx \SI{2,0}{V}$ par élément, \SI{12}{V} pour 6 éléments). Les nouvelles motos électriques et projets solaires intègrent des batteries Lithium-ion ($\ce{LiCoO2}/\ce{Li}$, \SI{3,7}{V}), plus légères et à haute densité d'énergie. Le recyclage de ces piles usagées (plomb, lithium, zinc) est un enjeu environnemental majeur pour nos villes (Douala, Yaoundé, Bafoussam).
\end{savaistu}

\begin{attention}[Les 5 erreurs qui coûtent des points au Probatoire]
\begin{enumerate}
\item \textbf{Inverser les pôles} : le pôle $+$ est le couple de potentiel le \emph{plus élevé} (cathode, réduction) ; le pôle $-$ est celui du potentiel le plus faible (anode, oxydation). Ne jamais confondre anode/cathode avec $+$/$-$ : dans une pile, anode = pôle $-$.
\item \textbf{Multiplier les potentiels par les coefficients stœchiométriques} : même si l'équation-bilan est $\ce{2Ag+ + Cu -> 2Ag + Cu^2+}$, on calcule toujours $E = 0,80 - 0,34 = \SI{0,46}{V}$, jamais $2 \times 0,80 - 0,34$.
\item \textbf{Faire apparaître les électrons dans l'équation-bilan} : les $e^-$ doivent s'éliminer en combinant les demi-équations. Une équation-bilan contenant des électrons est fausse.
\item \textbf{Oublier les conversions} dans les calculs quantitatifs : le temps en \emph{secondes} ($\SI{2}{h} = \SI{7200}{s}$), l'intensité en \emph{ampères} ($\SI{50}{mA} = \SI{0,050}{A}$), avant d'appliquer $Q = I\,\Delta t$ puis $n(e^-) = Q/F$.
\item \textbf{Mal gérer les signes dans la f.é.m.} : $E = E^{\circ}_{(+)} - E^{\circ}_{(-)}$ avec des potentiels souvent négatifs ; écrire $E = 0,34 - (-0,76) = \SI{1,10}{V}$ et non $\SI{-0,42}{V}$. Une f.é.m. est toujours \emph{positive} : un résultat négatif signale une inversion des pôles.
\end{enumerate}
\end{attention}

\newpage

\coursec{Exercices}

\courssub{Je vérifie mes connaissances}

\begin{exercice}[QCM : la pile Daniell]
Pour chaque affirmation, choisis la (ou les) bonne(s) réponse(s), puis corrige les propositions fausses.

Dans une pile Daniell en fonctionnement :
\begin{enumerate}
\item[\textbf{a)}] Les électrons circulent dans le pont salin, de l'électrode de cuivre vers l'électrode de zinc.
\item[\textbf{b)}] Les électrons circulent dans le circuit extérieur, de l'électrode de zinc vers l'électrode de cuivre.
\item[\textbf{c)}] L'électrode de zinc est le pôle négatif de la pile ; on l'appelle l'\emph{anode}.
\item[\textbf{d)}] À l'électrode de cuivre, il se produit une oxydation : \ce{Cu -> Cu^2+ + 2e-}.
\item[\textbf{e)}] Le pont salin permet le passage du courant dans les solutions et assure l'électroneutralité de chaque compartiment.
\item[\textbf{f)}] La f.é.m. standard de la pile Daniell vaut $E^{\circ} = E^{\circ}(\ce{Cu^2+}/\ce{Cu}) - E^{\circ}(\ce{Zn^2+}/\ce{Zn})$.
\end{enumerate}
\end{exercice}

\begin{exercice}[Vrai ou faux, avec justification]
Réponds par \textbf{vrai} ou \textbf{faux}, puis justifie brièvement chaque réponse.
\begin{enumerate}
\item Le potentiel standard de l'électrode normale à hydrogène vaut, par convention, $E^{\circ}(\ce{H+}/\ce{H2}) = \SI{0,00}{V}$ à toute température.
\item Plus le potentiel standard d'un couple est élevé, plus le réducteur du couple est fort.
\item Dans une pile, l'oxydation a toujours lieu à l'anode, qui constitue le pôle négatif.
\item La notation conventionnelle d'une pile s'écrit toujours avec le pôle négatif à gauche et le pôle positif à droite.
\item Une pile usuelle (pile saline, pile alcaline) ne contient pas de pont salin : elle ne peut donc pas fonctionner.
\end{enumerate}
\end{exercice}

\begin{exercice}[Texte à trous]
Recopie et complète le texte suivant avec les mots ou expressions de la liste : \emph{anode ; cathode ; oxydation ; réduction ; négatif ; positif ; pont salin ; f.é.m. ; hydrogène ; volt}.

La pile Daniell associe deux demi-piles : une lame de zinc plongée dans une solution de sulfate de zinc et une lame de cuivre plongée dans une solution de sulfate de cuivre. Les deux solutions sont reliées par un \ldots\ldots\ldots qui assure le passage du courant. L'électrode de zinc, siège d'une \ldots\ldots\ldots, est l'\ldots\ldots\ldots de la pile ; elle constitue son pôle \ldots\ldots\ldots. L'électrode de cuivre, siège d'une \ldots\ldots\ldots, est la \ldots\ldots\ldots ; elle constitue le pôle \ldots\ldots\ldots. La tension mesurée entre les deux électrodes lorsque la pile ne débite pas s'appelle la \ldots\ldots\ldots de la pile ; elle s'exprime en \ldots\ldots\ldots. Pour mesurer les potentiels d'oxydoréduction, on utilise une demi-pile de référence : l'électrode normale à \ldots\ldots\ldots.
\end{exercice}

\begin{exercice}[Définitions et conventions]
\begin{enumerate}
\item Définis : demi-pile ; potentiel standard d'oxydoréduction ; f.é.m. d'une pile.
\item Donne les conditions \emph{standard} (ou normales) dans lesquelles on définit $E^{\circ}$ d'un couple.
\item Écris la notation conventionnelle de la pile Daniell et précise la signification de chaque barre verticale et de la double barre.
\item Pourquoi dit-on que l'électrode normale à hydrogène est une électrode « de référence » ? Décris brièvement sa constitution.
\end{enumerate}
\end{exercice}

\courssub{Je m'entraîne}

\begin{exercice}[La pile fer--argent]
On réalise une pile en associant les couples \ce{Fe^2+}/\ce{Fe} et \ce{Ag+}/\ce{Ag}. On donne : $E^{\circ}(\ce{Fe^2+}/\ce{Fe}) = \SI{-0,44}{V}$ et $E^{\circ}(\ce{Ag+}/\ce{Ag}) = \SI{+0,80}{V}$.
\begin{enumerate}
\item Prévois la polarité de chaque électrode. Justifie.
\item Écris les équations des réactions aux électrodes, puis l'équation-bilan de fonctionnement de la pile.
\item Calcule la f.é.m. standard de cette pile.
\item Donne sa notation conventionnelle.
\item Indique le sens de circulation des électrons dans le circuit extérieur et le sens de déplacement des ions dans le pont salin (au nitrate de potassium).
\end{enumerate}
\end{exercice}

\begin{exercice}[Prévision de réactions à l'aide des potentiels]
On donne les potentiels standard : $E^{\circ}(\ce{Pb^2+}/\ce{Pb}) = \SI{-0,13}{V}$ ; $E^{\circ}(\ce{Cu^2+}/\ce{Cu}) = \SI{+0,34}{V}$ ; $E^{\circ}(\ce{Fe^2+}/\ce{Fe}) = \SI{-0,44}{V}$.
\begin{enumerate}
\item Classe les trois couples par potentiel croissant, puis les trois métaux par pouvoir réducteur croissant.
\item Une lame de plomb est plongée dans une solution de sulfate de cuivre(II). Une réaction se produit-elle ? Si oui, écris son équation-bilan et justifie par la règle du gamma.
\item Une lame de plomb est plongée dans une solution de sulfate de fer(II). Une réaction se produit-elle ? Justifie.
\item Le cuivre métallique peut-il réduire les ions \ce{Fe^2+} ? Justifie.
\end{enumerate}
\end{exercice}

\begin{exercice}[Schéma et mesures de f.é.m.]
Un élève de Première D réalise au laboratoire trois piles et mesure leur f.é.m. avec un voltmètre électronique. Il note :
\begin{itemize}
\item Pile 1 : \ce{Zn}/\ce{Zn^2+} $\parallel$ \ce{Cu^2+}/\ce{Cu} ; $E_1 = \SI{1,10}{V}$ ;
\item Pile 2 : \ce{Zn}/\ce{Zn^2+} $\parallel$ \ce{Ag+}/\ce{Ag} ; $E_2 = \SI{1,56}{V}$ ;
\item Pile 3 : \ce{Cu}/\ce{Cu^2+} $\parallel$ \ce{Ag+}/\ce{Ag} ; $E_3 = \SI{0,46}{V}$.
\end{itemize}
\begin{enumerate}
\item Schématise la pile 1 en précisant la nature des solutions, le pont salin, les pôles et le sens des électrons.
\item Sachant que $E^{\circ}(\ce{Zn^2+}/\ce{Zn}) = \SI{-0,76}{V}$, déduis de la pile 1 le potentiel standard du couple \ce{Cu^2+}/\ce{Cu}.
\item Déduis de la pile 2 le potentiel standard du couple \ce{Ag+}/\ce{Ag}.
\item Vérifie la cohérence des résultats à l'aide de la pile 3.
\end{enumerate}
\end{exercice}

\begin{exercice}[Construction d'une pile de f.é.m. imposée]
On dispose au laboratoire des métaux et solutions correspondant aux couples suivants : \ce{Mg^2+}/\ce{Mg} ($\SI{-2,37}{V}$), \ce{Zn^2+}/\ce{Zn} ($\SI{-0,76}{V}$), \ce{Pb^2+}/\ce{Pb} ($\SI{-0,13}{V}$), \ce{Cu^2+}/\ce{Cu} ($\SI{+0,34}{V}$).
\begin{enumerate}
\item Quelle association de couples donne la pile de plus grande f.é.m. ? Calcule sa valeur et donne la notation conventionnelle de cette pile.
\item Quelle association de couples donne une f.é.m. standard voisine de $\SI{0,47}{V}$ ? Précise les pôles et écris l'équation-bilan de fonctionnement.
\item Pour la pile choisie en 2., indique quel métal voit sa masse augmenter au cours du fonctionnement. Justifie.
\end{enumerate}
\end{exercice}

\courssub{Je me prépare au Probatoire}

\begin{exercice}[La pile Daniell au laboratoire de Nkongsamba]
Au laboratoire d'un lycée de Nkongsamba, un professeur demande à ses élèves de Première C de réaliser une pile Daniell. Ils disposent d'une lame de zinc, d'une lame de cuivre, d'une solution de sulfate de zinc(II) de concentration \SI{1,0}{mol.L^{-1}}, d'une solution de sulfate de cuivre(II) de concentration \SI{1,0}{mol.L^{-1}}, d'un tube en U rempli d'une solution gélifiée de nitrate de potassium et d'un voltmètre.

On donne : $E^{\circ}(\ce{Zn^2+}/\ce{Zn}) = \SI{-0,76}{V}$ ; $E^{\circ}(\ce{Cu^2+}/\ce{Cu}) = \SI{+0,34}{V}$.
\begin{enumerate}
\item \textbf{Schéma.} Fais le schéma annoté de la pile réalisée : nature des électrodes et des solutions, pont salin, voltmètre, pôles $+$ et $-$, sens des électrons et sens du courant conventionnel.
\item \textbf{Polarité.} Justifie, à l'aide des potentiels standard, que le zinc constitue le pôle négatif.
\item \textbf{Fonctionnement.} Écris l'équation de la réaction à chaque électrode en précisant s'il s'agit d'une oxydation ou d'une réduction, puis l'équation-bilan de fonctionnement de la pile.
\item \textbf{F.é.m.} Calcule la f.é.m. standard de cette pile. Que lit-on sur le voltmètre si la pile ne débite pas ?
\item \textbf{Notation.} Donne la notation conventionnelle de la pile.
\item \textbf{Évolution.} Après un long fonctionnement, comment évoluent la masse de la lame de zinc, la masse de la lame de cuivre et la couleur bleue de la solution de sulfate de cuivre ? Justifie chaque réponse.
\item \textbf{Rôle du pont salin.} Explique pourquoi la pile cesse de fonctionner si l'on retire le pont salin.
\end{enumerate}
\end{exercice}

\begin{exercice}[Identification d'un métal inconnu]
Un laboratoire de la SONARA récupère un métal inconnu noté $X$, associé aux ions \ce{X^2+}. Pour l'identifier, un technicien réalise des piles avec des électrodes de référence connues et mesure les f.é.m. suivantes (toutes les solutions sont à \SI{1,0}{mol.L^{-1}}) :

\begin{center}
\begin{tabularx}{\linewidth}{|l|X|c|}
\hline
\rowcolor{popblueL} \textbf{Pile} & \textbf{Couples associés} & \textbf{F.é.m. mesurée} \\
\hline
Pile A & \ce{Zn^2+}/\ce{Zn} et \ce{X^2+}/\ce{X} ; le zinc est le pôle $-$ & \SI{0,32}{V} \\
\hline
Pile B & \ce{X^2+}/\ce{X} et \ce{Cu^2+}/\ce{Cu} ; le métal $X$ est le pôle $-$ & \SI{0,78}{V} \\
\hline
\end{tabularx}
\end{center}

On donne : $E^{\circ}(\ce{Zn^2+}/\ce{Zn}) = \SI{-0,76}{V}$ ; $E^{\circ}(\ce{Cu^2+}/\ce{Cu}) = \SI{+0,34}{V}$ ; $E^{\circ}(\ce{Fe^2+}/\ce{Fe}) = \SI{-0,44}{V}$ ; $E^{\circ}(\ce{Pb^2+}/\ce{Pb}) = \SI{-0,13}{V}$ ; $E^{\circ}(\ce{Ni^2+}/\ce{Ni}) = \SI{-0,23}{V}$.
\begin{enumerate}
\item Schématise la pile A en précisant ses pôles et le sens de circulation des électrons.
\item À partir de la pile A, détermine le potentiel standard $E^{\circ}(\ce{X^2+}/\ce{X})$.
\item Vérifie ce résultat à l'aide de la pile B.
\item Parmi le fer, le plomb et le nickel, identifie le métal $X$. Justifie.
\item Écris les équations aux électrodes et l'équation-bilan de fonctionnement de la pile B.
\item Classe les quatre couples \ce{Zn^2+}/\ce{Zn}, \ce{X^2+}/\ce{X}, \ce{Cu^2+}/\ce{Cu} et \ce{Ag+}/\ce{Ag} ($E^{\circ} = \SI{+0,80}{V}$) par potentiel croissant, puis les métaux par pouvoir réducteur croissant.
\item Le métal $X$ peut-il être attaqué par une solution d'acide chlorhydrique (couple \ce{H+}/\ce{H2}, $E^{\circ} = \SI{0,00}{V}$) ? Et le cuivre ? Justifie.
\end{enumerate}
\end{exercice}

\begin{exercice}[Une pile Daniell pour éclairer un village]
Dans un village de la région de l'Adamaoua non raccordé au réseau ENEO, un jeune ingénieur utilise une pile Daniell de grande capacité pour alimenter une lampe. La pile débite un courant d'intensité constante $I = \SI{0,10}{A}$ pendant une durée $\Delta t = \SI{5,0}{h}$. L'électrode de zinc est la seule source d'électrons.

On donne : $M(\ce{Zn}) = \SI{65,4}{g.mol^{-1}}$ ; $M(\ce{Cu}) = \SI{63,5}{g.mol^{-1}}$ ; constante de Faraday $F = \SI{96500}{C.mol^{-1}}$ ; $E^{\circ}(\ce{Zn^2+}/\ce{Zn}) = \SI{-0,76}{V}$ ; $E^{\circ}(\ce{Cu^2+}/\ce{Cu}) = \SI{+0,34}{V}$.
\begin{enumerate}
\item Rappelle l'équation-bilan de fonctionnement de la pile Daniell et calcule sa f.é.m. standard.
\item Calcule la quantité d'électricité $Q$ fournie par la pile pendant les $\SI{5,0}{h}$ de fonctionnement.
\item Montre que la quantité de matière d'électrons échangée vaut $n(e^-) = \dfrac{Q}{F}$, puis calcule sa valeur.
\item En déduis la quantité de matière de zinc consommée, puis la masse de zinc disparue pendant ces $\SI{5,0}{h}$.
\item Calcule la masse de cuivre déposée sur l'électrode de cuivre pendant la même durée.
\item La lame de zinc initiale a une masse de \SI{10}{g}. Calcule la durée maximale théorique de fonctionnement de la pile à \SI{0,10}{A} avant épuisement complet du zinc.
\item En réalité, la lampe s'éteint avant cette durée. Propose deux raisons expérimentales possibles.
\end{enumerate}
\end{exercice}



\newpage

\coursec{Corrigés des exercices}

\courssub{Je vérifie mes connaissances}

\begin{corrige}[Exercice 1 — QCM : la pile Daniell]
\begin{enumerate}
\item[\textbf{a)}] \textbf{Faux.} Les électrons ne circulent \emph{jamais} dans le pont salin : celui-ci est un conducteur \emph{ionique}. Correction : dans le pont salin, ce sont les ions qui migrent (\ce{K+} vers le compartiment du cuivre, \ce{NO3-} vers celui du zinc). Les électrons circulent uniquement dans le circuit métallique extérieur, du zinc vers le cuivre.
\item[\textbf{b)}] \textbf{Vrai.} L'oxydation du zinc libère des électrons (\ce{Zn -> Zn^2+ + 2e-}) qui traversent le fil conducteur et le récepteur jusqu'à l'électrode de cuivre, où ils sont consommés par la réduction des ions \ce{Cu^2+}.
\item[\textbf{c)}] \textbf{Vrai.} Par définition, l'\cle{anode} est le siège de l'oxydation. Dans une pile (fonctionnement en générateur), l'anode est le pôle \textbf{négatif} : c'est elle qui fournit les électrons au circuit extérieur.
\item[\textbf{d)}] \textbf{Faux.} À l'électrode de cuivre (cathode, pôle $+$), il se produit une \emph{réduction} : \ce{Cu^2+ + 2e- -> Cu}. L'équation proposée \ce{Cu -> Cu^2+ + 2e-} est une oxydation : elle aurait lieu si le cuivre était l'anode, ce qui n'est pas le cas ici.
\item[\textbf{e)}] \textbf{Vrai.} Le pont salin (solution gélifiée de \ce{KNO3}) ferme le circuit électrique en assurant la conduction ionique entre les deux compartiments et maintient l'électroneutralité de chaque solution : les anions migrent vers l'anode (compensation des \ce{Zn^2+} formés), les cations vers la cathode (compensation des \ce{Cu^2+} consommés).
\item[\textbf{f)}] \textbf{Vrai.} $E^{\circ}_{\text{pile}} = E^{\circ}_{\text{cathode}} - E^{\circ}_{\text{anode}} = E^{\circ}(\ce{Cu^2+}/\ce{Cu}) - E^{\circ}(\ce{Zn^2+}/\ce{Zn}) = 0,34 - (-0,76) = \SI{1,10}{V}$.
\end{enumerate}
\textbf{Bilan :} les affirmations vraies sont \textbf{b, c, e, f}.
\end{corrige}

\begin{corrige}[Exercice 2 — Vrai ou faux, avec justification]
\begin{enumerate}
\item \textbf{Faux.} La convention $E^{\circ}(\ce{H+}/\ce{H2}) = \SI{0,00}{V}$ est fixée à \SI{25}{\celsius} (\SI{298}{K}), dans les conditions standard ($p(\ce{H2}) = \SI{1}{bar}$, $[\ce{H+}] = \SI{1}{mol.L^{-1}}$). Le potentiel de l'électrode à hydrogène dépend de la température : la convention n'est donc valable qu'à \SI{25}{\celsius}.
\item \textbf{Faux.} C'est l'inverse : plus le potentiel standard d'un couple est \emph{élevé}, plus l'\emph{oxydant} du couple est fort (et plus son réducteur conjugué est faible). Un réducteur fort appartient à un couple de potentiel standard très \emph{négatif} (ex. \ce{Mg^2+}/\ce{Mg}, $E^{\circ} = \SI{-2,37}{V}$).
\item \textbf{Vrai.} Dans une pile, l'oxydation a toujours lieu à l'anode, qui est le pôle négatif (elle cède des électrons au circuit extérieur). \emph{Nuance pour la suite :} en électrolyse (récepteur), l'anode reste le siège de l'oxydation mais c'est le pôle positif.
\item \textbf{Vrai.} C'est la convention internationale : le pôle négatif (anode) s'écrit à gauche, le pôle positif (cathode) à droite, la double barre représentant le pont salin.
\item \textbf{Faux.} Les piles usuelles (saline, alcaline, bouton) fonctionnent très bien : elles contiennent un \emph{équivalent} du pont salin, à savoir un électrolyte pâteux ou gélifié (\ce{NH4Cl} pour les piles salines, \ce{KOH} pour les piles alcalines) imprégnant un séparateur poreux. Celui-ci assure la conduction ionique et l'électroneutralité tout en empêchant le mélange direct des réactifs.
\end{enumerate}
\end{corrige}

\begin{corrige}[Exercice 3 — Texte à trous]
La pile Daniell associe deux demi-piles : une lame de zinc plongée dans une solution de sulfate de zinc et une lame de cuivre plongée dans une solution de sulfate de cuivre. Les deux solutions sont reliées par un \textbf{pont salin} qui assure le passage du courant. L'électrode de zinc, siège d'une \textbf{oxydation}, est l'\textbf{anode} de la pile ; elle constitue son pôle \textbf{négatif}. L'électrode de cuivre, siège d'une \textbf{réduction}, est la \textbf{cathode} ; elle constitue le pôle \textbf{positif}. La tension mesurée entre les deux électrodes lorsque la pile ne débite pas s'appelle la \textbf{f.é.m.} de la pile ; elle s'exprime en \textbf{volt}. Pour mesurer les potentiels d'oxydoréduction, on utilise une demi-pile de référence : l'électrode normale à \textbf{hydrogène}.
\end{corrige}

\begin{corrige}[Exercice 4 — Définitions et conventions]
\begin{enumerate}
\item \textbf{Définitions.}
\begin{itemize}
\item \cle{Demi-pile} : ensemble constitué d'un conducteur électronique (une électrode, par exemple une lame métallique) plongé dans un conducteur ionique (solution électrolytique) contenant les espèces d'un couple oxydant/réducteur. Exemple : une lame de zinc dans une solution de \ce{Zn^2+} constitue la demi-pile du couple \ce{Zn^2+}/\ce{Zn}.
\item \cle{Potentiel standard d'oxydoréduction} $E^{\circ}$ : différence de potentiel, mesurée à \SI{25}{\celsius} dans les conditions standard, entre l'électrode du couple étudié et l'électrode normale à hydrogène (ENH) prise comme référence ($E^{\circ} = \SI{0,00}{V}$). Il traduit quantitativement la force de l'oxydant du couple.
\item \cle{F.é.m.} (force électromotrice) : différence de potentiel entre les deux électrodes d'une pile lorsqu'elle ne débite aucun courant (circuit ouvert) : $E = E_{(+)} - E_{(-)} = E_{\text{cathode}} - E_{\text{anode}}$.
\end{itemize}
\item \textbf{Conditions standard :} température $\theta = \SI{25}{\celsius}$ ; pression des gaz $p = \SI{1}{bar}$ ; concentration \SI{1}{mol.L^{-1}} pour toutes les espèces dissoutes ; solides et liquides purs (activité égale à 1).
\item \textbf{Notation conventionnelle de la pile Daniell :}
\[ \ce{Zn(s) | Zn^2+(aq) || Cu^2+(aq) | Cu(s)} \]
La barre simple \texttt{|} marque une interface entre deux phases différentes (métal/solution) ; la double barre \texttt{||} représente le pont salin, c'est-à-dire la jonction entre les deux demi-piles. Le pôle négatif est écrit à gauche.
\item \textbf{L'ENH, électrode de référence :} on ne peut mesurer qu'une \emph{différence} de potentiel, jamais un potentiel absolu. Il faut donc un étalon : par convention, on attribue à l'ENH le potentiel $E^{\circ} = \SI{0,00}{V}$ à \SI{25}{\celsius}, et tous les autres potentiels sont mesurés par rapport à elle. \textbf{Constitution :} une lame de platine platiné (métal inerte jouant le rôle de conducteur et de catalyseur) plonge dans une solution acide telle que $[\ce{H3O+}] = \SI{1}{mol.L^{-1}}$, sur laquelle barbote du dihydrogène gazeux pur sous \SI{1}{bar} à \SI{25}{\celsius}. Demi-équation associée : \ce{2H+ + 2e- <=> H2(g)}.
\end{enumerate}
\end{corrige}

\courssub{Je m'entraîne}

\begin{corrige}[Exercice 5 — La pile fer–argent]
\begin{enumerate}
\item \textbf{Polarité.} On compare les potentiels standard :
\[ E^{\circ}(\ce{Ag+}/\ce{Ag}) = \SI{+0,80}{V} > E^{\circ}(\ce{Fe^2+}/\ce{Fe}) = \SI{-0,44}{V} \]
Le couple de potentiel le plus \emph{élevé} (\ce{Ag+}/\ce{Ag}) fournit la \textbf{cathode} (pôle $+$) : son oxydant \ce{Ag+} est réduit. Le couple de potentiel le plus \emph{faible} (\ce{Fe^2+}/\ce{Fe}) fournit l'\textbf{anode} (pôle $-$) : son réducteur \ce{Fe} est oxydé.
\item \textbf{Équations aux électrodes.}
\begin{itemize}
\item Anode ($-$), oxydation : \ce{Fe(s) -> Fe^2+(aq) + 2e-}
\item Cathode ($+$), réduction : \ce{Ag+(aq) + e- -> Ag(s)} \quad (à multiplier par 2)
\end{itemize}
\textbf{Équation-bilan de fonctionnement :}
\[ \ce{Fe(s) + 2Ag+(aq) -> Fe^2+(aq) + 2Ag(s)} \]
\item \textbf{F.é.m. standard :}
\[ E^{\circ} = E^{\circ}_{\text{cathode}} - E^{\circ}_{\text{anode}} = 0,80 - (-0,44) = \SI{+1,24}{V} \]
\item \textbf{Notation conventionnelle :}
\[ \ce{Fe(s) | Fe^2+(aq) || Ag+(aq) | Ag(s)} \]
\item \textbf{Circulation des porteurs de charge.}
\begin{itemize}
\item Dans le circuit extérieur, les \textbf{électrons} circulent de l'électrode de fer (anode, $-$) vers l'électrode d'argent (cathode, $+$) ; le courant conventionnel circule en sens inverse.
\item Dans le pont salin au nitrate de potassium : les cations \ce{K+} migrent vers le compartiment de l'argent (pour compenser la disparition des \ce{Ag+}) et les anions \ce{NO3-} migrent vers le compartiment du fer (pour compenser l'apparition des \ce{Fe^2+}).
\end{itemize}
\end{enumerate}
\begin{attention}[Points de vigilance]
Ne pas oublier de multiplier la demi-équation de l'argent par 2 avant d'écrire le bilan : les électrons ne doivent \emph{jamais} apparaître dans l'équation-bilan. La f.é.m. se calcule toujours par $E_{(+)} - E_{(-)}$, jamais par une somme.
\end{attention}
\end{corrige}

\begin{corrige}[Exercice 6 — Prévision de réactions à l'aide des potentiels]
\begin{enumerate}
\item \textbf{Classement des couples par potentiel croissant :}
\[ \ce{Fe^2+}/\ce{Fe}\;(-0,44) < \ce{Pb^2+}/\ce{Pb}\;(-0,13) < \ce{Cu^2+}/\ce{Cu}\;(+0,34)\ \text{(en V)} \]
\textbf{Métaux par pouvoir réducteur croissant :} le réducteur est d'autant plus fort que le potentiel du couple est plus bas :
\[ \ce{Cu} < \ce{Pb} < \ce{Fe} \]
\item \textbf{Lame de plomb dans une solution de \ce{CuSO4} : réaction, oui.} On place les couples sur un axe vertical de potentiels croissants vers le haut ; la règle du gamma indique que l'oxydant le plus fort (\ce{Cu^2+}, couple le plus haut) réagit avec le réducteur le plus fort (\ce{Pb}, couple le plus bas) :
\[ \ce{Pb(s) + Cu^2+(aq) -> Pb^2+(aq) + Cu(s)} \]
Vérification : $\Delta E^{\circ} = E^{\circ}(\ce{Cu^2+}/\ce{Cu}) - E^{\circ}(\ce{Pb^2+}/\ce{Pb}) = 0,34 - (-0,13) = \SI{+0,47}{V} > 0$ : la réaction est spontanée. On observe un dépôt rougeâtre de cuivre sur le plomb et la décoloration progressive de la solution bleue.
\item \textbf{Lame de plomb dans une solution de \ce{FeSO4} : pas de réaction.} Le plomb ($E^{\circ} = \SI{-0,13}{V}$) est un réducteur \emph{plus faible} que le fer ($E^{\circ} = \SI{-0,44}{V}$) : il ne peut pas réduire les ions \ce{Fe^2+}. La règle du gamma ne se referme pas ; $\Delta E^{\circ} = -0,44 - (-0,13) = \SI{-0,31}{V} < 0$.
\item \textbf{Le cuivre peut-il réduire \ce{Fe^2+} ? Non.} Le cuivre ($E^{\circ} = \SI{+0,34}{V}$) est un réducteur beaucoup plus faible que le fer : $\Delta E^{\circ} = -0,44 - 0,34 = \SI{-0,78}{V} < 0$. Aucune réaction spontanée entre \ce{Cu(s)} et \ce{Fe^2+(aq)}.
\end{enumerate}
\begin{methode}[Règle du gamma — réflexe]
Pour prévoir une réaction entre deux couples Ox$_1$/Red$_1$ et Ox$_2$/Red$_2$ : place les couples sur un axe vertical (potentiels croissants vers le haut), relie l'oxydant du couple du haut au réducteur du couple du bas : si le « gamma » se referme, la réaction est spontanée ; sinon, elle ne se produit pas.
\end{methode}
\end{corrige}

\begin{corrige}[Exercice 7 — Schéma et mesures de f.é.m.]
\begin{enumerate}
\item \textbf{Schéma de la pile 1 (pile Daniell).}

\begin{popfigure}
\centering
\begin{tikzpicture}[scale=0.95, >=Stealth]
% Bécher gauche : Zn
\draw[thick] (0,0) -- (0,-2.6) -- (2.2,-2.6) -- (2.2,0);
\draw[fill=popblue!20, draw=none] (0.08,-2.55) rectangle (2.12,-1.1);
\node[popdark] at (1.1,-1.9) {\footnotesize \ce{ZnSO4}};
\draw[thick, fill=popmuted] (0.85,-2.6) rectangle (1.25,0.4);
\node[above] at (1.05,0.4) {\footnotesize \ce{Zn} ($-$)};
% Bécher droit : Cu
\draw[thick] (5.6,0) -- (5.6,-2.6) -- (7.8,-2.6) -- (7.8,0);
\draw[fill=popblue!50, draw=none] (5.68,-2.55) rectangle (7.72,-1.1);
\node[white] at (6.7,-1.9) {\footnotesize \ce{CuSO4}};
\draw[thick, fill=poporange] (6.45,-2.6) rectangle (6.85,0.4);
\node[above] at (6.65,0.4) {\footnotesize \ce{Cu} ($+$)};
% Pont salin
\draw[thick, fill=popgreenL] (1.6,-1.4) -- (1.6,-0.6) .. controls (3.9,0.1) .. (6.2,-0.6) -- (6.2,-1.4) -- (5.7,-1.4) -- (5.7,-0.95) .. controls (3.9,-0.35) .. (2.1,-0.95) -- (2.1,-1.4) -- cycle;
\node[popdark] at (3.9,-0.55) {\footnotesize pont salin \ce{KNO3}};
% Voltmètre
\draw[thick] (1.05,0.4) -- (1.05,1.3) -- (3.4,1.3);
\draw[thick] (6.65,0.4) -- (6.65,1.3) -- (4.4,1.3);
\draw[thick, fill=white] (3.9,1.3) circle (0.5);
\node at (3.9,1.3) {V};
% Sens des électrons
\draw[->, very thick, popink] (1.6,1.75) -- (6.2,1.75);
\node[popink] at (3.9,2.05) {\footnotesize sens des électrons $e^-$};
\end{tikzpicture}
\legende{Pile Daniell : anode de zinc (pôle $-$), cathode de cuivre (pôle $+$), pont salin au \ce{KNO3}, voltmètre en dérivation.}
\end{popfigure}

\item \textbf{Potentiel du couple \ce{Cu^2+}/\ce{Cu}.} Pour la pile 1, le zinc est le pôle $-$ :
\[ E_1 = E^{\circ}(\ce{Cu^2+}/\ce{Cu}) - E^{\circ}(\ce{Zn^2+}/\ce{Zn}) \]
\[ E^{\circ}(\ce{Cu^2+}/\ce{Cu}) = E_1 + E^{\circ}(\ce{Zn^2+}/\ce{Zn}) = 1,10 + (-0,76) = \SI{+0,34}{V} \]
\item \textbf{Potentiel du couple \ce{Ag+}/\ce{Ag}.} Pour la pile 2, le zinc reste le pôle $-$ (son potentiel est le plus faible) :
\[ E^{\circ}(\ce{Ag+}/\ce{Ag}) = E_2 + E^{\circ}(\ce{Zn^2+}/\ce{Zn}) = 1,56 + (-0,76) = \SI{+0,80}{V} \]
\item \textbf{Vérification avec la pile 3.} Prévision :
\[ E_3 = E^{\circ}(\ce{Ag+}/\ce{Ag}) - E^{\circ}(\ce{Cu^2+}/\ce{Cu}) = 0,80 - 0,34 = \SI{+0,46}{V} \]
La valeur mesurée ($\SI{0,46}{V}$) coïncide exactement avec la valeur calculée : les trois mesures sont \textbf{cohérentes}. On remarque la relation d'additivité : $E_2 = E_1 + E_3$ ($1,10 + 0,46 = 1,56$~V).
\end{enumerate}
\end{corrige}

\begin{corrige}[Exercice 8 — Construction d'une pile de f.é.m. imposée]
\begin{enumerate}
\item \textbf{Pile de plus grande f.é.m.} La f.é.m. $E = E^{\circ}_{\text{cathode}} - E^{\circ}_{\text{anode}}$ est maximale en associant le couple de potentiel le plus \emph{haut} (cathode : \ce{Cu^2+}/\ce{Cu}, $+0,34$~V) au couple de potentiel le plus \emph{bas} (anode : \ce{Mg^2+}/\ce{Mg}, $-2,37$~V) :
\[ E^{\circ}_{\max} = 0,34 - (-2,37) = \SI{+2,71}{V} \]
Notation conventionnelle : \ce{Mg(s) | Mg^2+(aq) || Cu^2+(aq) | Cu(s)}.
\item \textbf{F.é.m. voisine de \SI{0,47}{V}.} On calcule les différences possibles entre couples voisins :

\begin{center}
\begin{tabularx}{\linewidth}{|l|X|c|}
\hline
\rowcolor{popblueL} \textbf{Association} & \textbf{Calcul} & \textbf{$E^{\circ}$ (V)} \\
\hline
\ce{Pb}/\ce{Cu} & $0,34 - (-0,13)$ & $0,47$ \\
\hline
\ce{Zn}/\ce{Pb} & $-0,13 - (-0,76)$ & $0,63$ \\
\hline
\ce{Zn}/\ce{Cu} & $0,34 - (-0,76)$ & $1,10$ \\
\hline
\ce{Mg}/\ce{Pb} & $-0,13 - (-2,37)$ & $2,24$ \\
\hline
\end{tabularx}
\end{center}

La combinaison cherchée est \textbf{\ce{Pb^2+}/\ce{Pb} et \ce{Cu^2+}/\ce{Cu}} : $E^{\circ} = \SI{+0,47}{V}$.
\textbf{Pôles :} anode ($-$) = plomb ; cathode ($+$) = cuivre.
\textbf{Équations :} anode : \ce{Pb(s) -> Pb^2+(aq) + 2e-} ; cathode : \ce{Cu^2+(aq) + 2e- -> Cu(s)}.
\textbf{Bilan :} \ce{Pb(s) + Cu^2+(aq) -> Pb^2+(aq) + Cu(s)}.
\item \textbf{Évolution des masses.} C'est la lame de \textbf{cuivre} dont la masse augmente : à la cathode, les ions \ce{Cu^2+} sont réduits et le cuivre métallique formé se dépose sur la lame (\ce{Cu^2+ + 2e- -> Cu(s)}). La lame de plomb, elle, s'amincit : le métal s'oxyde et passe en solution sous forme d'ions \ce{Pb^2+}.
\end{enumerate}
\end{corrige}

\courssub{Je me prépare au Probatoire}

\begin{corrige}[Exercice 9 — La pile Daniell au laboratoire de Nkongsamba]
\begin{enumerate}
\item \textbf{Schéma annoté.}

\begin{popfigure}
\centering
\begin{tikzpicture}[scale=0.95, >=Stealth]
% Bécher gauche : Zn
\draw[thick] (0,0) -- (0,-2.6) -- (2.2,-2.6) -- (2.2,0);
\draw[fill=popblue!15, draw=none] (0.08,-2.55) rectangle (2.12,-1.1);
\node[popdark] at (1.1,-1.95) {\footnotesize \ce{ZnSO4} \SI{1,0}{mol.L^{-1}}};
\draw[thick, fill=popmuted] (0.85,-2.6) rectangle (1.25,0.4);
\node[above] at (1.05,0.4) {\footnotesize \ce{Zn} : anode ($-$)};
% Bécher droit : Cu
\draw[thick] (6.4,0) -- (6.4,-2.6) -- (8.6,-2.6) -- (8.6,0);
\draw[fill=popblue!50, draw=none] (6.48,-2.55) rectangle (8.52,-1.1);
\node[white] at (7.5,-1.95) {\footnotesize \ce{CuSO4} \SI{1,0}{mol.L^{-1}}};
\draw[thick, fill=poporange] (7.25,-2.6) rectangle (7.65,0.4);
\node[above] at (7.45,0.4) {\footnotesize \ce{Cu} : cathode ($+$)};
% Pont salin
\draw[thick, fill=popgreenL] (1.6,-1.4) -- (1.6,-0.6) .. controls (4.3,0.1) .. (7.0,-0.6) -- (7.0,-1.4) -- (6.5,-1.4) -- (6.5,-0.95) .. controls (4.3,-0.35) .. (2.1,-0.95) -- (2.1,-1.4) -- cycle;
\node[popdark] at (4.3,-0.5) {\footnotesize pont salin : \ce{KNO3} gélifié};
% Voltmètre
\draw[thick] (1.05,0.4) -- (1.05,1.4) -- (3.8,1.4);
\draw[thick] (7.45,0.4) -- (7.45,1.4) -- (4.8,1.4);
\draw[thick, fill=white] (4.3,1.4) circle (0.5);
\node at (4.3,1.4) {V};
\node at (4.3,2.15) {\footnotesize $U = \SI{1,10}{V}$};
% Sens des électrons et du courant
\draw[->, very thick, popink] (1.6,1.85) -- (7.0,1.85);
\node[popink] at (3.0,2.15) {\footnotesize électrons $e^-$};
\draw[->, very thick, poppurple] (7.0,2.5) -- (1.6,2.5);
\node[poppurple] at (5.8,2.8) {\footnotesize courant $I$};
\end{tikzpicture}
\legende{Pile Daniell : les électrons vont du zinc vers le cuivre dans le circuit extérieur ; le courant conventionnel circule en sens inverse.}
\end{popfigure}

\item \textbf{Polarité.} $E^{\circ}(\ce{Zn^2+}/\ce{Zn}) = \SI{-0,76}{V} < E^{\circ}(\ce{Cu^2+}/\ce{Cu}) = \SI{+0,34}{V}$. Le zinc est le réducteur le plus fort : il s'oxyde en cédant des électrons ; l'électrode de zinc est donc l'anode, pôle \textbf{négatif}. Le cuivre constitue la cathode, pôle \textbf{positif}.
\item \textbf{Fonctionnement.}
\begin{itemize}
\item Anode ($-$), \textbf{oxydation} : \ce{Zn(s) -> Zn^2+(aq) + 2e-}
\item Cathode ($+$), \textbf{réduction} : \ce{Cu^2+(aq) + 2e- -> Cu(s)}
\item \textbf{Équation-bilan :} \ce{Zn(s) + Cu^2+(aq) -> Zn^2+(aq) + Cu(s)}
\end{itemize}
\item \textbf{F.é.m. standard :}
\[ E^{\circ} = E^{\circ}(\ce{Cu^2+}/\ce{Cu}) - E^{\circ}(\ce{Zn^2+}/\ce{Zn}) = 0,34 - (-0,76) = \SI{+1,10}{V} \]
Si la pile ne débite pas (circuit ouvert), le voltmètre indique la f.é.m., soit environ \SI{1,10}{V} (borne COM reliée au zinc).
\item \textbf{Notation conventionnelle :}
\[ \ce{Zn(s) | Zn^2+(aq) (\SI{1,0}{mol.L^{-1}}) || Cu^2+(aq) (\SI{1,0}{mol.L^{-1}}) | Cu(s)} \]
\item \textbf{Évolution après un long fonctionnement.}
\begin{itemize}
\item \textbf{Lame de zinc :} sa masse \emph{diminue} — le métal est consommé par l'oxydation \ce{Zn -> Zn^2+ + 2e-}.
\item \textbf{Lame de cuivre :} sa masse \emph{augmente} — le cuivre formé par réduction des \ce{Cu^2+} se dépose sur elle.
\item \textbf{Solution de sulfate de cuivre :} sa couleur bleue \emph{pâlit} progressivement, car la concentration des ions \ce{Cu^2+} (responsables de la couleur bleue) diminue au cours de la réduction.
\end{itemize}
\item \textbf{Rôle du pont salin.} Si l'on retire le pont salin, le circuit n'est plus fermé : aucune migration ionique n'est possible. Le compartiment anodique se chargerait positivement (accumulation de \ce{Zn^2+}) et le compartiment cathodique négativement (excès d'anions \ce{SO4^2-}) ; ces excès de charge s'opposent immédiatement à la poursuite des réactions : la pile cesse de débiter. Le pont salin assure la \cle{fermeture du circuit} et l'\cle{électroneutralité} des solutions grâce à la migration de \ce{K+} vers la cathode et de \ce{NO3-} vers l'anode.
\end{enumerate}
\textbf{Barème indicatif (sur 10) :} schéma annoté 2 pts ; polarité justifiée 1 pt ; demi-équations + bilan 2 pts ; f.é.m. 1 pt ; notation 1 pt ; évolutions 2 pts ; rôle du pont salin 1 pt.
\begin{attention}[Points de vigilance]
Le sens du courant conventionnel est \emph{opposé} à celui des électrons. La couleur bleue est due aux ions \ce{Cu^2+} en solution, pas au métal cuivre. Ne jamais écrire d'électrons dans l'équation-bilan.
\end{attention}
\end{corrige}

\begin{corrige}[Exercice 10 — Identification d'un métal inconnu (SONARA)]
\begin{enumerate}
\item \textbf{Schéma de la pile A.} Le zinc est le pôle $-$ (anode) ; le métal $X$ est donc le pôle $+$ (cathode). Les électrons circulent dans le circuit extérieur du zinc vers $X$.

\begin{popfigure}
\centering
\begin{tikzpicture}[scale=0.9, >=Stealth]
\draw[thick] (0,0) -- (0,-2.4) -- (2.0,-2.4) -- (2.0,0);
\draw[fill=popblue!15, draw=none] (0.08,-2.35) rectangle (1.92,-1.0);
\node[popdark] at (1.0,-1.7) {\footnotesize \ce{Zn^2+(aq)}};
\draw[thick, fill=popmuted] (0.75,-2.4) rectangle (1.15,0.4);
\node[above] at (0.95,0.4) {\footnotesize \ce{Zn} ($-$)};
\draw[thick] (5.4,0) -- (5.4,-2.4) -- (7.4,-2.4) -- (7.4,0);
\draw[fill=poppinkL, draw=none] (5.48,-2.35) rectangle (7.32,-1.0);
\node[popdark] at (6.4,-1.7) {\footnotesize \ce{X^2+(aq)}};
\draw[thick, fill=poppurple!60] (6.15,-2.4) rectangle (6.55,0.4);
\node[above] at (6.35,0.4) {\footnotesize $X$ ($+$)};
\draw[thick, fill=popgreenL] (1.5,-1.3) -- (1.5,-0.6) .. controls (3.7,0.05) .. (5.9,-0.6) -- (5.9,-1.3) -- (5.4,-1.3) -- (5.4,-0.9) .. controls (3.7,-0.3) .. (2.0,-0.9) -- (2.0,-1.3) -- cycle;
\node[popdark] at (3.7,-0.45) {\footnotesize pont salin};
\draw[thick] (0.95,0.4) -- (0.95,1.2) -- (3.2,1.2);
\draw[thick] (6.35,0.4) -- (6.35,1.2) -- (4.2,1.2);
\draw[thick, fill=white] (3.7,1.2) circle (0.45);
\node at (3.7,1.2) {\footnotesize V};
\draw[->, very thick, popink] (1.5,1.7) -- (5.9,1.7);
\node[popink] at (3.7,2.0) {\footnotesize $e^-$};
\end{tikzpicture}
\legende{Pile A : \ce{Zn} (anode, $-$) et $X$ (cathode, $+$) ; électrons de \ce{Zn} vers $X$.}
\end{popfigure}

\item \textbf{Potentiel standard de \ce{X^2+}/\ce{X} (pile A).} Le zinc étant le pôle $-$ :
\[ E_A = E^{\circ}(\ce{X^2+}/\ce{X}) - E^{\circ}(\ce{Zn^2+}/\ce{Zn}) \]
\[ E^{\circ}(\ce{X^2+}/\ce{X}) = E_A + E^{\circ}(\ce{Zn^2+}/\ce{Zn}) = 0,32 + (-0,76) = \SI{-0,44}{V} \]
\item \textbf{Vérification (pile B).} Le métal $X$ est le pôle $-$ :
\[ E_B = E^{\circ}(\ce{Cu^2+}/\ce{Cu}) - E^{\circ}(\ce{X^2+}/\ce{X}) \]
\[ E^{\circ}(\ce{X^2+}/\ce{X}) = E^{\circ}(\ce{Cu^2+}/\ce{Cu}) - E_B = 0,34 - 0,78 = \SI{-0,44}{V} \]
Les deux mesures concordent : $E^{\circ}(\ce{X^2+}/\ce{X}) = \SI{-0,44}{V}$.
\item \textbf{Identification.} Parmi les valeurs données, $E^{\circ}(\ce{Fe^2+}/\ce{Fe}) = \SI{-0,44}{V}$ : le métal inconnu $X$ est le \textbf{fer} \ce{Fe}.
\item \textbf{Fonctionnement de la pile B.}
\begin{itemize}
\item Anode ($-$), oxydation : \ce{Fe(s) -> Fe^2+(aq) + 2e-}
\item Cathode ($+$), réduction : \ce{Cu^2+(aq) + 2e- -> Cu(s)}
\item \textbf{Bilan :} \ce{Fe(s) + Cu^2+(aq) -> Fe^2+(aq) + Cu(s)}
\end{itemize}
\item \textbf{Classements.}
\begin{itemize}
\item Couples par potentiel croissant :
\[ \ce{Zn^2+}/\ce{Zn}\;(-0,76) < \ce{Fe^2+}/\ce{Fe}\;(-0,44) < \ce{Cu^2+}/\ce{Cu}\;(+0,34) < \ce{Ag+}/\ce{Ag}\;(+0,80)\ \text{(V)} \]
\item Métaux par pouvoir réducteur croissant (ordre inverse) : $\ce{Ag} < \ce{Cu} < \ce{Fe} < \ce{Zn}$.
\end{itemize}
\item \textbf{Attaque par l'acide chlorhydrique.} Un métal réduit les ions \ce{H+} (dégagement de \ce{H2}) si le potentiel de son couple est \emph{inférieur} à celui du couple \ce{H+}/\ce{H2} ($\SI{0,00}{V}$).
\begin{itemize}
\item Métal $X$ = fer : $E^{\circ}(\ce{Fe^2+}/\ce{Fe}) = \SI{-0,44}{V} < \SI{0,00}{V}$ : \textbf{oui}, le fer est attaqué : \ce{Fe(s) + 2H+ -> Fe^2+(aq) + H2(g)}.
\item Cuivre : $E^{\circ}(\ce{Cu^2+}/\ce{Cu}) = \SI{+0,34}{V} > \SI{0,00}{V}$ : \textbf{non}, le cuivre n'est pas attaqué par l'acide chlorhydrique (acide non oxydant).
\end{itemize}
\end{enumerate}
\textbf{Barème indicatif (sur 12) :} schéma 1,5 pt ; calcul pile A 1,5 pt ; vérification pile B 1,5 pt ; identification 1 pt ; équations pile B 2 pts ; classements 2 pts ; attaque acide 2,5 pts.
\begin{attention}[Points de vigilance]
Le signe de $E^{\circ}(\ce{X^2+}/\ce{X})$ dépend du sens de la formule $E = E_{(+)} - E_{(-)}$ : identifie d'abord quel métal est le pôle $+$ dans chaque pile avant de calculer. Une identification n'est valide que si les deux mesures donnent la même valeur.
\end{attention}
\end{corrige}

\begin{corrige}[Exercice 11 — Une pile Daniell pour éclairer un village]
\begin{enumerate}
\item \textbf{Équation-bilan et f.é.m. standard.}
\[ \ce{Zn(s) + Cu^2+(aq) -> Zn^2+(aq) + Cu(s)} \]
\[ E^{\circ} = E^{\circ}(\ce{Cu^2+}/\ce{Cu}) - E^{\circ}(\ce{Zn^2+}/\ce{Zn}) = 0,34 - (-0,76) = \SI{+1,10}{V} \]
\item \textbf{Quantité d'électricité fournie.}
\[ \Delta t = \SI{5,0}{h} = 5,0 \times 3600 = \SI{1,8e4}{s} \]
\[ Q = I \times \Delta t = 0,10 \times \num{1,8e4} = \SI{1,8e3}{C} \]
\item \textbf{Quantité de matière d'électrons.} Une mole d'électrons transporte une charge égale au Faraday $F$ ; la charge totale étant $Q = n(e^-) \times F$, on en déduit :
\[ n(e^-) = \dfrac{Q}{F} = \dfrac{\num{1,8e3}}{\num{9,65e4}} = \SI{1,87e-2}{mol} \]
\item \textbf{Zinc consommé.} D'après la demi-équation \ce{Zn -> Zn^2+ + 2e-}, $n(\ce{Zn}) = \dfrac{n(e^-)}{2}$ :
\[ n(\ce{Zn}) = \dfrac{\num{1,87e-2}}{2} = \SI{9,3e-3}{mol} \]
\[ m(\ce{Zn}) = n(\ce{Zn}) \times M(\ce{Zn}) = \num{9,3e-3} \times 65,4 = \SI{0,61}{g} \]
\item \textbf{Cuivre déposé.} D'après \ce{Cu^2+ + 2e- -> Cu}, $n(\ce{Cu}) = \dfrac{n(e^-)}{2} = \SI{9,3e-3}{mol}$ :
\[ m(\ce{Cu}) = n(\ce{Cu}) \times M(\ce{Cu}) = \num{9,3e-3} \times 63,5 = \SI{0,59}{g} \]
\item \textbf{Durée maximale théorique.} Tableau d'avancement pour l'épuisement du zinc :

\begin{center}
\begin{tabularx}{\linewidth}{|l|X|c|c|}
\hline
\rowcolor{popblueL} \textbf{État} & \textbf{Avancement $x$ (mol)} & $n(\ce{Zn})$ (mol) & $n(e^-)$ fournis (mol) \\
\hline
Initial & $0$ & $\dfrac{10}{65,4} = 0,153$ & $0$ \\
\hline
Final (Zn épuisé) & $x_{\max} = 0,153$ & $0$ & $2x_{\max} = 0,306$ \\
\hline
\end{tabularx}
\end{center}

\[ Q_{\max} = n_{\max}(e^-) \times F = 0,306 \times \num{9,65e4} = \SI{2,95e4}{C} \]
\[ \Delta t_{\max} = \dfrac{Q_{\max}}{I} = \dfrac{\num{2,95e4}}{0,10} = \SI{2,95e5}{s} \approx \SI{82}{h} \]
Soit environ \textbf{3,4 jours} de fonctionnement théorique à \SI{0,10}{A}.
\item \textbf{Raisons possibles d'un arrêt prématuré.}
\begin{itemize}
\item La concentration des ions \ce{Cu^2+} diminue au cours du fonctionnement : la f.é.m. réelle chute en dessous de sa valeur standard (équation de Nernst) et la tension devient insuffisante pour alimenter la lampe, avant même l'épuisement du zinc.
\item Le dépôt de cuivre peut devenir poudreux et mal adhérent, ou le pont salin s'épuiser (migration cumulative des ions) : la résistance interne de la pile augmente et la tension disponible $U = E - rI$ s'effondre.
\end{itemize}
\end{enumerate}
\textbf{Barème indicatif (sur 12) :} bilan + f.é.m. 2 pts ; $Q$ 1,5 pt ; $n(e^-)$ 1,5 pt ; $m(\ce{Zn})$ 2 pts ; $m(\ce{Cu})$ 1,5 pt ; durée max 2 pts ; discussion 1,5 pt.
\begin{attention}[Points de vigilance]
Convertis toujours la durée en \emph{secondes} avant d'appliquer $Q = I\Delta t$. N'oublie pas le facteur 2 entre $n(e^-)$ et $n(\ce{Zn})$ : la demi-équation fait intervenir deux électrons par atome de zinc. Vérifie la cohérence : $m(\ce{Cu})$ doit être légèrement inférieure à $m(\ce{Zn})$ car $M(\ce{Cu}) < M(\ce{Zn})$ pour une même quantité de matière.
\end{attention}
\end{corrige}

\newpage

\coursec{Fiche bilan}

\begin{popfigure}
\begin{tikzpicture}[
    node style/.style={draw, rounded corners, thick, text centered, text width=2.4cm, minimum height=1.3cm, minimum width=2.8cm, font=\small, inner sep=3pt},
    title style/.style={node style, fill=popdark, text=white, font=\small\bfseries, minimum width=10cm, text width=9cm},
    blue box/.style={node style, fill=popblueL, draw=popblue},
    orange box/.style={node style, fill=poporangeL, draw=poporange},
    green box/.style={node style, fill=popgreenL, draw=popgreen},
    purple box/.style={node style, fill=poppurpleL, draw=poppurple},
    pink box/.style={node style, fill=poppinkL, draw=poppink},
    gold box/.style={node style, fill=popgoldL, draw=popgold},
    arrow/.style={-Stealth, thick, draw=popdark, shorten >=2pt, shorten <=2pt}
]
% Titre central
\node[title style] (title) {Synthèse visuelle : Les piles électrochimiques};

% Ligne 1 : Pile Daniell & Fonctionnement
\node[blue box, below=1.2cm of title, xshift=-4.5cm] (daniell) {\textbf{Pile Daniell}\\\ce{Zn}|\ce{Zn^2+} (aq) \quad \ce{Cu^2+} (aq)|\ce{Cu}\\Pont salin : \ce{KNO3} (sat.)};
\node[orange box, right=1.2cm of daniell] (fonct) {\textbf{Fonctionnement}\\\cle{Anode} ($\ominus$) : Oxydation\\$\ce{Zn -> Zn^2+ + 2e-}$\\\cle{Cathode} ($\oplus$) : Réduction\\$\ce{Cu^2+ + 2e- -> Cu}$};

% Ligne 2 : Notation & Référence
\node[green box, below=1.5cm of daniell] (notat) {\textbf{Notation conventionnelle}\\$\ominus\ \ce{Zn}|\ce{Zn^2+}\ (aq)\ ||\ \ce{Cu^2+}\ (aq)|\ce{Cu}\ \oplus$\\Sens $\ominus \to \oplus$ dans le circuit extérieur};
\node[purple box, right=1.2cm of notat] (ref) {\textbf{Électrode de référence}\\\cle{Électrode normale à hydrogène (ENH)}\\\ce{Pt}|\ce{H2}(g, \SI{1}{bar})|\ce{H+}(aq, \SI{1}{mol.L^{-1}})\\\cle{$E^\circ(\ce{H+}/\ce{H2}) = \SI{0,00}{V}$ à \SI{25}{\celsius}}};

% Ligne 3 : Potentiels & f.é.m.
\node[pink box, below=1.5cm of notat] (pot) {\textbf{Potentiels standards $E^\circ$}\\\cle{Classification quantitative} des couples\\\textit{Plus $E^\circ$ grand $\Rightarrow$ Oxydant fort}\\\textit{Plus $E^\circ$ petit $\Rightarrow$ Réducteur fort}};
\node[gold box, right=1.2cm of pot] (fem) {\textbf{Force électromotrice (f.é.m.)}\\$e = E_{+} - E_{-} = E^\circ_{+} - E^\circ_{-}$\\Pile Daniell : $e = \SI{0,34}{V} - (-\SI{0,76}{V}) = \SI{1,10}{V}$\\$e > 0$ : réaction spontanée};

% Flèches de liaison
\draw[arrow] (daniell) -- (fonct);
\draw[arrow] (daniell) -- (notat);
\draw[arrow] (fonct) -- (ref);
\draw[arrow] (notat) -- (pot);
\draw[arrow] (ref) -- (fem);
\draw[arrow] (pot) -- (fem);
\draw[arrow, dashed] (daniell) -- (pot); % lien transversal
\end{tikzpicture}
\legende{Organigramme récapitulatif : de la pile Daniell à la prévision de la f.é.m. via l'échelle des potentiels standards.}
\end{popfigure}

\begin{definition}[Définitions fondamentales]
\begin{itemize}
  \item \cle{Pile électrochimique} : Dispositif qui transforme l'énergie chimique d'une réaction d'oxydoréduction \emph{spontanée} en énergie électrique utilisable.
  \item \cle{Demi-pile} : Association d'un conducteur (électrode, métallique ou inerte comme le Pt) et de la solution contenant les espèces du couple redox (ex. \ce{Zn^2+}/\ce{Zn}, \ce{Fe^3+}/\ce{Fe^2+}/\ce{Pt}).
  \item \cle{Anode} (pôle $\ominus$ dans une pile) : Siège de l'\textbf{oxydation} (perte d'électrons par le réducteur). \cle{Cathode} (pôle $\oplus$) : Siège de la \textbf{réduction} (gain d'électrons par l'oxydant).
  \item \cle{Circuit extérieur} : Les électrons circulent de l'anode ($\ominus$) vers la cathode ($\oplus$).
  \item \cle{Pont salin} : Rempli d'un électrolyte inerte (souvent \ce{KNO3} saturé), il assure la neutralité électrique des deux demi-piles par migration des anions vers l'anode et des cations vers la cathode, et ferme le circuit interne.
  \item \cle{Potentiel standard d'oxydoréduction $E^\circ$} : Potentiel d'un couple \ce{Ox}/\ce{Red} mesuré dans les conditions standard (\SI{25}{\celsius}, $c = \SI{1}{mol.L^{-1}}$ pour les espèces dissoutes, $p = \SI{1}{bar}$ pour les gaz) \emph{par rapport à l'électrode normale à hydrogène (ENH)} prise pour référence à \SI{0,00}{V}.
\end{itemize}
\end{definition}

\begin{propriete}[Équations de la pile Daniell — À connaître par cœur]
\begin{itemize}
  \item \textbf{Équation à l'anode (oxydation) :} \[ \ce{Zn(s) -> Zn^2+(aq) + 2e-} \]
  \item \textbf{Équation à la cathode (réduction) :} \[ \ce{Cu^2+(aq) + 2e- -> Cu(s)} \]
  \item \textbf{Équation-bilan de fonctionnement (réaction spontanée) :} \[ \ce{Zn(s) + Cu^2+(aq) -> Zn^2+(aq) + Cu(s)} \]
  \item \textbf{Notation conventionnelle (norme IUPAC) :} L'anode ($\ominus$) s'écrit à \textbf{gauche}, la cathode ($\oplus$) à \textbf{droite}. Les changements de phase sont séparés par une barre verticale $|$, le pont salin par deux barres $||$.
  \[ \ominus\ \ce{Zn(s)}|\ce{Zn^2+(aq)}||\ce{Cu^2+(aq)}|\ce{Cu(s)}\ \oplus \]
\end{itemize}
\end{propriete}

\begin{propriete}[Formules clés et valeurs de référence]
\begin{center}
\begin{tabularx}{\linewidth}{|l|X|l|}
\hline
\rowcolor{popblueL} \textbf{Grandeur / Notion} & \textbf{Expression / Valeur} & \textbf{Unité} \\
\hline
f.é.m. d'une pile (conditions standard) & $e = E_{+} - E_{-} = E^\circ_{+} - E^\circ_{-}$ & V \\
\hline
f.é.m. de la pile Daniell & $e = \SI{0,34}{V} - (-\SI{0,76}{V}) = \SI{1,10}{V}$ & V \\
\hline
Référence absolue (ENH) & $E^\circ(\ce{H+}/\ce{H2}) = \SI{0,00}{V}$ & V \\
\hline
Règle du gamma (prévision) & L'oxydant le plus fort (plus grand $E^\circ$) réagit avec le réducteur le plus fort (plus petit $E^\circ$) & --- \\
\hline
Polarité & Pôle $\ominus$ = Couple de \textbf{plus petit} $E^\circ$ (réducteur le plus fort) & --- \\
\hline
\end{tabularx}
\end{center}
\end{propriete}

\begin{methode}[Démarche type : Prévoir la polarité et calculer la f.é.m. d'une pile]
\begin{enumerate}
  \item \textbf{Identifier les deux couples} mis en jeu (ex. \ce{Zn^2+}/\ce{Zn} et \ce{Cu^2+}/\ce{Cu}).
  \item \textbf{Reporter leurs potentiels standards $E^\circ$} sur l'échelle (donnée au Probatoire ou apprise).
  \item \textbf{Appliquer la règle du gamma} : L'oxydant du couple au $E^\circ$ le plus grand oxyde le réducteur du couple au $E^\circ$ le plus petit.
  \item \textbf{Déduire la polarité} :
        \begin{itemize}
          \item Pôle $\ominus$ (Anode) = Métal/Électrode du couple de \textbf{plus petit $E^\circ$} (siège de l'oxydation).
          \item Pôle $\oplus$ (Cathode) = Métal/Électrode du couple de \textbf{plus grand $E^\circ$} (siège de la réduction).
        \end{itemize}
  \item \textbf{Écrire les demi-équations} équilibrées (atomes + charges) :
        \begin{itemize}
          \item Anode : $\text{Réducteur} \to \text{Oxydant} + n\ce{e-}$
          \item Cathode : $\text{Oxydant} + n\ce{e-} \to \text{Réducteur}$
        \end{itemize}
  \item \textbf{Calculer la f.é.m.} : $e = E^\circ_{\text{cathode}} - E^\circ_{\text{anode}}$ (toujours positif pour une pile).
  \item \textbf{Écrire l'équation-bilan} en additionnant les demi-équations (annuler les électrons).
  \item \textbf{Écrire la notation conventionnelle} ($\ominus$ à gauche, $\oplus$ à droite).
\end{enumerate}
\end{methode}

\begin{methode}[Détermination expérimentale des pôles et de la f.é.m.]
\begin{enumerate}
  \item Monter la pile (deux béchers, électrodes, pont salin en U rempli de \ce{KNO3} saturé).
  \item Brancher un \textbf{voltmètre} (mode DC, calibre adapté, ex. \SI{2}{V} ou \SI{20}{V}) : borne \textbf{COM (noire)} sur une électrode, borne \textbf{V/$\Omega$ (rouge)} sur l'autre.
  \item Lire la tension $U$.
  \item \textbf{Interprétation} :
        \begin{itemize}
          \item Si $U > 0$ (aiguille dévie dans le bon sens / affichage positif) : La borne \textbf{COM} est reliée au pôle $\ominus$ (anode), la borne \textbf{V/$\Omega$} au pôle $\oplus$ (cathode). La f.é.m. $e = U$.
          \item Si $U < 0$ : Inverser les bornes pour identifier les pôles.
        \end{itemize}
\end{enumerate}
\end{methode}

\begin{savaistu}[Piles, corrosion et contexte camerounais]
\begin{itemize}
  \item \textbf{Corrosion des toitures en tôle (zinc) :} C'est une pile de corrosion naturelle ! L'eau de pluie (électrolyte) met en contact le zinc (anode, $E^\circ = -\SI{0,76}{V}$) et des impuretés ferreuses/cuivre (cathode, $E^\circ > -\SI{0,76}{V}$). Le zinc s'oxyde (\ce{Zn -> Zn^2+ + 2e-}), la tôle se perce. La galvanisation (zinc sur fer) protège le fer \emph{par sacrifice} : le zinc s'oxyde à la place du fer ($E^\circ_{\ce{Zn}} < E^\circ_{\ce{Fe}}$).
  \item \textbf{Électrification rurale :} Les kits solaires individuels (très répandus au Nord et à l'Est) stockent l'énergie dans des \textbf{accumulateurs plomb-acide} (piles secondaires rechargeables) ou des batteries \textbf{Lithium-ion}. Comprendre la f.é.m. et la capacité (\SI{}{Ah}) est crucial pour dimensionner l'installation.
  \item \textbf{Industrie (SONARA, CIMENCAM) :} La protection cathodique des pipelines (SONARA) et des cuves utilise des anodes sacrificielles (Mg, Al, Zn) ou du courant imposé pour empêcher la corrosion de l'acier.
  \item \textbf{Pile alcaline (lampe de poche) :} Couple \ce{Zn}/\ce{ZnO} (anode) et \ce{MnO2}/\ce{Mn2O3} (cathode) dans un électrolyte \ce{KOH}. f.é.m. $\approx \SI{1,5}{V}$. Non rechargeable (réaction peu réversible).
\end{itemize}
\end{savaistu}

\begin{aretenir}[Checklist finale avant le Probatoire]
\begin{itemize}
  \item[$\square$] Je sais définir une pile, une demi-pile, l'anode, la cathode, le pont salin.
  \item[$\square$] Je sais schématiser une pile (Daniell ou autre) : électrodes, solutions, pont salin, voltmètre, sens des électrons ($\ominus \to \oplus$), migration des ions dans le pont.
  \item[$\square$] Je sais écrire les \textbf{demi-équations aux électrodes} (équilibrées atomes + charges) et l'\textbf{équation-bilan}.
  \item[$\square$] Je maîtrise la \textbf{notation conventionnelle} : $\ominus\ \text{Anode}|\text{Electrolyte}||\text{Electrolyte}|\text{Cathode}\ \oplus$.
  \item[$\square$] Je connais le rôle exact du pont salin (neutralité + fermeture circuit) et pourquoi on choisit \ce{KNO3} (ions de mobilités voisines, pas de précipité).
  \item[$\square$] Je sais que l'\textbf{ENH} (\ce{Pt}|\ce{H2}(1 bar)|\ce{H+}(1 M)) est la référence universelle ($E^\circ = \SI{0,00}{V}$).
  \item[$\square$] Je sais \textbf{placer les couples sur l'échelle des $E^\circ$} et appliquer la \textbf{règle du gamma} pour prévoir le sens spontané.
  \item[$\square$] Je détermine la \textbf{polarité} sans hésiter : $\ominus$ = plus petit $E^\circ$ ; $\oplus$ = plus grand $E^\circ$.
  \item[$\square$] Je calcule la \textbf{f.é.m.} : $e = E^\circ_{\oplus} - E^\circ_{\ominus}$ (résultat \textbf{toujours positif} en V).
  \item[$\square$] Je sais \textbf{manipuler un voltmètre} pour identifier les pôles : affichage positif $\Rightarrow$ COM sur $\ominus$.
  \item[$\square$] \textbf{Vigilance :} Je vérifie systématiquement l'équilibrage des charges dans mes demi-équations (nombre d'électrons) avant d'additionner pour le bilan.
\end{itemize}
\end{aretenir}

\findecours

\end{document}
